\documentclass[11pt,oneside]{book}
\usepackage[T1]{fontenc}
\usepackage[utf8]{inputenc}
\usepackage{lmodern}
\usepackage[a4paper,margin=1in]{geometry}
\usepackage{url}
\usepackage[hidelinks]{hyperref}
\usepackage{enumitem}
\emergencystretch=3em

\title{Medical Myths\\[0.4em]\large Evidence-Based Questions and Answers}
\author{Chaman Singh Verma}
\date{August 2026}

\begin{document}
\frontmatter
\hypersetup{pageanchor=false}
\maketitle
\hypersetup{pageanchor=true}
\chapter*{Medical disclaimer}
\phantomsection
\addcontentsline{toc}{chapter}{Medical disclaimer}
This book is an educational reference, not a substitute for a clinical examination, diagnosis, or individualized medical advice. It does not tell a reader whether to start, stop, or change a medicine, vaccine, supplement, diet, or treatment. Urgent symptoms require prompt local medical care.
\chapter*{Preface}
\phantomsection
\addcontentsline{toc}{chapter}{Preface}
Medical myths are more than harmless pieces of folklore. Some are repeated so often that they shape decisions about medicines, vaccines, pregnancy, food, emergencies, and when to seek care. Others begin with a genuine observation and then attach an explanation that the evidence does not support. This book examines those claims carefully and separates what is established from what is uncertain.

The book is organized into two main collections. \emph{Common Medical Myths} covers widely circulated health claims found across many countries and communities. \emph{Indian Myths} focuses on beliefs associated with particular regions, families, and social or traditional practices in India. The second collection is not a description of all Indians or Hindus. Beliefs vary widely, and cultural meaning should be discussed with respect even when a medical claim is unsupported.

Each entry follows the same simple pattern: Myth, Status, and Fact. ``False'' means that the claim is not supported as stated. ``True'' means that the claim is supported within the circumstances described. ``Uncertain'' means that evidence is limited, mixed, or dependent on context. These labels describe the claim, not the worth of the person who believes it.

The explanations aim to be useful rather than merely dismissive. They identify the relevant biology, clarify important exceptions, point out possible harms from delay or inappropriate treatment, and direct readers toward appropriate professional care when necessary. The bibliography provides the evidence base used for the collection, but medical knowledge changes; readers should confirm current recommendations with qualified clinicians and relevant public-health or professional organizations.

This is an educational reference, not a substitute for examination, diagnosis, or individualized medical advice. No entry should be used by itself to start, stop, or change a medicine, vaccine, supplement, diet, or treatment. In an emergency, seek local emergency care first.
\chapter*{Reader guide and registry map}
\addcontentsline{toc}{chapter}{Reader guide and registry map}
Search stable IDs for quick lookup. False rejects a claim as stated; Uncertain means evidence is mixed or context-dependent. Urgent symptoms require local emergency care.
\tableofcontents
\mainmatter
\chapter{Common Medical Myths}
\begin{enumerate}[leftmargin=0pt,labelwidth=0pt,labelsep=0pt,itemindent=0pt]
\item[] \textbf{GEN-001.} \textbf{Myth:} You need exactly 8 glasses of water every day.\newline
\textbf{Status:} False.\newline
\textbf{Fact:} Fluid requirements vary with diet, activity, climate, body size, and health. There is no single daily volume that fits every healthy person. Fluids also come from food, beverages, and normal metabolic processes. Thirst, urine output, activity, heat exposure, and illness all affect needs. Excessive intake can be harmful, particularly when it causes low blood sodium. People with heart, kidney, or liver disease may need individualized fluid advice.
\item[] \textbf{GEN-002.} \textbf{Myth:} Going outside with wet hair causes a cold.\newline
\textbf{Status:} False.\newline
\textbf{Fact:} Colds are caused by viruses \cite{ref4}. Wet hair can make a person feel cold, but it does not create a viral infection. Infection requires exposure to a susceptible pathogen and successful transmission. Cold symptoms can follow exposure to many respiratory viruses, not merely chilling. Seasonal behavior and indoor crowding may change transmission patterns. Persistent or severe symptoms should not be attributed to wet hair alone.
\item[] \textbf{GEN-003.} \textbf{Myth:} Cold weather itself causes the common cold.\newline
\textbf{Status:} False.\newline
\textbf{Fact:} Viral infection causes it, although seasonal conditions can affect transmission \cite{ref4}. Temperature alone is not the cause of a cold. Rhinoviruses and other viruses infect the respiratory tract and produce the illness. Cooler seasons may increase spread through indoor contact and may influence viral survival. Cold air can also irritate the nose and make symptoms feel worse. Prevention therefore focuses on hygiene, ventilation, vaccination where relevant, and avoiding infectious exposure.
\item[] \textbf{GEN-004.} \textbf{Myth:} Green mucus proves you have a bacterial infection.\newline
\textbf{Status:} False.\newline
\textbf{Fact:} Mucus color alone cannot reliably distinguish bacterial from viral infection. White blood cells and inflammatory proteins can change mucus color during ordinary viral illness. Color does not identify the pathogen or prove that antibiotics will help. Clinicians consider duration, severity, examination findings, and complications instead. Unnecessary antibiotics expose people to adverse effects and promote resistance. Breathing difficulty, dehydration, severe pain, or prolonged symptoms warrant medical assessment.
\item[] \textbf{GEN-005.} \textbf{Myth:} You should ``sweat out'' a fever.\newline
\textbf{Status:} False.\newline
\textbf{Fact:} Excessive bundling can worsen discomfort and dehydration. Sweating is one way the body loses heat, but forced sweating does not remove the infection causing a fever. Heavy blankets can increase discomfort and fluid loss. Comfortable clothing, appropriate fluids, and rest are usually more sensible. Fever-reducing medicine may be used for discomfort when appropriate, following age- and condition-specific guidance. Very young infants and people with severe symptoms need prompt clinical advice.
\item[] \textbf{GEN-006.} \textbf{Myth:} Fever is always dangerous.\newline
\textbf{Status:} False.\newline
\textbf{Fact:} Fever is usually part of the body's response to infection; severity depends on temperature, age, symptoms, and cause. Fever is a sign, not a diagnosis. Many uncomplicated infections cause short-lived fever without permanent harm. Risk depends on the person's age, appearance, hydration, immune status, temperature, and associated symptoms. The number alone cannot determine whether a serious illness is present. Confusion, breathing difficulty, stiff neck, dehydration, seizure, or persistent worsening require urgent assessment.
\item[] \textbf{GEN-007.} \textbf{Myth:} Antibiotics cure colds and flu.\newline
\textbf{Status:} False.\newline
\textbf{Fact:} Antibiotics treat susceptible bacteria, not viruses \cite{ref4,ref3}. Colds and influenza are viral illnesses, so an antibacterial drug does not eliminate their cause. Antibiotics may be appropriate when a separate bacterial infection is diagnosed. Taking them unnecessarily can cause diarrhea, allergic reactions, and other harms. It also selects for bacteria that resist treatment. The correct decision depends on clinical evaluation rather than the desire for a ``strong'' medicine.
\item[] \textbf{GEN-008.} \textbf{Myth:} Stronger antibiotics are always better.\newline
\textbf{Status:} False.\newline
\textbf{Fact:} The appropriate antibiotic depends on the organism, infection site, resistance patterns, and patient factors. Antibiotic choice is a matching problem, not a strength contest. A broad-spectrum drug may be less appropriate than a narrower agent when the organism is known. Dose and duration must also account for kidney function, allergies, interactions, and tissue penetration. More drug activity can mean more adverse effects without more benefit. Culture results and local guidelines often help clinicians choose safely.
\item[] \textbf{GEN-009.} \textbf{Myth:} If symptoms improve, an infection must be completely gone.\newline
\textbf{Status:} False.\newline
\textbf{Fact:} Clinical improvement does not necessarily establish eradication. Symptoms reflect inflammation and body function, not a direct measurement of pathogen presence. Some infections improve while organisms remain or complications develop. Conversely, symptoms can persist after the infection has cleared because tissues need time to recover. Follow-up depends on the infection, treatment, risk factors, and clinical response. People should follow the prescribed plan rather than stopping or extending treatment on symptoms alone.
\item[] \textbf{GEN-010.} \textbf{Myth:} Natural remedies are inherently safe.\newline
\textbf{Status:} False.\newline
\textbf{Fact:} Natural substances can cause toxicity, interactions, and adverse effects. ``Natural'' describes origin, not dose, purity, mechanism, or safety. Plants and supplements can affect the liver, kidneys, heart rhythm, blood clotting, pregnancy, and prescription-drug levels. Products may also be contaminated, mislabeled, or variable in strength. Evidence of traditional use does not establish effectiveness for a modern indication. Discussing supplements with a clinician or pharmacist is especially important before surgery or during pregnancy.
\item[] \textbf{GEN-011.} \textbf{Myth:} Cancer is always a death sentence.\newline
\textbf{Status:} False.\newline
\textbf{Fact:} Many cancers are curable, and many others can be controlled for years. Cancer is a group of diseases with widely different biology, stage, and treatment response. Some cancers are commonly cured when found and treated early. Others are managed as chronic illness for long periods, while some remain difficult to control. Prognosis depends on pathology, spread, molecular features, general health, and available treatment. A diagnosis should prompt individualized discussion rather than an assumption of inevitable death \cite{ref16}.
\item[] \textbf{GEN-012.} \textbf{Myth:} Cancer is one disease.\newline
\textbf{Status:} False.\newline
\textbf{Fact:} Cancer comprises hundreds of biologically different diseases. Cancer begins when cells acquire abnormal growth and survival behavior, but the pathways differ among organs and tumor types. Even two tumors in the same organ can have different mutations and treatment responses. Classification by tissue, grade, stage, and molecular features guides management. A treatment that works for one cancer may be ineffective or harmful for another. General cancer claims therefore need a specific cancer context.
\item[] \textbf{GEN-013.} \textbf{Myth:} Cancer is contagious.\newline
\textbf{Status:} False.\newline
\textbf{Fact:} Cancer itself generally is not transmissible between people, although some cancerassociated infections are. Malignant cells do not ordinarily spread from one person to another through ordinary contact, sex, food, or shared air. Some infections, including certain types of HPV, hepatitis viruses, and Helicobacter pylori, can increase cancer risk. Preventing or treating those infections can reduce risk without implying that the cancer itself is contagious. Rare transplant-related transmission is an exceptional medical circumstance. Cancer patients should not be isolated because of fear of ordinary contagion.
\item[] \textbf{GEN-014.} \textbf{Myth:} A biopsy causes cancer to spread.\newline
\textbf{Status:} False.\newline
\textbf{Fact:} Properly performed biopsies do not ordinarily cause clinically meaningful cancer spread. A biopsy removes or samples tissue so that pathology can identify what a mass is. Needle and surgical biopsy techniques are selected according to the organ, suspected disease, and diagnostic objective. Tumor-cell displacement is a theoretical concern in some procedures, but clinically important spread from standard biopsies is rare. Avoiding a needed biopsy can delay diagnosis and treatment. The procedure's risks should be discussed with the treating team rather than replaced by a myth.
\item[] \textbf{GEN-015.} \textbf{Myth:} Surgery exposes cancer to air and makes it spread.\newline
\textbf{Status:} False.\newline
\textbf{Fact:} Exposure to air does not cause cancer metastasis. Metastasis requires tumor cells to survive, enter appropriate pathways, and establish growth at another site. Opening tissue during a medically indicated operation does not create this process by itself. Surgeons plan incisions, handling, margins, and containment to reduce disease spread and obtain control. Surgery may be recommended precisely because it can remove localized cancer. Delaying indicated surgery because of this belief can reduce the chance of cure.
\item[] \textbf{GEN-016.} \textbf{Myth:} Sugar directly ``feeds'' cancer, so eliminating sugar cures it.\newline
\textbf{Status:} False.\newline
\textbf{Fact:} Cancer metabolism is complex; eliminating dietary sugar does not cure cancer. All cells use glucose and the body can make glucose even when dietary carbohydrate is restricted. Tumors differ in how they obtain and use energy, so a simple food prohibition cannot selectively starve cancer cells. A highly restrictive diet can cause weight loss, malnutrition, or treatment intolerance. Limiting excessive added sugar may support general health, but it is not a substitute for cancer therapy. Nutrition plans should account for diagnosis, treatment, weight, and symptoms.
\item[] \textbf{GEN-017.} \textbf{Myth:} Only smokers get lung cancer.\newline
\textbf{Status:} False.\newline
\textbf{Fact:} Nonsmokers can develop lung cancer too. Tobacco is a major lung-cancer risk factor, but it is not the only one. Radon, secondhand smoke, occupational exposures, air pollution, prior radiation, genetics, and chance can contribute. Some tumors in people who never smoked have specific molecular changes that guide targeted treatment. Smoking status affects risk but does not determine whether symptoms deserve evaluation. Persistent cough, blood in sputum, chest pain, or unexplained weight loss should not be dismissed in a nonsmoker.
\item[] \textbf{GEN-018.} \textbf{Myth:} Cancer always produces symptoms early.\newline
\textbf{Status:} False.\newline
\textbf{Fact:} Some cancers remain asymptomatic until relatively advanced. Early tumors may be too small or located where they do not disturb normal function. Some cancers are found through screening before symptoms appear, while others have no effective routine screening test. New, persistent, or changing symptoms still deserve evaluation, but absence of symptoms is not proof of absence. Screening recommendations depend on age, risk, and evidence of benefit. Overtesting without an indication can also produce false positives and unnecessary procedures.
\item[] \textbf{GEN-019.} \textbf{Myth:} Normal blood tests rule out cancer.\newline
\textbf{Status:} False.\newline
\textbf{Fact:} Many cancers can exist despite routine blood tests being normal. Routine blood tests measure selected body functions, not every tumor directly. Some cancers do not alter blood counts, liver tests, or chemistry values until late disease or treatment. Diagnosis may require examination, imaging, endoscopy, tissue sampling, or disease-specific testing. A normal result can be reassuring in context but cannot override a concerning symptom or examination finding. Clinicians interpret tests using the whole clinical picture.
\item[] \textbf{GEN-020.} \textbf{Myth:} Alternative medicine can replace proven cancer treatment.\newline
\textbf{Status:} False.\newline
\textbf{Fact:} No alternative therapy has been shown to replace effective evidence-based cancer treatment generally. Complementary practices may help some people manage symptoms, stress, or wellbeing when used alongside appropriate care. They should not be assumed to kill cancer because a product is natural, traditional, or supported by testimonials. Delaying surgery, radiation, systemic therapy, or monitoring can allow disease to progress. Supplements can also interact with cancer drugs or affect bleeding and anesthesia. Treatment choices should be discussed with an oncology team that can integrate goals, evidence, harms, and preferences.
\item[] \textbf{GEN-021.} \textbf{Myth:} Wearing a bra causes breast cancer.\newline
\textbf{Status:} False.\newline
\textbf{Fact:} Evidence does not support this. Breast cancer develops through genetic, hormonal, environmental, and tissue-level processes, not restriction of lymph flow by ordinary clothing. Bra use has not been established as a causal risk factor. A bra can cause pressure, skin irritation, or discomfort without causing malignancy. Breast symptoms should be evaluated based on the symptom itself, not blamed on underwear. Persistent lumps, skin changes, nipple changes, or discharge deserve clinical assessment.
\item[] \textbf{GEN-022.} \textbf{Myth:} Wearing a loose bra causes breast cancer.\newline
\textbf{Status:} False.\newline
\textbf{Fact:} Bra looseness has not been established as a breast-cancer cause. There is no accepted biological mechanism by which a loose garment would initiate breast cancer. Cancer risk is not determined by whether a bra is tight, loose, wired, or wireless. Comfort and skin health may influence clothing choices, but they are separate from cancer prevention. Avoiding medical evaluation because a bra seems responsible could delay diagnosis. Breast-cancer risk assessment should focus on age, family history, reproductive factors, genetics, and other established factors.
\item[] \textbf{GEN-023.} \textbf{Myth:} Every breast lump is cancer.\newline
\textbf{Status:} False.\newline
\textbf{Fact:} Many breast lumps are benign. Cysts, fibroadenomas, hormonal changes, infection, and other benign conditions can produce lumps. A benign explanation cannot be confirmed reliably by touch alone. Age, size, mobility, timing, skin changes, and imaging findings all contribute to evaluation. A new or persistent lump should be assessed rather than ignored. The correct response is examination and appropriate imaging, not automatic panic or reassurance.
\item[] \textbf{GEN-024.} \textbf{Myth:} Breast cancer lumps are always painful.\newline
\textbf{Status:} False.\newline
\textbf{Fact:} Many breast cancers are painless. Pain is neither required for cancer nor sufficient to rule it out. Hormonal changes, cysts, inflammation, injury, and benign breast disorders can all cause pain. Some cancers present as painless masses or as skin, nipple, or architectural changes. A persistent change should be assessed even when it does not hurt. Conversely, breast pain alone is often not cancer, so evaluation should be guided by the whole clinical picture.
\item[] \textbf{GEN-025.} \textbf{Myth:} A painless breast lump is harmless.\newline
\textbf{Status:} False.\newline
\textbf{Fact:} Pain does not reliably distinguish benign from malignant masses. The absence of pain provides little diagnostic information. Some malignant and benign masses are both painless, while some benign conditions are tender. Clinical examination cannot always determine the tissue type. Diagnostic mammography, ultrasound, follow-up, or biopsy may be recommended according to age and findings. Delaying evaluation because a lump is painless is unsafe.
\item[] \textbf{GEN-026.} \textbf{Myth:} Men cannot get breast cancer.\newline
\textbf{Status:} False.\newline
\textbf{Fact:} Men can develop breast cancer, although it is much less common. Male breast tissue can undergo malignant change despite the lower incidence. Risk can be influenced by age, genetics, hormone conditions, radiation, and family history. A firm lump, nipple retraction, discharge, or skin change in a man should not be dismissed as impossible cancer. Benign gynecomastia is also common and may mimic a mass. Evaluation is based on findings, not sex alone.
\item[] \textbf{GEN-027.} \textbf{Myth:} Mammography gives enough radiation to make screening more dangerous than useful.\newline
\textbf{Status:} False.\newline
\textbf{Fact:} Mammography uses a low radiation dose, and recommended screening has a favorable benefit-risk balance for appropriate populations. Mammography does involve ionizing radiation, so screening is not risk-free. Its value comes from finding some cancers earlier, when treatment may be more effective. Benefits and harms depend on age, breast density, risk, screening interval, false positives, and overdiagnosis. Screening recommendations therefore differ among organizations and countries. A person should discuss individualized screening rather than assume that all mammography is either harmless or harmful.
\item[] \textbf{GEN-028.} \textbf{Myth:} Ultrasound can completely replace mammography.\newline
\textbf{Status:} False.\newline
\textbf{Fact:} They provide different information and are often complementary. Ultrasound can characterize some masses and is useful in selected breasts and clinical situations. Mammography can reveal calcifications and architectural patterns that ultrasound may miss. Ultrasound also produces false positives and is not a universal replacement for screening mammography. The appropriate test depends on age, symptoms, breast density, pregnancy, and local guidance. A normal result on one modality does not always settle a concerning problem.
\item[] \textbf{GEN-029.} \textbf{Myth:} BI-RADS 5 is a definitive cancer diagnosis.\newline
\textbf{Status:} False.\newline
\textbf{Fact:} It indicates a likelihood of malignancy of at least 95 percent; histopathology establishes the diagnosis. BI-RADS is a standardized imaging assessment and recommendation system. Category 5 means the imaging appearance is highly suspicious and tissue diagnosis is usually recommended. It is not the same as a pathology report. A small proportion of highly suspicious lesions may still prove benign. Treatment decisions should therefore follow diagnostic confirmation and multidisciplinary interpretation.
\item[] \textbf{GEN-030.} \textbf{Myth:} Heart disease affects mostly men.\newline
\textbf{Status:} False.\newline
\textbf{Fact:} Cardiovascular disease is a major cause of illness and death in women as well. Cardiovascular disease affects people of all sexes, although risk, presentation, and timing can differ. Women may have typical chest pressure or less typical symptoms such as breathlessness, nausea, fatigue, or back and jaw discomfort. Social and clinical under-recognition can delay treatment. Risk factors include blood pressure, diabetes, lipids, smoking, age, family history, and pregnancy-related conditions. Prevention and evaluation should not be limited by gender stereotypes.
\item[] \textbf{GEN-031.} \textbf{Myth:} Heart attacks always cause dramatic chest pain.\newline
\textbf{Status:} False.\newline
\textbf{Fact:} Symptoms can be mild, atypical, or occasionally silent. Chest pressure is common but is not present in every heart attack. Symptoms may include breathlessness, sweating, nausea, unusual fatigue, or discomfort in the arm, back, jaw, or upper abdomen. Older adults and people with diabetes may have less typical or minimal symptoms. A person should not wait for severe pain before seeking emergency help. New, concerning symptoms require urgent assessment because early treatment can limit heart-muscle damage.
\item[] \textbf{GEN-032.} \textbf{Myth:} Young people cannot have heart attacks.\newline
\textbf{Status:} False.\newline
\textbf{Fact:} They can, although risk increases substantially with age. Age changes probability but does not make an event impossible. Smoking, inherited lipid disorders, diabetes, hypertension, obesity, stimulant use, inflammation, and clotting disorders can affect younger people. Some young patients also have unusual coronary anatomy or spontaneous artery injury. Severe or unexplained chest symptoms should be assessed according to presentation, not dismissed because of age. Risk-factor prevention remains useful throughout adulthood.
\item[] \textbf{GEN-033.} \textbf{Myth:} High blood pressure always causes symptoms.\newline
\textbf{Status:} False.\newline
\textbf{Fact:} Hypertension is frequently asymptomatic. Blood pressure can remain elevated for years while a person feels well. Headache or dizziness is nonspecific and cannot diagnose hypertension. The condition is detected by properly performed measurements, often repeated or confirmed outside the clinic. Untreated high pressure can damage arteries, heart, brain, kidneys, and eyes over time. Extremely high readings with neurological, chest, or breathing symptoms require urgent care.
\item[] \textbf{GEN-034.} \textbf{Myth:} If you feel fine, your blood pressure must be normal.\newline
\textbf{Status:} False.\newline
\textbf{Fact:} Symptoms are a poor indicator of blood pressure. Feeling well does not measure arterial pressure. Some people with normal readings feel unwell, while others with high readings have no symptoms. Accurate cuff size, positioning, rest, and repeated measurements matter. Home or ambulatory monitoring may clarify whether a clinic reading reflects a persistent problem. Medication decisions should be based on validated measurements and clinical context, not sensations alone.
\item[] \textbf{GEN-035.} \textbf{Myth:} Once blood pressure becomes normal on medication, hypertension is cured.\newline
\textbf{Status:} False.\newline
\textbf{Fact:} The medication may be what is keeping it controlled. A normal treated reading usually means control, not proof that the underlying tendency has disappeared. Stopping medication abruptly can allow pressure to rise and may create additional risk. Weight change, diet, exercise, sleep, kidney disease, and other factors can alter treatment needs. Some people can reduce therapy under clinician supervision, while others need long-term treatment. Any change should follow a monitoring plan.
\item[] \textbf{GEN-036.} \textbf{Myth:} All dietary cholesterol directly becomes blood cholesterol.\newline
\textbf{Status:} False.\newline
\textbf{Fact:} Blood lipid levels reflect multiple dietary, metabolic, and genetic factors. Digestion and liver metabolism regulate how dietary cholesterol affects circulating lipoproteins. Saturated and trans fats, overall dietary pattern, genetics, insulin resistance, and medications can be important. The effect of one food cannot be inferred from its cholesterol content alone. People with established cardiovascular disease or inherited lipid disorders may need individualized advice. A lipid panel and overall risk assessment are more informative than a single food rule \cite{ref10}.
\item[] \textbf{GEN-037.} \textbf{Myth:} All fats are bad for the heart.\newline
\textbf{Status:} False.\newline
\textbf{Fact:} Different types of dietary fat have different cardiovascular effects. Unsaturated fats from foods such as nuts, seeds, fish, and many plant oils can fit into heart-healthy patterns. Saturated and industrial trans fats have different risk profiles and should not be treated as identical to unsaturated fats. Total diet quality and replacement foods matter when changing fat intake. A food's fat content alone does not determine its health effect. Dietary advice should account for medical conditions, cultural foods, and sustainable habits.
\item[] \textbf{GEN-038.} \textbf{Myth:} A normal ECG proves the heart is completely healthy.\newline
\textbf{Status:} False.\newline
\textbf{Fact:} Many cardiovascular conditions can exist despite a normal resting ECG. A resting ECG records electrical activity during a short interval. It can identify some rhythm problems, conduction abnormalities, prior injury, or patterns of ischemia, but it cannot assess every structure or circumstance. Coronary disease, intermittent arrhythmia, valve disease, and early heart failure may require other tests. Test interpretation depends on symptoms, examination, risk, and timing. A normal result is useful information but not a universal clearance.
\item[] \textbf{GEN-039.} \textbf{Myth:} Coughing vigorously can reliably stop your own heart attack.\newline
\textbf{Status:} False.\newline
\textbf{Fact:} ``Cough CPR'' is not an appropriate self-treatment for a suspected heart attack. Forceful coughing is not a validated way to reopen a blocked coronary artery. It can delay emergency contact and distract from evidence-based care. A suspected heart attack requires calling emergency services and following dispatcher instructions. Collapse without normal breathing requires immediate CPR and an automated external defibrillator when available. Advice for a monitored hospital rhythm problem is not the same as home treatment for chest pain.
\item[] \textbf{GEN-040.} \textbf{Myth:} Eating sugar directly causes diabetes.\newline
\textbf{Status:} False.\newline
\textbf{Fact:} Diabetes results from complex genetic, metabolic, lifestyle, and other factors. Type 1 diabetes is an autoimmune disease and is not caused by eating sweets. Type 2 diabetes involves insulin resistance, beta-cell function, inherited risk, body composition, activity, sleep, medications, and environment. High-sugar diets can contribute excess energy and weight gain, which may increase type 2 risk, but one food does not create diabetes by itself. Diagnosis requires validated laboratory criteria rather than assumptions about diet.
\item[] \textbf{GEN-041.} \textbf{Myth:} People with diabetes can never eat carbohydrates.\newline
\textbf{Status:} False.\newline
\textbf{Fact:} Carbohydrate quantity and quality can be managed rather than universally eliminated. Carbohydrates are present in fruit, grains, beans, vegetables, milk, and many mixed foods. Their effect depends on portion, fiber, processing, medication, activity, and individual glucose response. Some people need structured carbohydrate planning, but complete avoidance is usually unnecessary and difficult to sustain. Dietary patterns should provide adequate nutrition while supporting glucose targets. A diabetes educator or dietitian can tailor the plan to treatment and culture.
\item[] \textbf{GEN-042.} \textbf{Myth:} Diabetes is always caused by obesity.\newline
\textbf{Status:} False.\newline
\textbf{Fact:} Body weight is only one risk factor, and this claim is particularly wrong for type 1 diabetes. Type 1 diabetes results primarily from autoimmune destruction of insulin-producing beta cells. Type 2 diabetes reflects inherited susceptibility, insulin resistance, beta-cell function, body composition, activity, sleep, medications, and social environment. People at many body sizes can develop either type, although population risks differ. Weight stigma can delay diagnosis and reduce quality of care. Diagnosis should be based on glucose testing and clinical assessment, not appearance.
\item[] \textbf{GEN-043.} \textbf{Myth:} Thin people cannot develop type 2 diabetes.\newline
\textbf{Status:} False.\newline
\textbf{Fact:} They can. Body size does not reveal insulin sensitivity, beta-cell reserve, family risk, or fat distribution. Some people with apparently normal weight have substantial metabolic risk or sarcopenic obesity. Other causes of diabetes can also be mistaken for type 2 disease. Symptoms and laboratory testing are more reliable than visual assumptions. A person with thirst, urination, weight change, recurrent infections, or risk factors may need testing regardless of weight.
\item[] \textbf{GEN-044.} \textbf{Myth:} Diabetes always causes obvious symptoms.\newline
\textbf{Status:} False.\newline
\textbf{Fact:} Type 2 diabetes can remain unnoticed for years. Glucose can rise gradually without dramatic thirst, urination, or weight loss. Some people first learn they have diabetes through screening or a complication. Symptoms may also be attributed to aging, stress, or another illness. Risk-based screening and periodic testing can identify disease earlier. New severe thirst, vomiting, confusion, or rapid breathing requires urgent medical care.
\item[] \textbf{GEN-045.} \textbf{Myth:} One normal fasting glucose proves you do not have diabetes.\newline
\textbf{Status:} False.\newline
\textbf{Fact:} Diagnosis depends on validated criteria and sometimes repeat testing. Glucose varies with fasting duration, illness, stress, medication, and time of day. A single measurement may miss post-meal hyperglycemia or impaired glucose regulation. Clinicians may use repeat fasting glucose, A1C, an oral glucose tolerance test, or random glucose with symptoms. The correct test depends on the clinical situation. A normal result is reassuring but cannot answer every question about future or intermittent dysglycemia.
\item[] \textbf{GEN-046.} \textbf{Myth:} Honey is harmless for blood glucose because it is natural.\newline
\textbf{Status:} False.\newline
\textbf{Fact:} Honey contains sugars and can raise glucose. Honey is primarily a concentrated source of simple carbohydrates. Its natural origin does not prevent digestion into glucose and other sugars. Small portions may fit some meal plans, but honey is not a glucose-free sweetener. People using insulin or glucose-lowering medicines should consider its carbohydrate content. ``Natural'' should not be used as a substitute for reading portions and monitoring response.
\item[] \textbf{GEN-047.} \textbf{Myth:} Brown sugar is substantially healthier for diabetes than white sugar.\newline
\textbf{Status:} False.\newline
\textbf{Fact:} Their effects on blood glucose are broadly similar. Brown sugar usually contains white sugar with a small amount of molasses. Both provide readily absorbed sucrose and similar calories per comparable amount. Color and processing do not make brown sugar a free food for diabetes. Flavor preference may help reduce total use, but it does not remove the glucose effect. Overall dietary pattern, portion, fiber, and treatment remain more important.
\item[] \textbf{GEN-048.} \textbf{Myth:} Starting insulin means you have ``failed.''\newline
\textbf{Status:} False.\newline
\textbf{Fact:} Insulin is an appropriate and sometimes essential treatment. Diabetes is a biological disease, not a moral test of willpower. Type 1 diabetes requires insulin for survival, and people with type 2 diabetes may need it as beta-cell function changes or during illness and pregnancy. Insulin can reduce dangerous hyperglycemia and protect against acute metabolic emergencies. Treatment may change over time without representing personal failure. Education about dosing, hypoglycemia, storage, and monitoring helps people use it safely.
\item[] \textbf{GEN-049.} \textbf{Myth:} Diabetes medicines make the pancreas permanently dependent on them.\newline
\textbf{Status:} False.\newline
\textbf{Fact:} Diabetes progression and treatment requirements reflect the underlying disease, not drug addiction. Medicines alter glucose physiology; they do not create a psychological dependence like an addictive substance. Some treatments reduce the workload on beta cells, while others replace hormones or improve insulin sensitivity. The need for more treatment can reflect disease progression, illness, pregnancy, or changing physiology. Stopping medication without guidance can cause dangerous hyperglycemia. Treatment should be adjusted through monitoring and clinician review.
\item[] \textbf{GEN-050.} \textbf{Myth:} Humans use only 10 percent of their brains.\newline
\textbf{Status:} False.\newline
\textbf{Fact:} Brain imaging and neurological evidence show widespread brain activity and function. Different tasks recruit different networks, and not every neuron is active at the same instant. That does not mean most brain tissue is permanently unused. Injury to seemingly small regions can produce major effects because brain functions are distributed and interconnected. The myth likely persists because people like the idea of hidden potential. Learning can improve performance, but it does not require activating a dormant 90 percent.
\item[] \textbf{GEN-051.} \textbf{Myth:} Adults cannot form new brain cells.\newline
\textbf{Status:} Uncertain.\newline
\textbf{Fact:} Adult neurogenesis has evidence in some brain regions, although its extent in humans remains an active research area. New neurons are clearly generated in some adult mammalian brain regions. Evidence for meaningful adult hippocampal neurogenesis in humans is substantial but remains debated because human tissue is difficult to study and methods differ. It is therefore inaccurate to say that adults never form new brain cells. At the same time, claims that large numbers of new neurons routinely replace old ones are not established. Exercise, sleep, stress, and aging may influence neuroplasticity, but no supplement has been proven to create a major increase in human neurogenesis. The safest conclusion is that adult brains retain plasticity, while the exact extent of neuron production in humans remains under study.
\item[] \textbf{GEN-052.} \textbf{Myth:} Memory works like a perfect video recording.\newline
\textbf{Status:} False.\newline
\textbf{Fact:} Memory is reconstructive and susceptible to distortion. Remembering involves reconstructing information from stored traces, context, expectations, and later cues. Rehearsal can strengthen memory, but it does not guarantee exactness. Confidence and vividness are not reliable measures of accuracy. Stress, suggestion, delay, and repeated questioning can alter recall. People may sincerely report details that were inferred or introduced after an event. This is why clinical and legal conclusions should consider corroborating evidence rather than memory confidence alone.
\item[] \textbf{GEN-053.} \textbf{Myth:} Mental illness is simply a lack of willpower.\newline
\textbf{Status:} False.\newline
\textbf{Fact:} Mental disorders involve complex biological, psychological, and environmental factors. Mental disorders can involve genetic vulnerability, brain development, physiology, trauma, social conditions, and learned behavior. Willingness may affect engagement with treatment, but it does not determine whether a disorder exists. Symptoms can impair motivation, concentration, judgment, sleep, and energy themselves. Psychotherapy, medication, social support, and practical changes can all be useful, depending on the condition. Blame and stigma can delay assessment and worsen isolation. Compassionate, evidence-based care is more helpful than treating symptoms as a character defect.
\item[] \textbf{GEN-054.} \textbf{Myth:} Depression is just ordinary sadness.\newline
\textbf{Status:} False.\newline
\textbf{Fact:} Clinical depression is a defined disorder involving multiple cognitive, emotional, behavioral, and physical symptoms. Sadness is a normal emotion that often changes with circumstances. Major depressive episodes involve persistent depressed mood or loss of interest together with symptoms such as sleep, appetite, energy, concentration, guilt, or suicidal thinking changes. Symptoms cause clinically meaningful distress or impairment and are not defined solely by how sad someone appears. Depression can occur without an obvious external trigger. It can affect children, adolescents, adults, and older people. Evaluation is important because effective treatments and urgent safety support are available.
\item[] \textbf{GEN-055.} \textbf{Myth:} People with depression can simply ``snap out of it.''\newline
\textbf{Status:} False.\newline
\textbf{Fact:} Depression is not corrected merely by deciding to be happier. Depression can alter motivation, reward processing, sleep, concentration, and physical energy. Encouragement and routine may support recovery, but they are not substitutes for assessment when symptoms persist or impair functioning. Psychotherapies, antidepressant medicines, activity-based approaches, and treatment of contributing medical conditions may help. Treatment choice should be individualized, especially for younger people, pregnancy, bipolar symptoms, or medication interactions. Asking directly about suicide does not create suicidal thoughts and can open a path to safety. Urgent help is needed for imminent self-harm risk.
\item[] \textbf{GEN-056.} \textbf{Myth:} Dementia is an inevitable part of normal aging.\newline
\textbf{Status:} False.\newline
\textbf{Fact:} Dementia is pathological, although its prevalence rises strongly with age. Normal aging may bring slower recall or occasional forgetfulness without major loss of independence. Dementia involves acquired decline in one or more cognitive domains that interferes with daily life. Its causes include Alzheimer disease, vascular injury, Lewy body disease, frontotemporal degeneration, and other conditions. Some contributors, such as medication effects, depression, sleep disorders, thyroid disease, or vitamin deficiency, may be treatable. Risk is influenced by age, genetics, vascular health, hearing, education, and lifestyle factors. New or worsening cognitive impairment deserves medical evaluation rather than dismissal as aging.
\item[] \textbf{GEN-057.} \textbf{Myth:} Alzheimer's disease and dementia are synonymous.\newline
\textbf{Status:} False.\newline
\textbf{Fact:} Alzheimer's is one cause of dementia. Dementia is a clinical syndrome involving cognitive decline that disrupts everyday function. Alzheimer disease is a specific neurodegenerative disease and is the most common cause of dementia, but it is not the only one. Vascular cognitive impairment, Lewy body disease, frontotemporal degeneration, and mixed pathologies can produce dementia. Different causes may have different early symptoms, tests, treatments, and support needs. A person can also have Alzheimer pathology together with vascular or other pathology. Accurate diagnosis helps guide prognosis, safety planning, and management.
\item[] \textbf{GEN-058.} \textbf{Myth:} All seizures involve violent shaking.\newline
\textbf{Status:} False.\newline
\textbf{Fact:} Some seizures cause staring, altered awareness, sensory phenomena, or subtle movements. Seizures arise from abnormal electrical activity and can affect movement, awareness, sensation, behavior, or autonomic function. Generalized convulsive seizures may cause stiffening and rhythmic jerking, but focal seizures may cause staring, unusual sensations, repetitive movements, or brief impaired awareness. Some events are difficult to recognize without a witness description or video. A first suspected seizure requires medical assessment because causes range from epilepsy to acute metabolic or structural problems. During a seizure, protect the person from injury and time the event. Emergency help is appropriate for prolonged, repeated, or dangerous seizures.
\item[] \textbf{GEN-059.} \textbf{Myth:} You should put something in the mouth of someone having a seizure.\newline
\textbf{Status:} False.\newline
\textbf{Fact:} This can cause injury; do not put objects in their mouth. A person cannot swallow their tongue during a seizure. Placing a spoon, fingers, or another object in the mouth can break teeth, obstruct breathing, or injure the rescuer. Clear nearby hazards and cushion the head if possible. Do not restrain the movements or give food, drink, or medicine until full alertness returns. When convulsions stop, turn the person onto their side and monitor breathing. Call emergency services if the seizure lasts five minutes or more, repeats, causes serious injury, occurs in water or pregnancy, or is a first known seizure.
\item[] \textbf{GEN-060.} \textbf{Myth:} Everyone needs a daily detox.\newline
\textbf{Status:} False.\newline
\textbf{Fact:} The liver, kidneys, lungs, gastrointestinal tract, and other systems continuously process and eliminate metabolic products. Healthy bodies already have organs that metabolize and excrete many substances. Commercial detox diets, teas, and supplements use an undefined idea of ``toxins'' and generally lack evidence for removing harmful substances beyond normal physiology. Restrictive regimens can cause dehydration, electrolyte disturbance, nutrient deficiency, drug interactions, or delayed treatment. Some products marketed as detoxes have contained undeclared pharmacologically active ingredients. Specific poisonings require poison-control or emergency medical guidance, not a general cleanse. Adequate nutrition, sleep, activity, vaccination, and appropriate treatment support health more reliably.
\item[] \textbf{GEN-061.} \textbf{Myth:} Detox juices remove unspecified toxins from the body.\newline
\textbf{Status:} False.\newline
\textbf{Fact:} Such claims generally lack credible clinical evidence. ``Toxins'' in detox advertising are usually not identified, measured, or linked to a demonstrated disease process. Juices may provide fluids and some micronutrients, but that is different from medically detoxifying the body. Very restrictive juice regimens can provide too little protein, fiber, or energy and may destabilize glucose. Some products also interact with medicines or contain contaminants. The body uses the liver, kidneys, lungs, gut, and skin to handle normal metabolic waste. Known poisonings require a specific medical treatment, not a generalized juice cleanse. \emph{(See also GEN-060.)}
\item[] \textbf{GEN-062.} \textbf{Myth:} Gluten is unhealthy for everyone.\newline
\textbf{Status:} False.\newline
\textbf{Fact:} Gluten avoidance is medically important for conditions such as celiac disease but is not necessary for everyone. Gluten is harmful in celiac disease and can cause symptoms in some people with wheat allergy or possible non-celiac sensitivity. For most people, however, wheat and other gluten-containing foods can be part of a nutritious diet. Removing them without a reason may reduce fiber and lead to unnecessary reliance on highly processed substitutes. People with suspected celiac disease should be evaluated before starting a gluten-free diet because testing can become less informative after gluten withdrawal. A gluten-free label does not automatically mean a food is lower in sugar, calories, or sodium. Treatment should match the diagnosed condition.
\item[] \textbf{GEN-063.} \textbf{Myth:} Carbohydrates are inherently unhealthy.\newline
\textbf{Status:} False.\newline
\textbf{Fact:} Health effects depend strongly on carbohydrate type, quantity, overall diet, and individual metabolism. Carbohydrates include fruits, vegetables, legumes, whole grains, milk sugars, and refined sweets, which do not have identical nutritional profiles. Fiber-rich carbohydrate foods can support bowel function, satiety, and cardiometabolic health. Refined products and sugar-sweetened drinks are easier to overconsume and often provide fewer protective nutrients. People with diabetes may need individualized carbohydrate planning, not universal avoidance. Total diet quality, portions, activity, and medical conditions all matter. Replacing carbohydrate indiscriminately with saturated fat is not automatically healthier.
\item[] \textbf{GEN-064.} \textbf{Myth:} Eating fat automatically makes you fat.\newline
\textbf{Status:} False.\newline
\textbf{Fact:} Body weight depends primarily on long-term energy balance and numerous biological and behavioral factors. Dietary fat is energy-dense, but it is not converted into body fat in a simple one-food, one-outcome process. Weight change reflects sustained energy intake and expenditure together with appetite, sleep, medications, genetics, and environment. Unsaturated fats from foods such as nuts, seeds, fish, and olive oil can fit within heart-healthy dietary patterns. Replacing fat with refined starch or added sugar may not improve health. Portion size and overall food pattern are more informative than labeling all fat as harmful. Individual medical advice may differ for specific conditions.
\item[] \textbf{GEN-065.} \textbf{Myth:} Skipping breakfast damages everyone's metabolism.\newline
\textbf{Status:} False.\newline
\textbf{Fact:} Meal timing effects vary; breakfast is not metabolically mandatory for every healthy adult. Studies of breakfast and weight do not show a universal rule that every person must eat early to maintain metabolism. Some people feel better and eat more regularly with breakfast, while others can meet nutritional needs with a later first meal. Breakfast skipping may be associated with poorer diet quality in some populations, but association does not prove that skipping caused the outcome. People taking glucose-lowering medicines, pregnant people, children, and those with eating-disorder risk may need individualized guidance. Food quality and total intake remain central. A consistent pattern that supports health and adherence is usually more useful than a rigid clock rule.
\item[] \textbf{GEN-066.} \textbf{Myth:} Eating after a particular hour automatically causes weight gain.\newline
\textbf{Status:} False.\newline
\textbf{Fact:} Total energy intake, food quality, circadian factors, and activity matter more than a universal cutoff time. Calories do not become uniquely fattening at midnight. Late eating can correlate with irregular schedules, larger portions, sleep disruption, or highly palatable snack foods, which may affect weight indirectly. Circadian biology may influence glucose and appetite responses, but there is no single cutoff that applies to everyone. Shift workers and people with medical conditions require practical, individualized schedules. Avoiding large meals immediately before sleep may help reflux or sleep for some people. Overall dietary pattern and sustained intake are more important than the clock alone.
\item[] \textbf{GEN-067.} \textbf{Myth:} Sea salt is dramatically healthier than table salt.\newline
\textbf{Status:} False.\newline
\textbf{Fact:} Both are primarily sodium chloride; sodium intake remains relevant. Sea salt and table salt contain similar amounts of sodium by weight. Sea salt may retain trace minerals, but these are usually too small to create a major health advantage. Table salt is often iodized, which helps prevent iodine deficiency. Excess sodium can raise blood pressure in susceptible people regardless of its source. Flavor and crystal size can change how much is used per spoon, but not the basic sodium physiology. Reading labels and limiting total sodium is more useful than choosing salt by marketing language.
\item[] \textbf{GEN-068.} \textbf{Myth:} Supplements can replace a balanced diet.\newline
\textbf{Status:} False.\newline
\textbf{Fact:} Supplements cannot generally reproduce all benefits of nutritious foods. Whole foods provide interacting nutrients, fiber, protein, water, and bioactive compounds that a single supplement does not replicate. Supplements are valuable when they treat a documented deficiency, meet a life-stage need, or are prescribed for a specific condition. They can also cause toxicity, interactions, or false reassurance when taken unnecessarily. Some products have variable quality and may not contain what their labels imply. A supplement cannot reliably compensate for a diet dominated by ultra-processed foods. Clinicians should review supplements when a person uses medicines or has kidney, liver, or pregnancy-related concerns.
\item[] \textbf{GEN-069.} \textbf{Myth:} More vitamins are always better.\newline
\textbf{Status:} False.\newline
\textbf{Fact:} Excessive doses of some vitamins can be harmful. Vitamins are essential in appropriate amounts, but benefit does not increase indefinitely with dose. Fat-soluble vitamins can accumulate, and high doses of several water-soluble vitamins also cause adverse effects. Excess vitamin A can harm the liver and fetus, while excess vitamin D can cause hypercalcemia and kidney injury. High-dose supplements can interfere with laboratory tests or medicines. Deficiency should be confirmed or reasonably suspected before replacement when possible. Food-first intake and evidence-based dosing are safer than megadose products.
\item[] \textbf{GEN-070.} \textbf{Myth:} Antibiotics prevent every infection.\newline
\textbf{Status:} False.\newline
\textbf{Fact:} They work only against susceptible bacteria and have no direct effect on most viruses. Antibiotics treat selected bacterial infections; they do not cure colds, influenza, or most viral sore throats. Even bacterial infections may not respond if the organism is resistant or the drug does not reach the site effectively. Unnecessary use can cause diarrhea, allergic reactions, drug interactions, and disruption of the microbiome. It also selects for resistant organisms that threaten the patient and community. Preventive antibiotics are appropriate only in defined situations such as selected operations or high-risk exposures. Diagnosis and treatment decisions should be guided by symptoms, examination, testing, and local recommendations.
\item[] \textbf{GEN-071.} \textbf{Myth:} Vaccines cause the disease they are designed to prevent.\newline
\textbf{Status:} False.\newline
\textbf{Fact:} Most vaccines cannot cause the target disease; some live-attenuated vaccines can rarely produce vaccine-associated illness in particular circumstances. Inactivated, subunit, mRNA, and most vector vaccines cannot reproduce the target infection. Live-attenuated vaccines contain weakened organisms and have specific contraindications and rare adverse events, but this is not equivalent to ordinary disease in healthy recipients. Temporary fever, soreness, or fatigue reflects immune activation. Vaccine protection is weighed against the substantially greater risks of many natural infections. Product-specific guidance matters, especially for pregnancy and immune suppression. A suspected serious reaction needs medical assessment rather than abandoning vaccination generally.
\item[] \textbf{GEN-072.} \textbf{Myth:} Vaccines cause autism.\newline
\textbf{Status:} False.\newline
\textbf{Fact:} Large bodies of evidence have found no causal relationship. Large epidemiologic studies across countries and vaccine schedules have not found that routine vaccination causes autism \cite{ref1,ref2}. The original claim was based on discredited research and has not been reproduced by high-quality evidence. Autism is a neurodevelopmental condition with complex genetic and developmental contributors. Temporal coincidence can be mistaken for causation because vaccines are given during early childhood when developmental differences may become noticeable. Vaccination prevents serious infections and does not cause autism. Concerns should be discussed with a qualified clinician using credible evidence.
\item[] \textbf{GEN-073.} \textbf{Myth:} Vaccination weakens the immune system.\newline
\textbf{Status:} False.\newline
\textbf{Fact:} Vaccination trains adaptive immunity against specific pathogens. Vaccines present antigens that allow immune memory to develop without the full risks of the disease. They do not generally exhaust or weaken the immune system. Mild short-term reactions are common signs of immune activation, not immune failure. Protection is specific and may wane or require boosters, so vaccination is not a guarantee against every infection. Delaying or avoiding vaccines leaves people vulnerable to preventable disease. Individual contraindications should be assessed rather than generalized to everyone.
\item[] \textbf{GEN-074.} \textbf{Myth:} Natural infection is always safer than vaccination.\newline
\textbf{Status:} False.\newline
\textbf{Fact:} Infection can carry substantially greater risks than immunization. Natural infection may produce immunity, but it also exposes the person to the pathogen's complications and allows transmission to others. Risks vary by organism, age, pregnancy, immunity, and chronic disease. Vaccines are tested and monitored to provide protection with a lower expected risk than the disease for most recommended populations. Some infections cause long-term disability even after apparent recovery. A vaccine may not prevent every infection but can reduce severe outcomes. ``Natural'' describes route, not safety.
\item[] \textbf{GEN-075.} \textbf{Myth:} Once vaccinated, infection becomes impossible.\newline
\textbf{Status:} False.\newline
\textbf{Fact:} No vaccine provides 100 percent protection. Vaccine effectiveness varies by product, pathogen, strain, time since vaccination, age, immune status, and exposure intensity. Breakthrough infections can occur, but vaccination often reduces severe disease, hospitalization, or transmission. Some vaccines primarily prevent serious illness rather than every mild infection. Updated doses may be recommended as pathogens evolve or immunity changes. A breakthrough case does not prove that the vaccine had no benefit. Prevention also includes ventilation, hygiene, and staying home when ill.
\item[] \textbf{GEN-076.} \textbf{Myth:} Hand sanitizer is always superior to soap and water.\newline
\textbf{Status:} False.\newline
\textbf{Fact:} Soap and water are preferable in several circumstances, including visibly dirty hands and certain pathogens. Alcohol-based sanitizer can rapidly reduce many germs when hands are not visibly dirty and the product contains an effective concentration. Soap and water physically remove dirt, grease, some chemicals, and organisms less susceptible to alcohol. Sanitizer is not a substitute after certain contamination exposures or when hands are visibly soiled. It should be rubbed over all surfaces until dry and kept away from children who may ingest it. Handwashing is especially important before eating and after toileting. Choice should reflect the exposure, not an absolute hierarchy.
\item[] \textbf{GEN-077.} \textbf{Myth:} You can tell whether an infection is viral or bacterial solely from fever severity.\newline
\textbf{Status:} False.\newline
\textbf{Fact:} Temperature alone cannot reliably determine the pathogen type. Both viral and bacterial infections can cause mild or high fever, and some serious infections cause little fever. Age, immune status, prior medicines, timing, and measurement method alter temperature. Diagnosis uses symptoms, examination, exposure, testing, and sometimes imaging or cultures. Antibiotics should not be started solely because a fever is high. Difficulty breathing, altered alertness, dehydration, severe pain, or rapidly worsening illness requires prompt care. Fever is a clinical sign, not a pathogen test.
\item[] \textbf{GEN-078.} \textbf{Myth:} Antibiotics should be taken whenever you have a sore throat.\newline
\textbf{Status:} False.\newline
\textbf{Fact:} Many sore throats are viral; bacterial causes require appropriate evaluation. Viruses cause most sore throats, and antibiotics do not shorten those illnesses. Group A streptococcal infection is one bacterial cause for which testing and treatment may be appropriate. Symptoms alone may not reliably distinguish causes, especially in children. Unnecessary antibiotics can cause allergy, diarrhea, interactions, and resistance. Hydration and symptom relief may help while evaluation is arranged. Severe swallowing difficulty, drooling, breathing problems, or dehydration is urgent.
\item[] \textbf{GEN-079.} \textbf{Myth:} Antibiotic resistance means your body has become resistant to antibiotics.\newline
\textbf{Status:} False.\newline
\textbf{Fact:} It is microorganisms that acquire or possess resistance. Bacteria can develop or acquire mechanisms that allow them to survive an antibiotic. The person may still respond to another drug, dose, or treatment approach, depending on the organism and infection site. Resistance can spread between bacteria and across communities through antibiotic exposure and transmission. Taking antibiotics only when indicated, at the prescribed dose and duration, helps reduce selection pressure. Stopping early or saving leftovers is unsafe, but duration should follow current clinical guidance rather than folklore. Suspected treatment failure needs medical review and sometimes culture.
\item[] \textbf{GEN-080.} \textbf{Myth:} Cracking your knuckles causes arthritis.\newline
\textbf{Status:} False.\newline
\textbf{Fact:} Habitual knuckle cracking has not been convincingly shown to cause osteoarthritis. The popping sound generally results from changes in pressure and gas bubbles within a joint, not bones grinding away. Observational studies have not shown that habitual knuckle cracking causes hand osteoarthritis. It can temporarily affect soft-tissue swelling or grip in some people, but these findings do not establish destructive arthritis. Forceful manipulation can still injure a joint or tendon. Persistent pain, swelling, locking, or loss of motion should not be attributed to harmless cracking. The absence of arthritis risk does not make painful cracking normal.
\item[] \textbf{GEN-081.} \textbf{Myth:} Exercise is bad for arthritic joints.\newline
\textbf{Status:} False.\newline
\textbf{Fact:} Appropriate exercise is an important component of management for many forms of arthritis. Regular, tailored activity can improve strength, joint range, function, mood, and cardiovascular health. It does not mean forcing an acutely inflamed or injured joint through severe pain. Low-impact aerobic exercise, strengthening, and flexibility work can be adjusted to disease and fitness level. Temporary mild soreness is different from escalating pain, swelling, or instability. Physical therapists can adapt loads and technique when movement is difficult. Avoiding all activity often causes deconditioning, which can increase disability.
\item[] \textbf{GEN-082.} \textbf{Myth:} Back pain always means a slipped disc.\newline
\textbf{Status:} False.\newline
\textbf{Fact:} Back pain has numerous possible causes. Common low-back pain can arise from muscles, joints, discs, nerves, or combinations of these structures. Imaging abnormalities are also common in people without pain, so a disc finding does not automatically identify the cause. Most uncomplicated episodes improve with conservative care. Neurologic deficits, fever, cancer history, major trauma, unexplained weight loss, or bowel and bladder changes require prompt assessment. Imaging is usually reserved for selected situations where it will change management. A careful history and examination are more useful than assuming every episode is disc disease.
\item[] \textbf{GEN-083.} \textbf{Myth:} Bed rest is the best treatment for ordinary low-back pain.\newline
\textbf{Status:} False.\newline
\textbf{Fact:} Prolonged bed rest can delay recovery; maintaining appropriate activity is generally preferred. Extended inactivity can reduce muscle conditioning, worsen stiffness, and increase fear of movement. Most people benefit from continuing ordinary activities as tolerated while avoiding movements that clearly aggravate symptoms. Short periods of relative rest may be reasonable during severe acute pain, but prolonged bed rest is not routine treatment. Gradual walking and individualized exercise often support recovery. Severe or progressive neurologic symptoms are a reason for medical assessment, not self-directed bed rest. Return to activity should be paced rather than forced.
\item[] \textbf{GEN-084.} \textbf{Myth:} Weight training is only for young people.\newline
\textbf{Status:} False.\newline
\textbf{Fact:} Properly designed resistance training can benefit older adults substantially. Older adults can improve strength, balance, function, bone loading, and independence with progressive resistance exercise. Programs should begin at an appropriate intensity and account for osteoporosis, arthritis, cardiovascular disease, and prior injury. Supervision is useful when technique or balance is uncertain. Training does not require maximal lifting or gym machines; bands, body weight, and household loads can be effective. Gradual progression and recovery reduce injury risk. New chest pain, fainting, severe shortness of breath, or acute joint pain warrants evaluation.
\item[] \textbf{GEN-085.} \textbf{Myth:} Running inevitably destroys your knees.\newline
\textbf{Status:} False.\newline
\textbf{Fact:} Recreational running has not been shown to inevitably cause knee osteoarthritis. Knee osteoarthritis is influenced by age, prior injury, obesity, genetics, alignment, occupation, and activity history. Recreational running alone is not a guaranteed cause of cartilage failure, and physical activity can support general health. Risk may differ with high training loads, poor recovery, or previous structural injury. Gradual increases, appropriate footwear, strength training, and attention to persistent pain can improve tolerance. A painful or swollen knee should be assessed rather than ignored. Exercise choices should be individualized rather than based on fear of running.
\item[] \textbf{GEN-086.} \textbf{Myth:} Muscle soreness proves a workout was effective.\newline
\textbf{Status:} False.\newline
\textbf{Fact:} An effective workout does not require delayed-onset muscle soreness. Delayed soreness reflects a response to unfamiliar or strenuous loading, especially eccentric exercise. It may occur after a useful session, but it is not a direct measure of strength, fitness, calorie expenditure, or long-term progress. A well-adapted program can produce improvements with little soreness. Severe soreness can limit movement, impair training, and rarely contribute to exertional rhabdomyolysis. Progress is better tracked through performance, function, recovery, and consistent training. Increasing workload gradually is safer than deliberately chasing pain.
\item[] \textbf{GEN-087.} \textbf{Myth:} You can selectively burn abdominal fat by doing sit-ups.\newline
\textbf{Status:} False.\newline
\textbf{Fact:} Local exercise strengthens muscle but does not reliably produce localized fat loss. Situps and other abdominal exercises can improve trunk strength and endurance. They do not direct the body to use fat from the exact area being exercised. Fat loss depends on sustained energy balance, genetics, hormones, sleep, and overall activity. Aerobic and resistance training can reduce total adiposity while improving body composition. Spot-reduction products and extreme diets have not solved this physiological limitation. Waist changes should be evaluated in the context of overall health, not one exercise or measurement.
\item[] \textbf{GEN-088.} \textbf{Myth:} Sweating more means you are burning substantially more fat.\newline
\textbf{Status:} False.\newline
\textbf{Fact:} Sweat primarily regulates body temperature. Sweating reflects heat production, ambient temperature, humidity, clothing, hydration, and individual physiology. A large short-term drop on the scale after heavy sweating is mostly fluid loss, not fat loss. Exercise can increase energy expenditure, but sweat volume is not a reliable measure of it. Deliberate dehydration can reduce performance and cause heat illness, kidney stress, or electrolyte disturbances. Replace fluids according to conditions and duration, with electrolytes when appropriate. Training intensity and long-term energy balance matter more than how wet a workout leaves you.
\item[] \textbf{GEN-089.} \textbf{Myth:} No pain, no gain.\newline
\textbf{Status:} False.\newline
\textbf{Fact:} Pain can indicate injury and is not necessary for effective exercise. Productive exercise may feel challenging, but sharp, escalating, or persistent pain is a warning rather than proof of progress. Pain can arise from tissue injury, excessive load, poor technique, or unrelated disease. Ignoring it can convert a manageable problem into a longer-lasting one. Training can be modified through range of motion, load, speed, frequency, or exercise selection. Sudden severe pain, deformity, neurologic symptoms, chest pain, or fainting needs urgent assessment. Sustainable progression depends on recovery as well as effort.
\item[] \textbf{GEN-090.} \textbf{Myth:} Reading in dim light permanently damages your eyes.\newline
\textbf{Status:} False.\newline
\textbf{Fact:} It may cause temporary eye strain but is not known to permanently damage healthy eyes. Dim lighting can make focusing and reading more tiring, leading to temporary eyestrain, dryness, or headache. These symptoms generally improve with rest, better lighting, blinking, and appropriate breaks. Reading does not use up a finite store of vision or permanently injure a healthy eye. Persistent blur, double vision, pain, flashes, or headaches should not be dismissed as simple strain. Children and adults may need an eye examination if they hold text unusually close or cannot see clearly. Good lighting improves comfort even though it is not a cure for refractive error.
\item[] \textbf{GEN-091.} \textbf{Myth:} Sitting close to a television permanently damages your eyes.\newline
\textbf{Status:} False.\newline
\textbf{Fact:} It does not ordinarily cause permanent ocular damage, although persistent close viewing in children can sometimes signal a vision problem. Modern screens do not emit the kind of radiation that ordinarily damages the eye at normal viewing distances. Sitting close may cause temporary discomfort or indicate that a child cannot see distant details clearly. Children who consistently sit very close, squint, or complain of headaches should receive age-appropriate vision screening. Breaks, blinking, and a comfortable viewing distance can reduce digital eyestrain. Screen use can also displace sleep, activity, and face-to-face interaction, which are separate health concerns. The viewing habit itself is not evidence of permanent eye injury.
\item[] \textbf{GEN-092.} \textbf{Myth:} Wearing glasses makes your eyes progressively weaker.\newline
\textbf{Status:} False.\newline
\textbf{Fact:} Corrective lenses do not weaken the eyes. Glasses bend incoming light so that images focus more accurately on the retina. They do not make the eye dependent on them or cause the underlying refractive error to progress. Vision may change over time because of normal development, aging, disease, or changes in the eye's shape. Removing glasses may make vision seem blurrier by contrast, but that is not damage. An incorrect prescription can cause discomfort or strain, not usually permanent weakening. Regular examinations are appropriate when vision changes or symptoms arise.
\item[] \textbf{GEN-093.} \textbf{Myth:} Cataracts can be permanently removed with eye drops.\newline
\textbf{Status:} False.\newline
\textbf{Fact:} Established cataracts are definitively treated surgically; currently marketed drops have not replaced cataract surgery. A cataract is clouding of the eye's natural lens, usually related to aging but also to injury, medicines, or disease. New drops may be studied, but no over-the-counter product has been established as a replacement for removing a visually significant cataract. Updated glasses and brighter lighting may help early symptoms temporarily. When vision interferes with reading, driving, work, or safety, cataract surgery can remove the cloudy lens and implant an artificial one. Eye drops prescribed around surgery are different from drops claiming to dissolve cataracts. Sudden vision loss needs urgent assessment because it is not typical of an uncomplicated cataract.
\item[] \textbf{GEN-094.} \textbf{Myth:} Acne is simply caused by dirty skin.\newline
\textbf{Status:} False.\newline
\textbf{Fact:} Acne involves follicles, sebum, inflammation, hormones, and microorganisms; poor hygiene is not the basic cause. Acne begins when hair follicles become blocked and inflamed, with contributions from sebum production, hormones, follicular cells, and Cutibacterium acnes. Scrubbing harder does not remove the underlying biology and can worsen irritation. Gentle cleansing can remove sweat, cosmetics, and excess oil without damaging the skin barrier. Evidence-based treatments include topical retinoids, benzoyl peroxide, selected antibiotics, hormonal therapy, or isotretinoin depending on severity. Picking can increase inflammation, scarring, and pigment changes. Persistent or scarring acne deserves professional treatment rather than harsher washing.
\item[] \textbf{GEN-095.} \textbf{Myth:} Toothpaste is a good treatment for acne.\newline
\textbf{Status:} False.\newline
\textbf{Fact:} It can irritate skin and is not a recommended acne treatment. Toothpaste is formulated for teeth, not the skin's barrier or hair follicles. Its flavoring agents, detergents, abrasives, and peroxide-like ingredients can cause irritant or allergic dermatitis. Redness or drying may make a spot look smaller without treating the follicular process. Safer acne ingredients include benzoyl peroxide, salicylic acid, and topical retinoids when appropriate. Treatment should be selected according to acne type, pregnancy status, skin sensitivity, and scarring risk. Wash off toothpaste exposure and seek care for swelling, blistering, or severe irritation.
\item[] \textbf{GEN-096.} \textbf{Myth:} A base tan protects you from harmful ultraviolet radiation.\newline
\textbf{Status:} False.\newline
\textbf{Fact:} A tan represents UV-induced skin injury and provides little protection. Melanin darkening is a response to ultraviolet injury, not a safe protective shield. A pre-existing tan provides only limited filtering and does not prevent DNA damage, photoaging, or skin cancer. Tanning beds also emit harmful ultraviolet radiation and are not a safer preparation for outdoor exposure. Broad-spectrum sunscreen, protective clothing, shade, and avoidance of intense midday exposure reduce risk more effectively. A sunburn is not required for UV damage to occur. Self-tanning products can change color without providing UV protection unless sunscreen is separately used.
\item[] \textbf{GEN-097.} \textbf{Myth:} Sunscreen is needed only on hot, sunny days.\newline
\textbf{Status:} False.\newline
\textbf{Fact:} Harmful UV exposure can occur on cool or cloudy days. Ultraviolet radiation can pass through clouds and does not depend on air temperature. Water, sand, snow, and light surfaces can reflect UV and increase exposure. Broad-spectrum sunscreen is one part of protection for exposed skin when UV levels are meaningful. Hats, tightly woven clothing, shade, sunglasses, and timing also matter. Sunscreen should be applied adequately and reapplied after swimming, sweating, or towel drying. People at high risk of skin cancer should follow individualized dermatology advice.
\item[] \textbf{GEN-098.} \textbf{Myth:} Wounds heal faster when routinely left open to breathe.\newline
\textbf{Status:} False.\newline
\textbf{Fact:} Many wounds heal better in an appropriately moist, protected environment. A clean, lightly moist wound environment can support cell migration and prevent a surface scab from becoming excessively hard. A suitable dressing also protects against friction, contamination, and repeated disruption. Moist does not mean dirty, waterlogged, or sealed over an infected wound. Clean running water or saline, gentle drying of surrounding skin, and a non-stick dressing are commonly appropriate for minor wounds. Deep, gaping, contaminated, bitten, burned, or heavily bleeding wounds need medical assessment. Redness spreading from a wound, pus, fever, or worsening pain may indicate infection.
\item[] \textbf{GEN-099.} \textbf{Myth:} Hydrogen peroxide or alcohol should routinely be poured into every wound.\newline
\textbf{Status:} False.\newline
\textbf{Fact:} These agents can damage viable tissue; uncomplicated wounds are generally cleaned with appropriate irrigation such as clean running water or saline. Strong antiseptics can irritate tissue and impair cells involved in repair when repeatedly applied to an open wound. Gentle irrigation removes dirt and bacteria without the same degree of tissue toxicity. Soap can be used on surrounding skin, but harsh scrubbing inside the wound is unnecessary. A clean dressing and observation are usually adequate for a minor uncomplicated cut. Tetanus vaccination status should be checked after contaminated or puncture wounds. Heavy bleeding, retained foreign material, loss of function, deep wounds, animal or human bites, and infection signs require medical care.

\section{Reproductive, Pediatric, Respiratory, Dermatologic, and Aging Myths}

The following entries extend the registry with additional reproductive, pediatric, respiratory, dermatologic, aging, and treatment myths. The supporting source set includes women's-health and general medical myth guidance; individual claims should still be checked against current local recommendations \cite{ref13,ref11,ref12,ref14}.
\item[] \textbf{REP-001.} \textbf{Myth:} Everyone should routinely undergo huge panels of blood tests.\newline
\textbf{Status:} False.\newline
\textbf{Fact:} Indiscriminate testing can produce false positives and unnecessary investigations. Screening is useful when a test has a reasonable balance of benefit and harm for a defined population. Testing without a clinical question increases the chance of borderline or false-positive results. Those results can lead to anxiety, repeat procedures, incidental findings, and overtreatment. Recommended screening depends on age, sex, pregnancy, risk factors, symptoms, and local guidelines. A focused history and examination help determine which tests are appropriate. People with concerning symptoms should seek evaluation rather than ordering an indiscriminate panel.
\item[] \textbf{REP-002.} \textbf{Myth:} The vagina needs internal washing or douching.\newline
\textbf{Status:} False.\newline
\textbf{Fact:} Routine internal cleansing is unnecessary and can disrupt the vaginal environment. The vagina has self-regulating secretions and a microbial environment that helps maintain health. Douching can irritate tissue and alter vaginal flora, and it has been associated with problems such as bacterial vaginosis and pelvic infection. Washing the external vulva gently with water or mild unscented cleanser is generally sufficient. Fragrance products can also trigger irritation. Odor, itching, pain, bleeding, or unusual discharge warrants assessment rather than masking symptoms. Douching does not prevent pregnancy or sexually transmitted infections.
\item[] \textbf{REP-003.} \textbf{Myth:} Vaginal discharge always indicates infection.\newline
\textbf{Status:} False.\newline
\textbf{Fact:} Physiological discharge is normal. Normal discharge can vary in amount, thickness, and appearance across the menstrual cycle, pregnancy, and sexual arousal. Infection becomes more likely when discharge is accompanied by new odor, itching, burning, pelvic pain, sores, or bleeding. These symptoms are not specific enough to diagnose the cause by appearance alone. Bacterial vaginosis, candidiasis, trichomoniasis, cervicitis, and noninfectious irritation require different management. Self-treating repeatedly can delay the correct diagnosis. Testing and examination are appropriate when symptoms are new, persistent, recurrent, or associated with pregnancy.
\item[] \textbf{REP-004.} \textbf{Myth:} Pregnancy cannot occur during menstruation.\newline
\textbf{Status:} False.\newline
\textbf{Fact:} It is possible, depending on cycle timing and sperm survival. Ovulation does not occur on an identical calendar day for every person or every cycle. Sperm can remain viable in the reproductive tract for several days, and bleeding is not always a true menstrual period. Short or irregular cycles can bring ovulation relatively soon after bleeding begins. Therefore, intercourse during bleeding can result in pregnancy, although probability varies. Contraception should be used according to the person's pregnancy-prevention goals, not assumed safe days. Unusual bleeding should be evaluated when it is heavy, painful, or recurrent.
\item[] \textbf{REP-005.} \textbf{Myth:} Pregnancy cannot occur the first time someone has intercourse.\newline
\textbf{Status:} False.\newline
\textbf{Fact:} It can. Pregnancy depends on sperm reaching an egg during a fertile window, not on how many times intercourse has occurred. A first sexual encounter can coincide with ovulation and result in pregnancy. Withdrawal is less reliable than consistent contraception because pre-ejaculatory fluid may contain sperm and timing is difficult. Condoms also reduce sexually transmitted infection risk when used correctly and consistently. Emergency contraception may be an option after unprotected intercourse, depending on timing and local availability. Pregnancy testing is appropriate after a missed period or as directed by the test instructions.
\item[] \textbf{REP-006.} \textbf{Myth:} Oral contraceptives permanently cause infertility.\newline
\textbf{Status:} False.\newline
\textbf{Fact:} Fertility generally returns after discontinuation. Combined and progestin-only pills prevent pregnancy temporarily by altering ovulation, cervical mucus, or the uterine environment. Most people resume ovulation and fertility soon after stopping, although timing varies. A delay in conception is not proof that the pill caused infertility; age, underlying conditions, partner factors, and chance are also important. Some pills can make pre-existing cycle irregularity less visible while they are being used. Persistent absence of periods or difficulty conceiving deserves an appropriate evaluation. Contraceptive choice should consider medical history and personal priorities.
\item[] \textbf{REP-007.} \textbf{Myth:} Menopause is a disease.\newline
\textbf{Status:} False.\newline
\textbf{Fact:} It is a normal biological transition. Menopause is defined retrospectively after twelve months without menstruation when another cause is not present. The transition can cause hot flashes, sleep disruption, vaginal symptoms, mood changes, and changes in bleeding patterns. Symptoms may be mild, severe, brief, or prolonged, and treatment can be considered when they impair quality of life. Menopause does not remove the need for contraception until pregnancy is no longer possible and does not eliminate cardiovascular or bone health needs. Bleeding after menopause is not automatically normal and should be evaluated. Care should be individualized rather than dismissive or pathologizing.
\item[] \textbf{REP-008.} \textbf{Myth:} Cervical cancer is hereditary in the same way as many inherited cancer syndromes.\newline
\textbf{Status:} False.\newline
\textbf{Fact:} Persistent high-risk HPV infection is the principal cause of cervical cancer. Most cervical cancers arise after persistent infection with a high-risk human papillomavirus type. Host immune response, smoking, immune suppression, screening history, and other factors influence progression. A family history may affect risk or screening discussions, but it does not make cervical cancer a simple inherited syndrome in the usual sense. HPV vaccination and evidence-based cervical screening prevent many cancers by stopping infection or detecting precancer \cite{ref9}. A positive HPV test does not mean cancer is present. Follow-up recommendations depend on the result and prior screening history.
\item[] \textbf{REP-009.} \textbf{Myth:} Pregnant women should literally ``eat for two.''\newline
\textbf{Status:} False.\newline
\textbf{Fact:} Energy requirements increase, but nowhere near doubling for most pregnancies. Pregnancy increases nutritional needs for protein, iron, folate, iodine, and other nutrients, but calorie needs change by trimester and individual circumstances. Excessive weight gain can increase risks such as gestational diabetes, hypertension, cesarean delivery, and fetal complications. Inadequate intake can also be harmful, so intentional weight loss diets are not generally appropriate without specialist guidance. Recommended weight gain depends partly on pre-pregnancy body mass index and should be discussed with prenatal clinicians. Food quality and prenatal supplements matter more than doubling portions. Severe nausea, vomiting, or restrictive eating needs medical support.
\item[] \textbf{REP-010.} \textbf{Myth:} Exercise is inherently dangerous during pregnancy.\newline
\textbf{Status:} False.\newline
\textbf{Fact:} Appropriate physical activity is recommended in most uncomplicated pregnancies. For many pregnant people, moderate activity supports cardiovascular fitness, mood, sleep, and metabolic health. Activity should be adapted to symptoms, fitness, pregnancy complications, and clinician advice. High-risk conditions such as significant bleeding, placenta problems, preterm labor risk, or certain cardiac disease may require restrictions. Hydration, temperature control, and avoiding fall or abdominal-trauma risks are important. Warning signs such as bleeding, fluid leakage, chest pain, fainting, painful contractions, or severe shortness of breath require prompt assessment. Pregnancy exercise is individualized, not universally forbidden.
\item[] \textbf{REP-011.} \textbf{Myth:} Heartburn during pregnancy reliably predicts a baby with lots of hair.\newline
\textbf{Status:} False.\newline
\textbf{Fact:} It is not a clinically useful predictor. Pregnancy heartburn commonly results from hormonal relaxation of the lower esophageal sphincter and pressure from the growing uterus. Fetal hair growth and maternal reflux are not a dependable diagnostic pair. Some research has found an association in particular samples, but it is not accurate enough to predict an individual baby's hair. Heartburn can be reduced by smaller meals, avoiding personal triggers, remaining upright after eating, and clinician-approved treatment. Severe pain, vomiting blood, difficulty swallowing, or weight loss requires medical review. Ultrasound and postnatal observation, not heartburn, answer fetal or newborn questions.
\item[] \textbf{REP-012.} \textbf{Myth:} The shape of the abdomen predicts the baby's sex.\newline
\textbf{Status:} False.\newline
\textbf{Fact:} It does not. Abdominal appearance reflects uterine position, fetal position, muscle tone, body shape, gestational age, and number of prior pregnancies. These factors vary widely and do not encode fetal sex. Reliable determination, when desired and appropriate, uses validated prenatal testing or imaging rather than visual folklore. No method is completely risk-free or suitable for every pregnancy. Sex prediction should not replace clinical assessment of fetal growth or maternal symptoms. Cultural assumptions about appearance can also create unnecessary anxiety or pressure.
\item[] \textbf{REP-013.} \textbf{Myth:} Fetal heart rate reliably predicts fetal sex.\newline
\textbf{Status:} False.\newline
\textbf{Fact:} It does not. Fetal heart rate changes with gestational age, movement, sleep, medications, maternal conditions, and the immediate physiologic state. Average rates overlap substantially across fetal sexes. A single reading is used to assess well-being in context, not to determine sex. Abnormal or persistently concerning patterns require obstetric evaluation regardless of sex. Home devices are not a substitute for clinical monitoring and can create false reassurance or alarm. Sex-related claims based on a heart-rate cutoff are not clinically reliable.
\item[] \textbf{REP-014.} \textbf{Myth:} Breastfeeding completely prevents pregnancy.\newline
\textbf{Status:} False.\newline
\textbf{Fact:} Lactational amenorrhea works as contraception only under specific conditions. Lactational amenorrhea can be effective when the baby is under six months, menstruation has not returned, and feeding is exclusive or nearly exclusive at frequent intervals. Effectiveness falls when feeds are spaced, formula or solids are added, pumping patterns change, or bleeding returns. Ovulation can occur before the first postpartum period. People who wish to avoid pregnancy should discuss a compatible contraceptive method with a clinician. Breastfeeding remains valuable even when contraception is also used. Pregnancy concerns, pain, or low milk supply warrant support rather than assumptions.
\item[] \textbf{REP-015.} \textbf{Myth:} Infant formula and breast milk are biologically identical.\newline
\textbf{Status:} False.\newline
\textbf{Fact:} Formula provides nutrition but does not reproduce all biological components of human milk. Commercial formula is regulated to provide appropriate nutrients for infant growth when prepared correctly. Human milk also contains living cells, antibodies, enzymes, hormones, and other components that vary over time and are not fully duplicated by formula. Feeding decisions are personal and may be shaped by health, access, work, medication, and preference. Formula feeding is not a moral failure and can be done safely with correct preparation. Breastfeeding challenges deserve practical support without shaming. Infants need adequate nutrition, safe handling, and responsive care by whichever feeding method is used.
\item[] \textbf{REP-016.} \textbf{Myth:} Babies need additional water during their first six months.\newline
\textbf{Status:} False.\newline
\textbf{Fact:} Exclusively breastfed healthy infants generally do not need additional water. Breast milk or properly prepared infant formula supplies both fluid and nutrition for most healthy infants before complementary foods begin. Extra water can displace calories and nutrients and, in large amounts, disturb sodium balance. Hot weather does not automatically change this recommendation; feeding frequency may be adjusted with professional guidance. Formula must be mixed exactly as directed because dilution can cause dangerous electrolyte problems. Infants with fever, vomiting, diarrhea, kidney disease, or other medical needs require clinician-specific advice. Signs of dehydration need prompt evaluation.
\item[] \textbf{REP-017.} \textbf{Myth:} Teething routinely causes high fever.\newline
\textbf{Status:} False.\newline
\textbf{Fact:} Teething may cause mild symptoms but should not be assumed to explain significant fever. Teething can coincide with drooling, gum discomfort, irritability, and a small temperature change. A high fever, marked lethargy, breathing difficulty, persistent vomiting, or significant diarrhea is not safely explained by teething alone. Infants are also exposed to infections during the same developmental period, which can create misleading timing. Cooling measures and clinician-approved analgesia may help gum discomfort; unsafe gels and necklaces should be avoided. The child's age, appearance, hydration, and measured temperature guide evaluation. Persistent or concerning symptoms require pediatric advice.
\item[] \textbf{REP-018.} \textbf{Myth:} Picking up a crying newborn will inevitably spoil the baby.\newline
\textbf{Status:} False.\newline
\textbf{Fact:} Responsive caregiving does not spoil a newborn. Newborns cry because of hunger, discomfort, fatigue, temperature, overstimulation, and the need for regulation. Prompt, safe comforting helps build predictable responses and does not create adult-style manipulation. Holding, feeding, rocking, and skin-to-skin contact can support bonding and physiologic regulation. Caregivers also need rest and practical support; placing a baby safely on the back in an empty crib is appropriate when overwhelmed. Never shake a baby. A sudden change in crying, weak feeding, fever, breathing difficulty, or unusual sleepiness needs medical evaluation.
\item[] \textbf{REP-019.} \textbf{Myth:} Children are simply miniature adults medically.\newline
\textbf{Status:} False.\newline
\textbf{Fact:} Physiology, disease patterns, and drug dosing differ substantially. Children differ from adults in airway size, fluid balance, metabolism, immune development, communication, and normal vital signs. Diseases can present differently, and some treatments are age- or weight-dependent. Adult medicines and doses should not be shared without professional instruction. Growth and development are central parts of pediatric assessment. A child's inability to describe symptoms does not mean symptoms are absent. Pediatric care uses age-appropriate examination, dosing, prevention, and safety planning.
\item[] \textbf{REP-020.} \textbf{Myth:} A chubby child is necessarily a healthy child.\newline
\textbf{Status:} False.\newline
\textbf{Fact:} Growth should be assessed using appropriate age- and sex-specific growth measures. Healthy growth cannot be judged reliably by appearance alone. Clinicians consider height, weight, body-mass-index percentile, growth trajectory, puberty, nutrition, activity, family history, and medical conditions. Excess adiposity can increase risks for hypertension, sleep apnea, insulin resistance, orthopedic problems, and stigma, while inadequate nutrition can also impair growth. Weight-focused criticism can worsen eating and mental-health outcomes. Families benefit from supportive, behavior-centered guidance rather than shame. Unexpected growth acceleration, slowdown, or symptoms deserves pediatric evaluation.
\item[] \textbf{REP-021.} \textbf{Myth:} Antibiotics help children recover from every fever faster.\newline
\textbf{Status:} False.\newline
\textbf{Fact:} They do not treat viral infections. Fever is a sign of many infections and inflammatory conditions, not a diagnosis that automatically calls for antibiotics. Most childhood respiratory infections are viral and improve with fluids, rest, and symptom care. Antibiotics do not shorten viral illness and can cause diarrhea, allergy, or resistant bacteria. Some bacterial infections do benefit from prompt treatment, so persistent or concerning symptoms need assessment. Age, appearance, breathing, hydration, pain, and duration matter more than fever alone. Caregivers should use prescribed medicines exactly as directed and not share leftovers.
\item[] \textbf{REP-022.} \textbf{Myth:} Children cannot have high blood pressure.\newline
\textbf{Status:} False.\newline
\textbf{Fact:} Pediatric hypertension occurs. Children can develop elevated blood pressure from primary risk factors, kidney disease, endocrine conditions, medications, sleep apnea, or other causes. Interpretation depends on age, sex, and height, so adult cutoffs are not always appropriate. A single anxious or technically poor measurement does not establish chronic hypertension. Repeated, correctly sized-cuff measurements are usually needed. Persistent hypertension can affect the heart, kidneys, brain, and blood vessels over time. Pediatric evaluation should identify both the cause and any modifiable risks.
\item[] \textbf{REP-023.} \textbf{Myth:} Children cannot develop type 2 diabetes.\newline
\textbf{Status:} False.\newline
\textbf{Fact:} They can. Type 2 diabetes is increasingly recognized in adolescents and, less commonly, younger children. Risk is influenced by genetics, insulin resistance, body composition, puberty, medications, and social determinants, not simply by personal choice. Some children with diabetes have normal or lower body weight, and type 1 diabetes can occur at any weight. Symptoms such as excessive thirst, frequent urination, fatigue, weight change, or recurrent infections require evaluation. Early diagnosis can reduce acute and long-term complications. Treatment is individualized and should avoid blame.
\item[] \textbf{REP-024.} \textbf{Myth:} Vaccines overload a child's immune system.\newline
\textbf{Status:} False.\newline
\textbf{Fact:} The immune system routinely handles vastly more antigens than those contained in vaccines. Vaccines expose the immune system to selected antigens that train protection against specific pathogens. Modern vaccines often contain fewer total antigens than older formulations despite covering more diseases. Infants encounter countless environmental microbes every day through breathing, feeding, and skin contact. Recommended schedules are designed around when protection is needed and when immune responses are reliable. Temporary fever or soreness can occur, while serious reactions are rare. Delaying vaccines prolongs vulnerability without providing a proven immune benefit.
\item[] \textbf{REP-025.} \textbf{Myth:} Fever always causes brain damage in children.\newline
\textbf{Status:} False.\newline
\textbf{Fact:} Ordinary infection-related fever generally does not cause brain damage. Fever is a regulated immune response and is different from dangerous hyperthermia caused by heat exposure or certain toxic states. Typical fevers from infection do not reach temperatures that injure the brain. Children can nevertheless become dehydrated, uncomfortable, or seriously ill from the underlying infection. Fever medicines are used for comfort and should be dosed by weight according to instructions. Infants, children with altered alertness, breathing difficulty, stiff neck, dehydration, or prolonged fever need medical assessment. A number alone does not determine severity \cite{ref8}.
\item[] \textbf{REP-026.} \textbf{Myth:} Every febrile seizure means epilepsy.\newline
\textbf{Status:} False.\newline
\textbf{Fact:} Febrile seizures and epilepsy are different entities. Febrile seizures occur in some young children during fever and are usually brief and generalized. Most children with a simple febrile seizure do not develop epilepsy. Risk is higher in selected circumstances, including prolonged, focal, recurrent, or atypical events and certain family or developmental histories. During a seizure, protect the child from injury, place nothing in the mouth, and time the event. Emergency care is needed for prolonged or repeated seizures, breathing problems, injury, or a first uncertain event. Follow-up depends on the seizure features and examination.
\item[] \textbf{REP-027.} \textbf{Myth:} Aspirin is an ordinary fever medicine for children.\newline
\textbf{Status:} False.\newline
\textbf{Fact:} Aspirin is generally avoided in children and adolescents with certain viral illnesses because of the association with Reye syndrome. Aspirin can cause serious adverse effects in children, including gastrointestinal bleeding and toxicity. Its use during or after some viral infections has been associated with Reye syndrome, a rare but severe illness affecting the liver and brain. Caregivers should not substitute adult aspirin for pediatric fever medicines without medical instruction. Acetaminophen or ibuprofen may be appropriate in specific ages and situations, but dosing must follow weight-based guidance and contraindications. Never give multiple products containing the same ingredient. Seek advice when fever is severe, prolonged, or accompanied by concerning symptoms \cite{ref8}.
\item[] \textbf{REP-028.} \textbf{Myth:} Asthma is primarily a psychological disease.\newline
\textbf{Status:} False.\newline
\textbf{Fact:} Asthma is a chronic inflammatory airway disease. Asthma involves variable airway narrowing, inflammation, and increased bronchial responsiveness. Stress, strong emotions, or anxiety can trigger or amplify symptoms, but they do not make asthma imaginary. Wheeze, cough, chest tightness, and breathlessness have multiple possible causes and require appropriate assessment. Inhaled anti-inflammatory treatment and bronchodilators reduce risk and symptoms for many patients. Avoiding stigma helps people use medicines and action plans consistently. Severe breathlessness, inability to speak normally, cyanosis, or poor response to rescue medicine is an emergency.
\item[] \textbf{REP-029.} \textbf{Myth:} People with asthma should not exercise.\newline
\textbf{Status:} False.\newline
\textbf{Fact:} With appropriate management, physical activity is encouraged. Exercise can provoke bronchoconstriction in some people, but regular activity improves fitness and quality of life. A clinician can assess control, teach inhaler technique, and provide a pre-exercise or rescue plan when indicated. Warm-up, environmental adjustments, and treatment of allergic or airway triggers may help. Avoiding all exercise can worsen deconditioning and confidence. Symptoms that are frequent, severe, or absent during ordinary activities may indicate poor control. Exercise plans should be adapted rather than abandoned.
\item[] \textbf{REP-030.} \textbf{Myth:} Asthma inhalers are addictive.\newline
\textbf{Status:} False.\newline
\textbf{Fact:} Standard asthma inhalers do not cause addiction. Inhaled corticosteroids reduce airway inflammation and are not addictive drugs. Reliever bronchodilators can produce physiologic effects such as tremor or a fast heartbeat, but these are not the compulsive behavior or withdrawal pattern of addiction. Needing a reliever frequently may signal uncontrolled asthma rather than dependence. Stopping controller therapy abruptly can allow inflammation and symptoms to return. Correct technique and a written action plan improve treatment safety. Medication changes should be discussed with a clinician, especially after severe attacks.
\item[] \textbf{REP-031.} \textbf{Myth:} Using an inhaler means your asthma has become severe.\newline
\textbf{Status:} False.\newline
\textbf{Fact:} Inhaled therapy is a standard route of asthma treatment across severity levels. Inhalers deliver medicines directly to the airways and are used for intermittent as well as persistent asthma. A controller inhaler may be prescribed because inflammation needs prevention, not because a person is in immediate danger. Frequent use of a reliever inhaler can indicate poor control and should prompt review. Severity and control are related but not identical clinical concepts. Correct technique, adherence, triggers, and lung function help guide treatment. Inhaled therapy should not be stopped solely because the device looks serious.
\item[] \textbf{REP-032.} \textbf{Myth:} Tuberculosis is always obvious from symptoms.\newline
\textbf{Status:} False.\newline
\textbf{Fact:} Latent TB infection produces no symptoms, and active disease can initially be subtle. Latent tuberculosis infection means immune control of the bacteria without symptoms or contagious pulmonary disease. Active TB can cause cough, fever, night sweats, weight loss, fatigue, or coughing blood, but not every person has the full pattern. Symptoms may develop gradually and can resemble other illnesses. Testing is guided by exposure, risk, symptoms, and local public-health recommendations. Untreated pulmonary TB can spread to others, so suspected cases require timely evaluation. Treatment differs for latent infection and active disease and must be completed as prescribed.
\item[] \textbf{REP-033.} \textbf{Myth:} Pneumonia occurs only in cold weather.\newline
\textbf{Status:} False.\newline
\textbf{Fact:} It can occur at any time of year. Pneumonia is infection or inflammation of lung tissue caused by viruses, bacteria, fungi, aspiration, and other processes. Seasonal circulation changes exposure patterns but does not make pneumonia a winter-only illness. Risk is increased by age, chronic disease, immune suppression, smoking, aspiration, and some healthcare exposures. Symptoms can include cough, fever, breathlessness, chest pain, confusion, or low oxygen, and older adults may lack fever. Vaccination and hand hygiene reduce some causes but do not prevent every case. Breathing difficulty, blue lips, confusion, or severe weakness requires urgent care.
\item[] \textbf{REP-034.} \textbf{Myth:} All pneumonia is bacterial.\newline
\textbf{Status:} False.\newline
\textbf{Fact:} Viruses, bacteria, fungi, and other causes exist. Pneumonia has multiple causes, and the responsible organism cannot always be identified from symptoms alone. Antibiotics help selected bacterial cases but do not treat viral or fungal disease. Some viral pneumonias can be severe, particularly in infants, older adults, pregnant people, and those with chronic or immune conditions. Testing and imaging may help, though treatment often begins before a definitive organism is known. Antimicrobial choice should reflect likely cause, severity, resistance, and patient factors. Worsening breathing or oxygenation is more urgent than knowing the label immediately.
\item[] \textbf{REP-035.} \textbf{Myth:} Smoking only damages the lungs.\newline
\textbf{Status:} False.\newline
\textbf{Fact:} It increases cardiovascular, cancer, reproductive, vascular, and many other risks. Tobacco smoke injures blood vessels, promotes inflammation and clotting, and exposes the body to carcinogens. Smoking contributes to heart disease, stroke, chronic lung disease, multiple cancers, infertility, pregnancy complications, osteoporosis, and impaired wound healing. Secondhand smoke also harms people nearby. Risk rises with exposure, but there is no safe tobacco product or need to wait for symptoms before quitting. Benefits begin soon after cessation and continue over time. Evidence-based cessation support can include counseling, nicotine replacement, and prescribed medicines.
\item[] \textbf{REP-036.} \textbf{Myth:} Occasional smoking is harmless.\newline
\textbf{Status:} False.\newline
\textbf{Fact:} Even low levels of tobacco exposure increase health risks. Cardiovascular risk does not fall in a perfectly linear way with cigarette count, so a few cigarettes are not equivalent to zero. Occasional smoking can maintain nicotine dependence and expose the cardiovascular and respiratory systems to toxic chemicals. It also increases the chance of progressing to regular use and exposes others to secondhand smoke. Cutting down may be a step toward quitting, but complete cessation provides the greatest health benefit. Vaping or switching products should not be assumed risk-free. Support is available even for people who smoke infrequently.
\item[] \textbf{REP-037.} \textbf{Myth:} Hookah is safer than cigarettes because smoke passes through water.\newline
\textbf{Status:} False.\newline
\textbf{Fact:} Water does not remove the major health hazards of tobacco smoke. Hookah smoke can contain nicotine, carbon monoxide, fine particles, and carcinogens. A session may last much longer than a cigarette and can result in substantial smoke exposure. Charcoal adds carbon monoxide and other toxic substances. Sharing mouthpieces can also transmit infections unless appropriate hygiene is used. Flavor, cooling, or water filtration does not make tobacco safe. Hookah users who want to stop should receive the same kind of cessation support offered to cigarette users.
\item[] \textbf{REP-038.} \textbf{Myth:} Shaving makes hair grow back thicker.\newline
\textbf{Status:} False.\newline
\textbf{Fact:} Shaving does not change follicle thickness; blunt-cut hairs can temporarily appear coarser. Shaving cuts hair at the skin surface and does not alter the follicle, pigment, growth rate, or number of hairs. Newly emerging stubble has a blunt tip that can feel rougher and look darker than a tapered, sun-lightened hair. Hormonal changes, age, medicines, and skin conditions can change hair growth independently of shaving. Repeated shaving can cause irritation, ingrown hairs, or folliculitis if technique is poor. Persistent new facial or body hair growth may need evaluation for endocrine or medication causes.
\item[] \textbf{REP-039.} \textbf{Myth:} Cutting hair makes it grow faster.\newline
\textbf{Status:} False.\newline
\textbf{Fact:} Hair growth occurs at the follicle, not the cut end. Hair shafts are nonliving structures produced by follicles beneath the skin. Trimming removes damaged or split ends and can make hair look fuller or healthier, but it does not accelerate follicle growth. Growth rate varies with genetics, age, hormones, nutrition, and health. Sudden shedding or patchy loss is not corrected by more frequent cutting. Gentle care, adequate nutrition, and diagnosis of underlying disease are more relevant. A dermatologist can assess persistent shedding, scarring, or scalp inflammation.
\item[] \textbf{REP-040.} \textbf{Myth:} Dandruff necessarily means poor hygiene.\newline
\textbf{Status:} False.\newline
\textbf{Fact:} It is associated with scalp biology and conditions such as seborrheic dermatitis. Dandruff involves flaking and inflammation influenced by sebum, skin-barrier factors, and the scalp yeast Malassezia. It is not proof that someone is dirty or does not wash. Shampoo frequency and product choice can affect symptoms, but excessive scrubbing can irritate the scalp. Antidandruff shampoos with ingredients such as ketoconazole, zinc pyrithione, selenium sulfide, or salicylic acid may help when used correctly. Thick plaques, hair loss, pain, or failure to improve may indicate psoriasis, dermatitis, or infection. Diagnosis and treatment should be tailored to the scalp condition.
\item[] \textbf{REP-041.} \textbf{Myth:} Hair loss always comes from stress.\newline
\textbf{Status:} False.\newline
\textbf{Fact:} Genetics, hormones, autoimmune disease, medications, nutrition, and numerous other factors can contribute. Stress can trigger temporary shedding, often called telogen effluvium, after a delay of several weeks or months. It is only one possible cause of hair loss. Pattern hair loss, alopecia areata, thyroid disease, iron deficiency, medication effects, infection, and scarring disorders require different evaluations and treatments. Blaming stress alone can delay diagnosis of a treatable condition. Sudden patchy loss, scalp inflammation, scarring, or loss of eyebrows warrants assessment. A careful history and scalp examination are more useful than commercial hair claims.
\item[] \textbf{REP-042.} \textbf{Myth:} Acne should be aggressively scrubbed away.\newline
\textbf{Status:} False.\newline
\textbf{Fact:} Excessive scrubbing can worsen irritation. Acne forms within follicles, so abrasive scrubbing cannot remove its underlying drivers. Harsh friction can damage the skin barrier, increase inflammation, and worsen redness or post-inflammatory pigmentation. Gentle cleansing once or twice daily is usually sufficient. Evidence-based topical treatments work through mechanisms such as reducing follicular plugging, inflammation, or bacterial load. Products should be introduced gradually to limit irritation. Severe, painful, or scarring acne needs professional care rather than stronger scrubbing.
\item[] \textbf{REP-043.} \textbf{Myth:} Chocolate invariably causes acne.\newline
\textbf{Status:} False.\newline
\textbf{Fact:} The diet-acne relationship is more complex than a single food causing acne in everyone. Acne is driven mainly by follicular biology, hormones, sebum, inflammation, and genetics. Some studies suggest that high-glycemic diets or certain dairy patterns may influence acne in susceptible people, but evidence does not support blaming every chocolate food equally. Chocolate products also differ in sugar, milk, and fat content, making the category biologically imprecise. Restricting a single food is not a substitute for proven treatment. People can track personal triggers without adopting unnecessarily restrictive diets. Persistent acne should be treated according to severity and scarring risk.
\item[] \textbf{REP-044.} \textbf{Myth:} A wound itching means it is infected.\newline
\textbf{Status:} False.\newline
\textbf{Fact:} Itching can occur during normal healing. Healing skin can itch as inflammation settles and new tissue forms. Mild, localized itching without worsening pain, spreading redness, pus, fever, or loss of function is often compatible with normal repair. Scratching can reopen the wound or introduce bacteria, so a protective dressing and gentle care are useful. Infection is assessed from the overall pattern, not itch alone. Increasing warmth, swelling, drainage, red streaks, or systemic illness requires medical attention. People with diabetes, poor circulation, or immune suppression need a lower threshold for review.
\item[] \textbf{REP-045.} \textbf{Myth:} All moles are harmless.\newline
\textbf{Status:} False.\newline
\textbf{Fact:} Most are benign, but melanoma can arise in or resemble a mole. Many moles are stable and noncancerous, but melanoma can develop in an existing mole or appear as a new lesion. Warning changes include asymmetry, irregular borders, varied color, increasing diameter, evolution, bleeding, or a lesion that looks unlike the person's others. These signs are prompts for assessment, not a diagnosis by themselves. Self-examination can help identify change, while photographs assist comparison. Avoid cutting or burning a mole at home. A clinician may use dermoscopy or biopsy when indicated.
\item[] \textbf{REP-046.} \textbf{Myth:} Darker skin cannot develop skin cancer.\newline
\textbf{Status:} False.\newline
\textbf{Fact:} Skin cancer occurs across skin tones. Melanin provides some protection against ultraviolet injury but does not eliminate cancer risk. Melanoma and other skin cancers can occur in people with richly pigmented skin and may be diagnosed later because the risk is underestimated. Acral sites such as palms, soles, and nails deserve attention, as do changing or non-healing lesions anywhere. Sun protection remains useful, although it should be combined with awareness that not all cancers are caused by sun exposure. New, changing, bleeding, or persistent lesions should be examined. Screening advice should reflect individual risk and skin history.
\item[] \textbf{REP-047.} \textbf{Myth:} Sunburn turns into a healthy tan.\newline
\textbf{Status:} False.\newline
\textbf{Fact:} Both sunburn and tanning represent ultraviolet skin injury; a tan is not a healthy protective response. Sunburn is an acute inflammatory response to UV damage, while tanning reflects increased melanin after injury. Both processes can damage cellular DNA and contribute to premature aging and skin cancer. A later tan does not reverse the initial injury or provide reliable protection. Repeated burns are particularly harmful, but tanning without visible redness is not risk-free. Shade, clothing, hats, and broad-spectrum sunscreen reduce exposure. Painful blistering burns, extensive burns, or burns with systemic symptoms need medical care.
\item[] \textbf{REP-048.} \textbf{Myth:} Older adults cannot learn new skills.\newline
\textbf{Status:} False.\newline
\textbf{Fact:} Learning and neuroplasticity persist throughout life. Aging can slow processing speed or make some forms of recall less efficient, but the brain retains capacity for learning and adaptation. Practice, meaningful repetition, sleep, physical activity, hearing and vision correction, and social engagement support performance. Older adults can learn languages, technologies, music, exercise skills, and new routines. Difficulty is not proof of inability or dementia. Sudden or progressive cognitive change that affects daily function needs assessment. Teaching should allow adequate time and use accessible strategies rather than age-based assumptions.
\item[] \textbf{REP-049.} \textbf{Myth:} Dietary supplements can reliably stop normal aging.\newline
\textbf{Status:} False.\newline
\textbf{Fact:} Normal aging has no proven supplement remedy. Aging unfolds across many organ systems and molecular pathways, so no single nutrient, hormone, or extract can be expected to halt it. Preparations sold as longevity boosters sometimes contain ingredients that interact with medicines or harm the liver and kidneys at high doses. Correcting a documented vitamin or mineral deficiency is legitimate care, but it is not the same as reversing age itself. Randomized trials measuring outcomes such as survival, function, or disease prevention, rather than testimonials, are the standard of proof. Habits with stronger supporting evidence include regular activity, adequate sleep, sensible nutrition, vaccination, social connection, and control of blood pressure, glucose, and lipids. Anyone considering supplements alongside chronic disease or regular medicines should review them with a clinician.
\item[] \textbf{REP-050.} \textbf{Myth:} If a treatment has been used for centuries, that proves it works.\newline
\textbf{Status:} False.\newline
\textbf{Fact:} Historical use is not a substitute for controlled evidence of efficacy and safety. Long use may indicate acceptability or cultural importance, but ineffective practices can persist through tradition, selective memory, and lack of alternatives. A treatment must be evaluated for a defined condition, dose, benefit, harms, interactions, and quality. Modern trials can reveal effects that historical experience misses, including placebo effects and delayed natural recovery. Traditional practices may contain useful interventions, but each claim should be tested rather than accepted wholesale or rejected wholesale. Natural products can also be toxic or contaminated. Patients should disclose all treatments so risks and interactions can be assessed. The legacy numbering of these entries is retained temporarily for source traceability. A final evidence audit must assign stable identifiers only after duplicate removal, provenance checks, and confirmation that no source entry is missing.

\section{Kidney, Nutrition, Cardiovascular, Eye, Breast, Sexual, and First-Aid Myths}
\item[] \textbf{KID-001.} \textbf{Myth:} Drinking enormous quantities of water cleans the kidneys.\newline
\textbf{Status:} False.\newline
\textbf{Fact:} Excessive water intake can cause dangerous hyponatremia. The kidneys continuously filter blood and regulate water and electrolytes without a special flush. Drinking beyond thirst does not scrub away kidney disease or prevent every stone. Excessive intake can dilute blood sodium and cause headache, confusion, seizures, or death. Fluid needs vary with climate, activity, diet, kidney function, heart disease, and medicines. People with a clinician-prescribed fluid limit should follow it. Hydration should be adequate and individualized, not forced to an arbitrary volume.
\item[] \textbf{KID-002.} \textbf{Myth:} Clear urine always means perfect hydration.\newline
\textbf{Status:} False.\newline
\textbf{Fact:} Urine color is influenced by fluid intake and other factors and is not a complete hydration assessment. Pale-yellow urine often accompanies adequate hydration, but completely clear urine may reflect excessive intake. Color is affected by foods, vitamins, medicines, blood, infection, and kidney concentrating ability. Thirst, symptoms, weather, exercise, and medical conditions also matter. People with kidney, heart, liver, or endocrine disorders may not follow generic hydration rules. Persistent dark, bloody, cloudy, or unusually foamy urine needs assessment. A single urine color cannot prove that every body system is healthy.
\item[] \textbf{KID-003.} \textbf{Myth:} Kidney disease always causes kidney-area pain.\newline
\textbf{Status:} False.\newline
\textbf{Fact:} Chronic kidney disease is frequently painless. Early chronic kidney disease often has no symptoms because remaining kidney function compensates. It may be detected through blood pressure, creatinine and estimated filtration, urine albumin, or other testing. Flank pain is more associated with selected problems such as stones, infection, obstruction, or bleeding and is not required for kidney disease. Diabetes and hypertension are major risk factors. Swelling, blood in urine, reduced urine, fever, or severe pain warrant assessment. Screening should be risk-based rather than waiting for pain.
\item[] \textbf{KID-004.} \textbf{Myth:} Kidney stones always require surgery.\newline
\textbf{Status:} False.\newline
\textbf{Fact:} Many small stones pass spontaneously. Management depends on stone size, location, composition, obstruction, infection, pain control, kidney function, and the likelihood of passage. Small uncomplicated stones may pass with observation and medical therapy. Larger, persistent, obstructing, infected, or highly painful stones may require a procedure. Fever with suspected obstruction is an emergency because infection can spread rapidly. Imaging and urine testing help guide decisions. Follow-up is important because passing one stone does not eliminate recurrence risk.
\item[] \textbf{KID-005.} \textbf{Myth:} All kidney stones are caused by calcium-rich foods.\newline
\textbf{Status:} False.\newline
\textbf{Fact:} Stone formation has multiple causes; restricting dietary calcium can sometimes be counterproductive. Calcium oxalate stones are common, but stones can also contain uric acid, struvite, cystine, or other compounds. Genetics, dehydration, sodium, animal protein, bowel disease, medicines, infection, and urine chemistry influence risk. Normal dietary calcium can bind oxalate in the gut and may reduce absorption into urine. Unsupervised severe calcium restriction can therefore worsen some stone patterns. A stone analysis and metabolic evaluation can identify prevention strategies. Adequate fluid and lower sodium are often useful, but advice must be individualized.
\item[] \textbf{KID-006.} \textbf{Myth:} Normal urination proves the kidneys are normal.\newline
\textbf{Status:} False.\newline
\textbf{Fact:} Significant kidney dysfunction can occur despite apparently normal urine output. Kidney function includes filtration, electrolyte regulation, acid-base balance, hormone production, and removal of waste, not merely making urine. A person can produce a normal or large volume while filtration is reduced or urine protein is leaking. Early disease is often silent and requires blood and urine testing for detection. Sudden reduction in urine can indicate acute injury, but normal output does not exclude it. Risk assessment is important in diabetes, hypertension, cardiovascular disease, family history, and use of kidney-stressing medicines. Results should be interpreted over time by a clinician.
\item[] \textbf{KID-007.} \textbf{Myth:} You should not swim for 30 minutes after eating.\newline
\textbf{Status:} False.\newline
\textbf{Fact:} Eating does not create a dangerous swimming-related cramp risk in healthy people. There is no established universal waiting period after a meal before swimming. Some people may feel discomfort, reflux, or sluggishness after a large meal and can choose to rest for comfort. Drowning prevention depends more on supervision, swimming ability, life jackets, alcohol avoidance, and safe water conditions. Children require close, active supervision near water regardless of meal timing. Sudden intense exercise after a heavy meal may be unpleasant without being a predictable drowning mechanism. Follow local safety rules and personal tolerance.
\item[] \textbf{KID-008.} \textbf{Myth:} Swallowed chewing gum remains in your stomach for seven years.\newline
\textbf{Status:} False.\newline
\textbf{Fact:} It usually passes through the gastrointestinal tract. The gum base is not fully digested, but the digestive tract normally moves it onward with other intestinal contents. An occasional piece usually leaves in stool within days. Swallowing large amounts or gum together with constipation can rarely contribute to obstruction, particularly in children. Gum should not be swallowed routinely, but the seven-year claim is unsupported. Persistent abdominal pain, vomiting, bloating, or inability to pass stool needs medical attention. Keep gum and other small objects away from young children.
\item[] \textbf{KID-009.} \textbf{Myth:} Holding in a sneeze is always harmless.\newline
\textbf{Status:} False.\newline
\textbf{Fact:} Forcefully suppressing a sneeze can rarely cause injury. A sneeze rapidly expels air and secretions to clear irritants from the nose and airway. Pinching the nose and closing the mouth can sharply increase pressure in the upper airway and ears. Serious injuries are rare but reported, including eardrum or throat injury and complications in vulnerable tissues. It is safer to sneeze into a tissue or the elbow while keeping distance from others. Hand hygiene reduces spread of respiratory pathogens. Recurrent sneezing with wheeze, facial pain, fever, or blood should be assessed.
\item[] \textbf{KID-010.} \textbf{Myth:} Sneezing means you are getting a cold.\newline
\textbf{Status:} False.\newline
\textbf{Fact:} Allergies, irritants, infections, and other causes can trigger sneezing. Sneezing is a protective reflex and is not specific to a particular diagnosis. Allergic rhinitis often causes repeated sneezing with itching and watery discharge, while irritants can cause abrupt symptoms without infection. Viral infections may add sore throat, cough, fever, or fatigue, but symptoms overlap. Treatment depends on the cause and can include trigger avoidance, saline, antihistamines, or medical evaluation. Persistent unilateral symptoms, recurrent nosebleeds, or breathing difficulty deserve review. A sneeze alone cannot determine whether antibiotics or other treatment is needed.
\item[] \textbf{KID-011.} \textbf{Myth:} Your heart stops when you sneeze.\newline
\textbf{Status:} False.\newline
\textbf{Fact:} It does not. A sneeze can briefly change pressure in the chest and alter blood flow returning to the heart. The heartbeat may change timing or rhythm for a moment, which can feel like a pause. It does not normally stop the heart. The heart's electrical system continues to coordinate contraction through the sneeze. Fainting, chest pain, sustained palpitations, or breathlessness should not be blamed on a simple reflex. Such symptoms require medical assessment.
\item[] \textbf{KID-012.} \textbf{Myth:} Humans have only five senses.\newline
\textbf{Status:} False.\newline
\textbf{Fact:} We also perceive balance, body position, temperature, pain, and other sensory modalities. The traditional five categories are sight, hearing, smell, taste, and touch, but physiology recognizes additional sensory systems. Vestibular organs detect head movement and balance, while proprioceptors signal limb position and movement. Nociception detects potentially damaging stimuli, and thermoreception detects temperature. Internal receptors also monitor stretch, blood chemistry, and organ state. The exact number depends on how sensory modalities are defined. The five-sense list is a teaching simplification, not a complete account of human perception.
\item[] \textbf{KID-013.} \textbf{Myth:} Fingernails continue growing after death.\newline
\textbf{Status:} False.\newline
\textbf{Fact:} Tissue dehydration makes surrounding skin retract, creating this appearance. Nail growth requires living cells in the nail matrix, circulation, and metabolism. After death, these processes stop. Dehydration causes skin and soft tissue to shrink or retract, making nails and hair appear more prominent. This visual change is not continued growth. Similar apparent changes can occur because of body position, lighting, or handling. Claims about postmortem growth should therefore be interpreted as changes in surrounding tissue, not new keratin production.
\item[] \textbf{KID-014.} \textbf{Myth:} You lose most body heat through your head.\newline
\textbf{Status:} False.\newline
\textbf{Fact:} Heat loss depends mainly on how much skin is exposed. The head can lose substantial heat when it is uncovered, because it is vascular and often exposed. It is not a special heat-loss outlet that accounts for a fixed majority of total loss. Any uncovered skin contributes according to area, temperature gradient, wind, moisture, and clothing. A hat helps in cold conditions by reducing exposed surface area. In infants, older adults, and people with impaired thermoregulation, all clothing and environment factors matter. Preventing hypothermia requires protecting the whole body, not only the head.
\item[] \textbf{KID-015.} \textbf{Myth:} Sitting on a cold floor causes piles or hemorrhoids.\newline
\textbf{Status:} False.\newline
\textbf{Fact:} Cold surfaces are not an established cause. Hemorrhoids are swollen vascular cushions associated with pressure, straining, constipation, pregnancy, prolonged sitting on the toilet, and other factors. A cold floor may be uncomfortable but does not create the underlying venous changes. Persistent rectal bleeding should not be assumed to be hemorrhoids, especially in older adults or when bowel habits change. Fiber, fluids, avoiding straining, and shorter toilet time may help selected symptoms. Severe pain, heavy bleeding, fever, or a mass needs assessment. Diagnosis is important because other conditions can mimic hemorrhoids.
\item[] \textbf{KID-016.} \textbf{Myth:} Sitting on a toilet seat commonly transmits serious infections.\newline
\textbf{Status:} False.\newline
\textbf{Fact:} This is an uncommon transmission route for most infections people worry about. Many pathogens require direct contact, respiratory droplets, sexual contact, blood exposure, food, or a suitable surface-to-mouth route. Intact skin on a toilet seat is an effective barrier for most organisms. Rare exceptions or contaminated bodily fluids do not justify treating every public seat as a major infection source. Hand hygiene remains important because bathroom surfaces can transfer organisms to hands. Open wounds, unusual sores, or exposure to blood need specific advice. Fear of toilet seats should not replace vaccination, safer sex, food safety, or handwashing.
\item[] \textbf{KID-017.} \textbf{Myth:} Drinking water during meals ruins digestion.\newline
\textbf{Status:} False.\newline
\textbf{Fact:} Normal amounts of water do not impair digestion. Water mixes with food and digestive secretions and is absorbed as part of normal gastrointestinal physiology. It does not dilute stomach acid enough to stop digestion in a healthy person. Drinking with meals may help swallowing and contribute to hydration. Some people with reflux, early satiety, or gastroparesis may prefer smaller amounts for comfort. Large volumes can cause temporary fullness but not a general digestive failure. Persistent pain, vomiting, swallowing difficulty, or unintended weight loss needs evaluation.
\item[] \textbf{KID-018.} \textbf{Myth:} Drinking water immediately after eating causes obesity.\newline
\textbf{Status:} False.\newline
\textbf{Fact:} Water contains no calories. Body fat increases when energy intake exceeds expenditure over time, not because of the timing of calorie-free water. Water can increase fullness and sometimes support lower calorie intake when substituted for sugary beverages. The amount of fluid needed varies with climate, activity, food, illness, and medical conditions. Forced overhydration is unsafe, but ordinary drinking with or after meals is not a cause of obesity. Sweetened drinks are a different issue because they add energy. Weight concerns should be addressed through the overall dietary pattern and health context.
\item[] \textbf{KID-019.} \textbf{Myth:} Fruit must never be eaten after a meal.\newline
\textbf{Status:} False.\newline
\textbf{Fact:} There is no universal medical requirement for this. The digestive tract processes mixed meals through coordinated stomach and intestinal functions rather than a rigid food-order system. Fruit can be eaten before, during, or after meals according to preference and tolerance. Some people with reflux, bloating, diabetes, or bowel disorders may find particular portions or timing more comfortable. Whole fruit provides fiber and micronutrients regardless of clock position. Avoiding fruit unnecessarily can reduce diet quality. Persistent symptoms after eating should be evaluated instead of attributed to a universal fruit rule.
\item[] \textbf{KID-020.} \textbf{Myth:} Fruit eaten with other food ferments dangerously in your stomach.\newline
\textbf{Status:} False.\newline
\textbf{Fact:} Normal gastrointestinal physiology does not work this way. The stomach's acid and enzymes begin digestion, while the small intestine absorbs nutrients and the colon contains most normal microbial fermentation. Food does not become dangerously fermented merely because fruit is eaten with other foods. Gas or bloating can result from fermentable carbohydrates, intolerance, altered motility, or underlying bowel disorders. Individual triggers are real but do not support a universal food-combining law. Balanced mixed meals are compatible with normal digestion. Severe pain, vomiting, bleeding, or weight loss needs medical review.
\item[] \textbf{KID-021.} \textbf{Myth:} Milk and fish eaten together cause vitiligo.\newline
\textbf{Status:} False.\newline
\textbf{Fact:} There is no evidence that this food combination causes vitiligo. Vitiligo is an autoimmune disorder in which melanocytes are damaged or destroyed. It is not caused by a toxic reaction between fish and dairy products. Genetics, immune regulation, oxidative stress, and other factors may influence susceptibility. Nutritional deficiencies can affect general health but do not make this food pairing a recognized cause. Avoiding nutritious foods on this basis can create unnecessary restriction. New depigmented patches should be assessed by a clinician rather than managed through food myths.
\item[] \textbf{KID-022.} \textbf{Myth:} Eating curd or yogurt at night is inherently harmful.\newline
\textbf{Status:} False.\newline
\textbf{Fact:} There is no general medical prohibition against it. Yogurt can be a source of protein, calcium, and other nutrients at any time of day. Some people experience reflux, bloating, lactose intolerance, or allergy after dairy, but these are individual conditions rather than a universal night-time effect. Fermented yogurt is not automatically mucus-forming or infection-causing. Portion, added sugar, and overall diet quality matter more than the clock. People with milk allergy should avoid the relevant proteins and seek dietary advice. Persistent symptoms need diagnosis instead of a blanket timing rule.
\item[] \textbf{KID-023.} \textbf{Myth:} Bananas cause colds.\newline
\textbf{Status:} False.\newline
\textbf{Fact:} Respiratory viruses cause colds. The common cold is caused by respiratory viruses transmitted through droplets, aerosols, hands, and contaminated surfaces. Eating a banana does not create a viral infection. A cold food may briefly feel uncomfortable for someone with a sore throat, but sensation is not causation. Bananas provide carbohydrate, fiber, potassium, and other nutrients for many people. Allergy or intolerance is possible in specific individuals and is a separate issue. Fever, breathing difficulty, or persistent symptoms should be evaluated on clinical grounds.
\item[] \textbf{KID-024.} \textbf{Myth:} Citrus fruits cause respiratory infections.\newline
\textbf{Status:} False.\newline
\textbf{Fact:} They do not cause viral infections. Citrus fruits cannot create the viruses that cause colds or influenza. They provide vitamin C and other nutrients, although eating citrus does not guarantee prevention of infection. Acidic foods may irritate a sore mouth or reflux in some people without causing the illness itself. Infection risk is influenced by exposure, vaccination where available, ventilation, hand hygiene, and immune factors. Citrus allergy is uncommon but possible. Persistent respiratory symptoms require assessment rather than avoidance of an entire fruit group.
\item[] \textbf{KID-025.} \textbf{Myth:} Spicy food causes appendicitis.\newline
\textbf{Status:} False.\newline
\textbf{Fact:} It is not an established cause. Appendicitis usually results from obstruction and inflammation of the appendix, with contributions from lymphoid swelling, stool, or other material. Spices are not a recognized routine cause. Spicy food can cause burning, reflux, diarrhea, or abdominal discomfort in some people, which may confuse symptom interpretation. Progressive pain that moves to the right lower abdomen, fever, vomiting, or pain with movement needs urgent medical assessment. Do not use food avoidance or laxatives to treat suspected appendicitis. Delayed treatment can allow perforation and serious infection.
\item[] \textbf{KID-026.} \textbf{Myth:} Swallowing fruit seeds commonly causes appendicitis.\newline
\textbf{Status:} False.\newline
\textbf{Fact:} This is extraordinarily uncommon and is not a routine cause. Most small seeds pass through the gastrointestinal tract without causing obstruction. Appendicitis is usually related to inflammation and blockage from a range of causes, not ordinary fruit consumption. Rare retained foreign bodies have been reported, but case reports do not establish a common risk. Avoiding all seeded fruit is unnecessary for most people. Severe or progressive abdominal pain, fever, vomiting, or localized tenderness requires urgent evaluation regardless of what was eaten. Never delay care while trying to identify a single food culprit.
\item[] \textbf{KID-027.} \textbf{Myth:} Obesity simply means someone lacks self-control.\newline
\textbf{Status:} False.\newline
\textbf{Fact:} Body weight is influenced by genetics, physiology, environment, medications, behavior, and socioeconomic factors. Appetite regulation, metabolism, sleep, stress, endocrine disease, medicines, food availability, and social conditions all influence weight. Personal behavior matters, but it operates within biological and environmental constraints. Stigma can discourage people from seeking care and is not an effective treatment. Health assessment includes blood pressure, glucose, lipids, fitness, sleep, mental health, and function rather than a moral judgment. Evidence-based care may include nutrition, activity, behavioral support, medicines, or surgery. Goals should be individualized and respectful.
\item[] \textbf{KID-028.} \textbf{Myth:} Thin automatically means healthy.\newline
\textbf{Status:} False.\newline
\textbf{Fact:} Body size alone does not establish metabolic or cardiovascular health. A person with a low or average body mass can have hypertension, diabetes, high cholesterol, malnutrition, anemia, or significant disease. Body-composition distribution and fitness are not fully captured by appearance. Some illnesses cause unintentional weight loss and may be missed when thinness is treated as inherently healthy. Preventive care, nutrition, sleep, activity, and mental health matter across body sizes. Weight loss without trying needs assessment. Clinicians should avoid assumptions based solely on body shape.
\item[] \textbf{KID-029.} \textbf{Myth:} Being overweight automatically means someone is unhealthy.\newline
\textbf{Status:} False.\newline
\textbf{Fact:} Individual health requires assessment beyond body weight alone. Higher weight can increase risk for some conditions, but risk varies with fat distribution, fitness, genetics, metabolic markers, sleep, and access to care. A weight category is not a complete diagnosis or a measure of worth. People can have healthy behaviors and normal tests at different body sizes, while disease can occur at any size. Weight stigma can worsen stress, avoidance of care, and disordered eating. Clinicians should assess health directly and discuss goals collaboratively. Treatment should focus on meaningful outcomes rather than shame.
\item[] \textbf{KID-030.} \textbf{Myth:} You can lose substantial body fat from one particular area by exercising it.\newline
\textbf{Status:} False.\newline
\textbf{Fact:} Spot reduction is not a reliable mechanism of fat loss. Local exercises strengthen the muscles beneath a body region but do not control where stored fat is mobilized. Fat distribution is influenced by genetics, sex hormones, age, energy balance, and overall physiology. Resistance training can change shape through muscle growth even when local fat loss does not occur. Sustainable weight management uses a combination of nutrition, activity, sleep, and behavior support when desired. Cosmetic procedures are distinct from exercise and carry their own risks. Marketing claims promising targeted fat melting should be treated skeptically. \emph{(See also GEN-087.)}
\item[] \textbf{KID-031.} \textbf{Myth:} Starvation diets are the fastest healthy way to lose weight.\newline
\textbf{Status:} False.\newline
\textbf{Fact:} Severe restriction can cause nutritional deficiencies and loss of lean tissue. Extreme restriction may produce rapid scale changes, but much of the early loss can be water and glycogen rather than fat. It can reduce protein, micronutrient, and fiber intake and promote loss of muscle. Hunger, fatigue, gallstones, electrolyte abnormalities, binge-restrict cycles, and medication complications are possible. Safer weight management uses an eating pattern that can be maintained and preserves nutrition and function. Some medically supervised very-low-calorie diets exist for selected patients, but they require monitoring. Eating-disorder history and medical conditions must be considered.
\item[] \textbf{KID-032.} \textbf{Myth:} Carbohydrates eaten at night automatically become fat.\newline
\textbf{Status:} False.\newline
\textbf{Fact:} Metabolism does not operate according to such a simple clock rule. Carbohydrate is metabolized according to total intake, activity, insulin sensitivity, glycogen stores, and the overall energy balance. Late meals can be associated with weight gain when they add extra calories, disturb sleep, or involve energy-dense foods. That association does not mean every nighttime carbohydrate is converted directly to fat. Meal timing may matter for some metabolic outcomes, but there is no universal ban. Portion, food quality, symptoms, work schedule, and treatment should guide timing. People with diabetes may need individualized advice about evening medicines and meals.
\item[] \textbf{KID-033.} \textbf{Myth:} Lemon water melts body fat.\newline
\textbf{Status:} False.\newline
\textbf{Fact:} It does not directly burn fat. Lemon water supplies fluid and flavor but has no special mechanism that dissolves stored adipose tissue. It may help someone replace a sugary drink, indirectly reducing calorie intake. The citric acid does not selectively accelerate fat breakdown in a clinically meaningful way. Excess acidity can worsen reflux or erode dental enamel if consumed frequently. Sustainable fat loss depends on overall energy balance, behavior, activity, sleep, and biology. No beverage can substitute for a safe, individualized plan.
\item[] \textbf{KID-034.} \textbf{Myth:} Hot water melts abdominal fat.\newline
\textbf{Status:} False.\newline
\textbf{Fact:} Water temperature does not selectively remove body fat. Body fat is stored in adipose cells and is mobilized through metabolic signaling, not local contact with hot liquid. Drinking hot water may feel soothing and contributes to hydration, but it does not melt tissue inside the abdomen. Very hot drinks can burn the mouth or esophagus. Weight change after sweating or reduced intake is not proof of targeted fat loss. Exercise and dietary patterns can alter total body composition over time. Persistent abdominal enlargement can have medical causes unrelated to fat and should be assessed when new or concerning.
\item[] \textbf{KID-035.} \textbf{Myth:} Apple-cider vinegar rapidly burns fat.\newline
\textbf{Status:} False.\newline
\textbf{Fact:} Evidence does not support dramatic fat-loss claims. Small studies of vinegar may show modest effects on appetite or metabolic measures, but they do not establish rapid, clinically meaningful fat burning. Vinegar can irritate the esophagus, worsen reflux, erode enamel, lower potassium, and interact with medicines. Concentrated products are especially hazardous. Replacing a high-calorie dressing or drink may alter total intake, but that is a behavioral effect rather than a fat-melting property. Proven weight-management care should not be replaced by vinegar. Dilution does not eliminate all risks.
\item[] \textbf{KID-036.} \textbf{Myth:} Sweating off 1 kg means you have lost 1 kg of fat.\newline
\textbf{Status:} False.\newline
\textbf{Fact:} Acute weight loss through sweating is predominantly water. One kilogram of body fat represents a large energy deficit that cannot ordinarily be created by a single hot workout. Sweating reduces body water and may temporarily lower scale weight. Rehydration restores much of that loss, which is appropriate and necessary. Deliberate dehydration can impair circulation, kidney function, thermoregulation, and athletic performance. Fat loss should be assessed over longer periods and alongside nutrition and activity patterns. Heat illness symptoms such as confusion, fainting, or cessation of sweating require emergency care. \emph{(See also GEN-088.)}
\item[] \textbf{KID-037.} \textbf{Myth:} People with diabetes must completely avoid fruit.\newline
\textbf{Status:} False.\newline
\textbf{Fact:} Appropriate portions of whole fruit can generally be incorporated into a diabetes eating plan. Whole fruit contains carbohydrate but also fiber, water, vitamins, and phytochemicals. Blood-glucose response varies by fruit, portion, ripeness, meal context, and individual treatment. Avoiding all fruit is usually unnecessary and can reduce diet variety. People using insulin or medicines that cause hypoglycemia may need carbohydrate planning and monitoring. Juice and dried fruit are more concentrated and easier to overconsume. A dietitian or diabetes clinician can help match portions to goals and glucose patterns.
\item[] \textbf{KID-038.} \textbf{Myth:} Jaggery is safe for diabetes because it is not refined sugar.\newline
\textbf{Status:} False.\newline
\textbf{Fact:} Jaggery contains substantial amounts of sugar and raises blood glucose. Jaggery is less processed than white sugar but remains rich in sucrose and other simple carbohydrates. Small amounts of minerals do not cancel its glycemic and calorie contribution. It should be counted as an added or concentrated sugar in a diabetes eating plan. Marketing terms such as natural, traditional, or unrefined do not establish safety. Portion size and the total meal determine the response, which can be monitored when appropriate. Diabetes medicines should not be changed to accommodate a sweetener without advice.
\item[] \textbf{KID-039.} \textbf{Myth:} Dates do not affect blood glucose because they are natural.\newline
\textbf{Status:} False.\newline
\textbf{Fact:} Dates contain carbohydrates and sugars. Dates are a nutritious fruit for many people, but natural sugar still contributes carbohydrate and energy. Fresh and dried forms differ in water content and therefore in carbohydrate concentration per bite or serving. Portion size, pairing with other foods, and individual glucose response matter. People with diabetes do not necessarily need to avoid dates, but they should account for them in their plan. Large portions can raise glucose substantially. The word natural does not mean carbohydrate-free or unlimited.
\item[] \textbf{KID-040.} \textbf{Myth:} Fruit juice is equivalent to whole fruit.\newline
\textbf{Status:} False.\newline
\textbf{Fact:} Whole fruit generally provides more intact fiber and produces different satiety and glycemic effects. Juicing removes or reduces much of the fiber structure and makes it easier to consume multiple fruits quickly. Juice can therefore deliver substantial sugar and calories with less chewing and fullness. Some juices retain micronutrients, but that does not make them nutritionally identical to whole fruit. Portion size and added sugar are important, especially for children and people with diabetes. Whole fruit is usually the more filling choice. Juice may fit in a diet, but it should not be treated as unlimited health water.
\item[] \textbf{KID-041.} \textbf{Myth:} Bitter gourd can cure diabetes.\newline
\textbf{Status:} False.\newline
\textbf{Fact:} It has not been demonstrated to cure diabetes. Bitter gourd contains compounds that may affect glucose in laboratory or small human studies. That evidence does not establish a cure, reliable dosing, or protection from complications. Diabetes is a chronic metabolic condition requiring monitoring and individualized treatment. Replacing prescribed therapy with bittergourd products can cause severe hyperglycemia or delay emergency care. Supplements may also interact with medicines or lower glucose unexpectedly. Foods can be part of a healthy diet, but cure claims require robust clinical evidence.
\item[] \textbf{KID-042.} \textbf{Myth:} Fenugreek can replace diabetes medication.\newline
\textbf{Status:} False.\newline
\textbf{Fact:} Evidence does not support using it as a general replacement for proven therapy. Fenugreek may modestly influence glucose in some studies, but products, doses, and study quality vary. It has not been shown to replace insulin or established diabetes medicines across patients. It can cause gastrointestinal symptoms, allergy, hypoglycemia, and drug interactions. People who stop effective treatment based on supplement claims risk acute and chronic complications. Any supplement should be disclosed to the diabetes team. Monitoring is especially important when combining it with glucose-lowering medicines.
\item[] \textbf{KID-043.} \textbf{Myth:} Once glucose becomes normal, diabetes has necessarily disappeared.\newline
\textbf{Status:} False.\newline
\textbf{Fact:} Normal values may represent successful treatment or remission rather than permanent cure. Normal glucose can result from medication, nutrition, activity, weight change, recovery from illness, or temporary physiology. Type 2 diabetes may enter remission in some circumstances, but recurrence remains possible and requires follow-up. Type 1 diabetes requires ongoing insulin even if readings temporarily improve. Stopping medication without advice can lead to dangerous hyperglycemia. A1C, home readings, symptoms, and the underlying diagnosis are interpreted together. ``Normal today'' is not proof that risk has permanently vanished.
\item[] \textbf{KID-044.} \textbf{Myth:} You can reliably judge your glucose from how you feel.\newline
\textbf{Status:} False.\newline
\textbf{Fact:} Hyperglycemia and hypoglycemia can sometimes occur with few or misleading symptoms. Symptoms such as thirst, fatigue, sweating, tremor, confusion, or headache are not specific and may be absent. People can adapt to chronic hyperglycemia or lose warning symptoms of hypoglycemia. Meter or continuous-glucose monitoring provides information that sensation cannot. Treatment plans should specify when to test and how to respond. Confusion, seizure, unconsciousness, vomiting, or severe breathing changes is an emergency. Do not drive or exercise intensely when serious glucose abnormality is suspected.
\item[] \textbf{KID-045.} \textbf{Myth:} Only people taking insulin develop hypoglycemia.\newline
\textbf{Status:} False.\newline
\textbf{Fact:} Several non-insulin glucose-lowering drugs can also cause it. Sulfonylureas and some related medicines can stimulate insulin release and produce hypoglycemia. Risk is also influenced by missed meals, alcohol, kidney disease, exercise, and drug interactions. Other therapies have lower risk when used alone but can contribute in combinations. Symptoms may include sweating, shaking, hunger, palpitations, confusion, or unusual behavior. Patients should know their treatment-specific action plan and carry an appropriate rapid carbohydrate source when advised. Recurrent episodes require medication and cause review.
\item[] \textbf{KID-046.} \textbf{Myth:} Diabetes affects only blood sugar.\newline
\textbf{Status:} False.\newline
\textbf{Fact:} Chronic diabetes can affect blood vessels, nerves, kidneys, eyes, heart, and other organs. Persistently abnormal glucose interacts with blood pressure, lipids, inflammation, and vascular biology. Complications can affect the retina, kidneys, nerves, feet, heart, brain, and immune response. Risk varies by duration, control, genetics, smoking, blood pressure, and other factors. Screening can detect kidney, eye, foot, and cardiovascular problems before severe symptoms occur. Good management reduces risk but does not make complications impossible. Diabetes care therefore includes more than checking a single glucose number.
\item[] \textbf{KID-047.} \textbf{Myth:} Headache is a reliable way to know your blood pressure is high.\newline
\textbf{Status:} False.\newline
\textbf{Fact:} Most hypertension causes no symptoms. Blood pressure can be markedly elevated without headache, dizziness, or any outward sign. Headache has many causes and does not diagnose hypertension. Very severe elevation with neurologic, visual, chest, or breathing symptoms can represent an emergency. Accurate measurement with a validated cuff is more useful than symptom guessing. Repeated high readings need clinical assessment even when the person feels well. Do not stop or double medicine based only on a headache.
\item[] \textbf{KID-048.} \textbf{Myth:} Feeling dizzy means your blood pressure must be low.\newline
\textbf{Status:} False.\newline
\textbf{Fact:} Dizziness has numerous possible causes. Dizziness can result from inner-ear disorders, dehydration, anemia, arrhythmia, medications, low glucose, anxiety, neurologic disease, or blood-pressure changes. Some people with high blood pressure feel normal, while others with normal readings feel dizzy. Sitting or lying down safely and checking relevant measurements may help immediate safety. Fainting, chest pain, severe headache, weakness, trouble speaking, or persistent symptoms needs urgent evaluation. Medication review is important when symptoms begin after a change. A blood-pressure reading alone does not explain every sensation.
\item[] \textbf{KID-049.} \textbf{Myth:} One high blood-pressure measurement means you definitely have chronic hypertension.\newline
\textbf{Status:} False.\newline
\textbf{Fact:} Diagnosis usually requires properly obtained repeated or out-of-office measurements, except in particular clinical situations. Blood pressure changes with stress, pain, caffeine, exercise, talking, posture, and cuff size. A single reading can be an important signal but usually needs confirmation under standardized conditions. Home or ambulatory monitoring can identify sustained hypertension and white-coat or masked patterns. Extremely high readings with symptoms require urgent care rather than waiting for routine confirmation. Repeated measurements should be taken after rest with the arm supported and the correct cuff. Clinicians interpret the pattern alongside overall cardiovascular risk.
\item[] \textbf{KID-050.} \textbf{Myth:} Blood pressure should always be exactly 120/80 mmHg.\newline
\textbf{Status:} False.\newline
\textbf{Fact:} Blood pressure naturally varies, and treatment targets depend on clinical circumstances. Blood pressure changes across the day and responds to posture, activity, sleep, emotion, illness, and medicines. Risk categories use ranges, not a single universally perfect pair of numbers. Treatment goals depend on age, pregnancy, kidney disease, diabetes, cardiovascular risk, side effects, and guideline context. A lower reading is not automatically better if it causes fainting or organ under-perfusion. Measurements need correct technique and repeated interpretation. Individual targets should be agreed with the treating clinician.
\item[] \textbf{KID-051.} \textbf{Myth:} If blood pressure becomes normal, antihypertensive medication should automatically be stopped.\newline
\textbf{Status:} False.\newline
\textbf{Fact:} The medication may be responsible for the normal reading. A normal reading during treatment may show that the treatment is working, not that hypertension has permanently disappeared. Stopping medicines abruptly can allow blood pressure to rise and may be dangerous with some drug classes. In selected people, weight change, improved health, pregnancy, or other circumstances may justify supervised dose reduction. Decisions require repeated measurements and review of kidney function, symptoms, and cardiovascular risk. Do not double, skip, or stop medication based on one reading. Any change should be planned with the prescribing clinician. \emph{(See also GEN-035.)}
\item[] \textbf{KID-052.} \textbf{Myth:} People with hypertension must completely eliminate salt.\newline
\textbf{Status:} False.\newline
\textbf{Fact:} Sodium reduction is often beneficial; absolute elimination is not required or feasible. Sodium can increase blood pressure in salt-sensitive people, and reducing excess intake often helps. The body still requires sodium for fluid balance, nerves, and muscles. Most dietary sodium comes from packaged, restaurant, and processed foods rather than the saltshaker alone. A sustainable pattern emphasizes less highly processed food and label awareness rather than dangerous zero-sodium restriction. Individual goals differ with kidney disease, heart failure, sweating, and medicines. Dietary changes complement rather than automatically replace prescribed treatment.
\item[] \textbf{KID-053.} \textbf{Myth:} Heart disease always produces warning symptoms.\newline
\textbf{Status:} False.\newline
\textbf{Fact:} Significant cardiovascular disease can be clinically silent. Coronary disease, high blood pressure, rhythm disorders, and structural heart disease may be present without obvious symptoms. Some people develop atypical symptoms such as fatigue, breathlessness, nausea, jaw discomfort, or reduced exercise tolerance. Diabetes and older age can blunt warning signs. Risk assessment, preventive care, and appropriate testing are therefore important even when someone feels well. New chest pressure, severe breathlessness, fainting, or neurologic symptoms requires emergency care. Absence of symptoms is not proof of a healthy heart.
\item[] \textbf{KID-054.} \textbf{Myth:} A normal cholesterol result guarantees you will not have a heart attack.\newline
\textbf{Status:} False.\newline
\textbf{Fact:} Cardiovascular risk depends on many factors. Heart attacks can result from plaque, clotting, inflammation, spasm, or other mechanisms that are not captured by one lipid panel. Risk also depends on blood pressure, smoking, diabetes, kidney disease, age, family history, and prior disease. A ``normal'' value may be interpreted differently according to baseline risk and treatment. Cholesterol management is one part of prevention, not a guarantee. Healthy habits and prescribed medicines reduce risk without making it zero. Emergency symptoms require immediate care regardless of a past laboratory result.
\item[] \textbf{KID-055.} \textbf{Myth:} Normal blood pressure guarantees your heart is healthy.\newline
\textbf{Status:} False.\newline
\textbf{Fact:} It excludes neither coronary disease nor many other cardiac disorders. Blood pressure is an important risk factor but does not measure coronary arteries, heart valves, rhythm, pumping function, or congenital structure. A person can have normal pressure and still have diabetes, high cholesterol, smoking-related risk, or an arrhythmia. Conversely, one elevated reading does not define the whole cardiovascular picture. Symptoms, examination, family history, and targeted tests guide evaluation. Exercise tolerance and preventive care also matter. Chest pain, palpitations, fainting, or unexplained breathlessness should be assessed.
\item[] \textbf{KID-056.} \textbf{Myth:} Glaucoma always causes eye pain.\newline
\textbf{Status:} False.\newline
\textbf{Fact:} The common open-angle form is typically painless. Open-angle glaucoma can progressively damage the optic nerve without pain or early awareness. Vision loss often begins in the peripheral field and may be compensated for by the other eye. Acute angle-closure glaucoma is different and can cause severe eye pain, halos, headache, nausea, and rapid visual loss. Both symptom patterns require attention, but lack of pain cannot exclude glaucoma. Eye-pressure measurement alone is insufficient; optic-nerve examination and visual fields are important. Risk-based eye examinations help detect disease before irreversible loss.
\item[] \textbf{KID-057.} \textbf{Myth:} Good vision means you cannot have glaucoma.\newline
\textbf{Status:} False.\newline
\textbf{Fact:} Peripheral visual-field damage can develop before central vision is affected. Early glaucoma may leave reading and central vision apparently normal while peripheral nerve fibers are lost. The brain can adapt to gradual change, and one eye may compensate for the other. By the time a person notices major visual loss, damage may be advanced and permanent. Screening intervals depend on age, family history, eye pressure, ethnicity, and other risks. Examination includes the optic nerve and often visual-field or imaging tests. Any sudden visual change or painful red eye needs urgent assessment.
\item[] \textbf{KID-058.} \textbf{Myth:} Glaucoma and cataract are the same disease.\newline
\textbf{Status:} False.\newline
\textbf{Fact:} They involve different structures and disease processes. A cataract is clouding of the eye's lens and commonly causes gradual blur, glare, and reduced contrast. Glaucoma is damage to the optic nerve, often associated with pressure-related risk, and produces characteristic field loss. A person can have either disease or both. Cataract surgery removes the cloudy lens but does not reverse established optic-nerve damage from glaucoma. Glaucoma treatment aims to lower risk of progression through medicines, laser, or surgery. Diagnosis determines the appropriate treatment.
\item[] \textbf{KID-059.} \textbf{Myth:} Cataract must become completely ``ripe'' before surgery.\newline
\textbf{Status:} False.\newline
\textbf{Fact:} Modern cataract surgery does not generally require waiting for a mature cataract. Surgery is usually considered when cataract-related vision loss interferes with daily activities, safety, work, or quality of life. Waiting until the lens is extremely advanced can increase visual disability and may make surgery more technically difficult in some cases. The decision depends on symptoms, examination, ocular health, and patient preference, not an old maturity rule. Glasses or lighting may help early symptoms but cannot remove the cataract. Other eye diseases can limit the expected benefit and should be discussed. A specialist can explain timing and risks.
\item[] \textbf{KID-060.} \textbf{Myth:} Cataract surgery permanently restores youthful vision in every patient.\newline
\textbf{Status:} False.\newline
\textbf{Fact:} Outcomes also depend on the retina, optic nerve, cornea, and other ocular conditions. Cataract surgery removes the cloudy lens and often improves clarity, glare, and contrast. It cannot repair retinal disease, optic-nerve damage, corneal scarring, or amblyopia. Many people still need glasses for reading, distance refinement, or astigmatism after surgery. Presbyopia continues because the natural focusing system has been replaced by an artificial lens with selected optical properties. Results are usually good but not identical for every eye. Expectations should be based on examination and discussion with the surgeon.
\item[] \textbf{KID-061.} \textbf{Myth:} Glaucoma damage can routinely be reversed with eye drops.\newline
\textbf{Status:} False.\newline
\textbf{Fact:} Treatment primarily aims to prevent or slow additional damage. Glaucoma damages retinal ganglion cells and the optic nerve, and established loss is generally permanent. Pressure-lowering drops can reduce the risk of further damage but do not restore lost nerve fibers. Some visual complaints may improve when another reversible factor is treated, but that is not reversal of glaucoma. Adherence and correct technique are important because treatment may have no immediate subjective effect. Laser or surgery may be needed when drops are insufficient or unsuitable. Regular visual-field and optic-nerve monitoring shows whether control is adequate.
\item[] \textbf{KID-062.} \textbf{Myth:} Eye pressure alone determines whether someone has glaucoma.\newline
\textbf{Status:} False.\newline
\textbf{Fact:} Optic-nerve structure, visual fields, and other factors are critical. Elevated intraocular pressure is a major risk factor, but some people with high pressure never develop optic-nerve damage. Conversely, normal-pressure glaucoma can occur with pressures in a statistically typical range. Diagnosis considers the nerve, visual field, corneal thickness, angle anatomy, family history, age, and change over time. A pressure reading can vary with technique and time of day. Treatment decisions should not be based on a single number alone. Comprehensive eye examination is more informative.
\item[] \textbf{KID-063.} \textbf{Myth:} Normal eye pressure means glaucoma is impossible.\newline
\textbf{Status:} False.\newline
\textbf{Fact:} Normal-tension glaucoma exists. Glaucoma is defined by characteristic optic-nerve and visual-field damage, not by a pressure threshold alone. In normal-tension glaucoma, damage occurs despite measured pressures within the usual population range. Vascular factors, nighttime pressure, corneal properties, and individual susceptibility may contribute. Repeated examinations can reveal progression that a single pressure reading misses. People with a strong family history or suspicious nerve appearance may need closer monitoring. Normal pressure is reassuring but not an absolute exclusion.
\item[] \textbf{KID-064.} \textbf{Myth:} Carrots can correct refractive errors and eliminate glasses.\newline
\textbf{Status:} False.\newline
\textbf{Fact:} Adequate vitamin A supports ocular health but does not correct myopia or astigmatism. Refractive errors occur when the eye's shape or optical surfaces prevent light from focusing on the retina. Carrots provide beta-carotene, which the body can use to make vitamin A, but they do not reshape the eye. Vitamin A deficiency can cause serious ocular disease, especially night blindness, but extra intake beyond adequacy does not improve ordinary myopia. Glasses, contact lenses, and selected procedures correct focus. Excess vitamin A supplements can be toxic. A balanced diet supports eyes without replacing an eye examination.
\item[] \textbf{KID-065.} \textbf{Myth:} Wearing somebody else's glasses permanently damages your eyes.\newline
\textbf{Status:} False.\newline
\textbf{Fact:} They can cause blur or discomfort but generally do not permanently damage adult eyes. Glasses change the path of light before it enters the eye and do not alter the eye's structure. Another person's prescription may cause blur, headache, eyestrain, or difficulty judging distance. Children should receive their own accurate prescription because uncorrected vision can affect development and learning. Sharing glasses is not a substitute for an examination, especially when vision changes suddenly. Pain, double vision, flashes, or loss of vision needs prompt care. Temporary discomfort should resolve after removing the unsuitable glasses.
\item[] \textbf{KID-066.} \textbf{Myth:} Breast pain usually means breast cancer.\newline
\textbf{Status:} False.\newline
\textbf{Fact:} Breast pain alone is rarely caused by cancer. Breast pain is commonly related to hormonal cycling, cysts, inflammation, trauma, medication, chest-wall muscles, or an ill-fitting support garment. Cancer can occur without pain, and rarely may cause painful symptoms, so the symptom cannot diagnose or exclude it. Persistent focal pain, a new lump, skin or nipple change, discharge, or an axillary mass deserves assessment. Evaluation depends on age, examination, risk, and imaging indications. Wearing a well-fitted bra and tracking cycle-related symptoms may help some people. Reassurance should follow appropriate assessment rather than dismissing all pain.
\item[] \textbf{KID-067.} \textbf{Myth:} Breast cancer always presents as a lump.\newline
\textbf{Status:} False.\newline
\textbf{Fact:} It can produce skin, nipple, architectural, or imaging changes without a palpable mass. Breast cancer may cause a mass, but it can also present with skin dimpling, nipple inversion, bloody discharge, thickening, asymmetry, or an imaging abnormality. Some cancers are found through screening before a person can feel anything. Most breast changes are benign, but a new or persistent change should be evaluated. A normal self-examination does not replace recommended screening. Diagnostic imaging and biopsy, when indicated, establish the diagnosis. Do not wait for a lump if other concerning changes appear.
\item[] \textbf{KID-068.} \textbf{Myth:} A movable breast lump can never be malignant.\newline
\textbf{Status:} False.\newline
\textbf{Fact:} Mobility does not reliably exclude malignancy. Benign fibroadenomas are often smooth and mobile, but malignant lesions can also be mobile, particularly when small. Texture, mobility, pain, and age cannot reliably classify a lump without appropriate evaluation. New, enlarging, hard, fixed, or otherwise concerning lumps need clinical examination. Imaging choice depends on age, pregnancy, breast density, and the finding. Biopsy is used when imaging cannot establish benignity. Do not repeatedly manipulate a lump or delay care because it moves.
\item[] \textbf{KID-069.} \textbf{Myth:} Breast cancer occurs only after menopause.\newline
\textbf{Status:} False.\newline
\textbf{Fact:} Younger women can develop it. Risk generally increases with age, but breast cancer can occur before menopause. Younger patients may have denser breasts, different tumor biology, pregnancy-related changes, and distinct treatment or fertility concerns. Family history, inherited variants, prior radiation, and other factors can increase risk. A new lump or persistent change should be evaluated at any age rather than dismissed as hormonal. Screening recommendations vary by risk and country. Risk assessment helps determine whether earlier or additional surveillance is appropriate.
\item[] \textbf{KID-070.} \textbf{Myth:} Breastfeeding guarantees protection against breast cancer.\newline
\textbf{Status:} False.\newline
\textbf{Fact:} It can reduce risk but does not eliminate it. Breastfeeding is associated with a modest reduction in breast-cancer risk, with effects influenced by duration and reproductive history. It does not prevent all breast cancers or replace screening and evaluation of symptoms. Risk also depends on age, genetics, hormones, alcohol, body composition, and other factors. Breastfeeding decisions should be supported for their many possible benefits without presenting them as cancer insurance. A lump, skin change, discharge, or persistent focal symptom needs assessment during or after lactation. Imaging can be performed when clinically indicated.
\item[] \textbf{KID-071.} \textbf{Myth:} Breast injury causes breast cancer.\newline
\textbf{Status:} False.\newline
\textbf{Fact:} Trauma is not an established cause, although injury can draw attention to a pre-existing lump. A blow can cause bruising, pain, fat necrosis, or a hematoma. These changes may feel like a lump and can prompt discovery of a lesion that was already present. Trauma does not create the genetic and cellular process that defines breast cancer. A new lump that persists after the injury should still be examined because benign and malignant causes can coexist. Imaging may be needed to distinguish scar-like change from another lesion. Worsening swelling, redness, fever, or severe pain needs care.
\item[] \textbf{KID-072.} \textbf{Myth:} Deodorant or antiperspirant causes breast cancer.\newline
\textbf{Status:} False.\newline
\textbf{Fact:} Current evidence does not establish this causal relationship. Deodorants reduce odor, while antiperspirants reduce sweating through local ingredients. There is no established mechanism or consistent evidence that ordinary use causes breast cancer. Aluminum exposure from antiperspirants has not been shown to create a breast-cancer risk in humans. Irritant or allergic skin reactions can occur and are a separate problem. People with a new breast or armpit lump should seek evaluation rather than changing cosmetic products alone. Avoid applying irritating products to broken skin.
\item[] \textbf{KID-073.} \textbf{Myth:} Mammography detects every breast cancer.\newline
\textbf{Status:} False.\newline
\textbf{Fact:} No imaging test has 100 percent sensitivity. Mammography can detect many cancers before they are palpable, but sensitivity varies with breast density, tumor features, age, technique, and interval between screens. Some cancers are hidden or develop between scheduled examinations. Screening reduces mortality at a population level but does not guarantee an individual cancer-free result. New symptoms need diagnostic evaluation even after a recent normal screen. Ultrasound, MRI, and biopsy may be used in selected circumstances. Recommendations should reflect age, risk, density, and local guidance.
\item[] \textbf{KID-074.} \textbf{Myth:} A normal mammogram means a new breast symptom can always be ignored.\newline
\textbf{Status:} False.\newline
\textbf{Fact:} Persistent or concerning symptoms can require further assessment despite negative imaging. Screening mammography is designed for people without a specific symptom and may not explain a new palpable or skin finding. Diagnostic imaging can target the symptom and may use additional views or ultrasound. Clinical examination remains important because imaging and examination each have limitations. Most symptoms are benign, but persistence, growth, bloody discharge, retraction, or skin changes deserve review. Do not repeatedly self-check to the point of injury or delay care for reassurance from an old report. Follow-up should be based on the current finding.
\item[] \textbf{KID-075.} \textbf{Myth:} BI-RADS 3 means ``early cancer.''\newline
\textbf{Status:} False.\newline
\textbf{Fact:} It means probably benign, with an expected malignancy probability of no more than 2 percent. BI-RADS is a standardized radiology assessment and recommendation system, not a cancer stage. Category 3 usually indicates a finding that is very likely benign but merits short-interval follow-up to document stability. It is not the same as BI-RADS 4 or 5, which generally prompt tissue diagnosis. Follow-up matters because a small residual risk remains. Patients should ask what interval and modality are recommended. The category applies to the imaging finding, not to the person's overall cancer identity.
\item[] \textbf{KID-076.} \textbf{Myth:} Masturbation causes blindness.\newline
\textbf{Status:} False.\newline
\textbf{Fact:} It does not. Masturbation is a common sexual behavior and has no established causal link to blindness or loss of visual acuity. Myths arose from older moral and medical beliefs rather than controlled evidence. Temporary fatigue, muscle discomfort, or emotional distress can occur depending on circumstances, but these are not blindness mechanisms. Eye disease has causes such as refractive error, infection, glaucoma, retinal disease, trauma, or vascular problems. Sudden visual loss is an emergency regardless of recent sexual activity. Sexual-health questions can be discussed confidentially without shame.
\item[] \textbf{KID-077.} \textbf{Myth:} Masturbation causes permanent physical weakness.\newline
\textbf{Status:} False.\newline
\textbf{Fact:} It does not. Normal sexual activity does not deplete a meaningful store of blood, protein, or life force. Temporary relaxation or tiredness after orgasm is a normal physiologic response, not permanent weakness. Persistent fatigue may reflect sleep problems, anemia, endocrine disease, medication, depression, infection, or other conditions. Compulsive behavior that causes injury or interferes with life is a behavioral concern, not evidence of semen-related depletion. Pain, bleeding, or skin injury warrants medical advice. Education should distinguish normal sexuality from symptoms needing care.
\item[] \textbf{KID-078.} \textbf{Myth:} Semen loss inherently causes severe nutritional depletion.\newline
\textbf{Status:} False.\newline
\textbf{Fact:} Normal ejaculation does not cause clinically meaningful nutritional depletion. Semen volume is small and its nutrient content is not sufficient to cause a major loss of protein, vitamins, or minerals in a healthy person. The body continuously produces reproductive fluids and replaces them through ordinary physiology. Weakness after ejaculation is more likely related to sleep, anxiety, illness, exertion, or expectation than to nutrient loss. Restrictive diets or supplements marketed to replace semen nutrients are unnecessary. Persistent symptoms should be assessed on their own merits. Shame-based explanations can delay appropriate care.
\item[] \textbf{KID-079.} \textbf{Myth:} Vasectomy causes impotence.\newline
\textbf{Status:} False.\newline
\textbf{Fact:} Vasectomy does not ordinarily impair erections or testosterone production. Vasectomy blocks sperm transport from the testes but does not remove the testes or stop testosterone production. Libido, erections, orgasm, and ejaculation usually remain intact, although semen no longer contains sperm. Short-term pain, bruising, swelling, or anxiety can temporarily affect sexual activity. Chronic post-vasectomy pain is uncommon but recognized and requires evaluation. Vasectomy is intended to be permanent contraception, not protection from STIs. Semen analysis is required after the procedure because fertility does not end immediately.
\item[] \textbf{KID-080.} \textbf{Myth:} Vasectomy immediately makes a man sterile.\newline
\textbf{Status:} False.\newline
\textbf{Fact:} Residual sperm remain for a period; semen testing is used to confirm success. Vasectomy interrupts sperm transport but does not instantly clear sperm already beyond the treated area. Another contraceptive method is needed until postoperative semen testing confirms success according to local guidance. Timing varies, so a calendar or a single ejaculation cannot prove sterility. The procedure does not protect against sexually transmitted infections. Persistent pain, swelling, fever, or wound problems should be reported. Follow-up testing is part of the procedure, not an optional formality.
\item[] \textbf{KID-081.} \textbf{Myth:} Contraceptive pills protect against sexually transmitted infections.\newline
\textbf{Status:} False.\newline
\textbf{Fact:} They do not. Hormonal pills prevent pregnancy by changing ovulation and reproductive physiology; they do not block infectious secretions or skin-to-skin transmission. Condoms and internal condoms reduce risk of many STIs when used correctly, though they do not cover every exposed area. Vaccination, testing, mutual agreements, and treatment of partners also matter. Many infections are asymptomatic, so absence of symptoms is not reassurance. Contraception and STI prevention are related but separate goals. Sexual-health plans should address both explicitly.
\item[] \textbf{KID-082.} \textbf{Myth:} You can always recognize an STI by symptoms.\newline
\textbf{Status:} False.\newline
\textbf{Fact:} Many STIs can be asymptomatic. Chlamydia, gonorrhea, HPV, herpes, HIV, syphilis, and other infections may cause no symptoms or only mild, nonspecific changes. Symptoms can include discharge, sores, pelvic pain, rash, burning, bleeding, or swollen glands, but they do not identify the organism reliably. Testing is based on exposure, anatomy, symptoms, pregnancy, and local recommendations. Untreated infection can cause complications and transmission even when a person feels well. Barrier protection reduces risk but does not eliminate it. A clinician can discuss confidential testing and treatment.
\item[] \textbf{KID-083.} \textbf{Myth:} Washing immediately after intercourse prevents pregnancy.\newline
\textbf{Status:} False.\newline
\textbf{Fact:} It does not. Washing the external genitals does not remove sperm that have entered the vagina or reproductive tract. Douching can irritate tissue and may increase infection risk without providing contraception. Pregnancy prevention requires a method used before or around intercourse, such as condoms, hormonal contraception, an intrauterine device, or another suitable option. Emergency contraception may help after unprotected intercourse when used within the recommended window. Washing can support comfort and hygiene but is not a pregnancy treatment. Pregnancy testing follows the timing of the method and menstrual cycle.
\item[] \textbf{KID-084.} \textbf{Myth:} Urinating after intercourse prevents pregnancy.\newline
\textbf{Status:} False.\newline
\textbf{Fact:} Urine exits through the urethra, not the vagina, and does not remove sperm from the reproductive tract. Urination may reduce the chance of some urinary infections for some people, but it does not prevent conception. The urinary and reproductive tracts have separate openings and functions. Contraception must be selected according to pregnancy-prevention goals and used consistently. Emergency contraception is time-sensitive after unprotected sex. Burning, urgency, fever, pelvic pain, or blood in urine may indicate infection or another condition. Do not rely on urination as a contraceptive method.
\item[] \textbf{KID-085.} \textbf{Myth:} Infertility is usually the woman's problem.\newline
\textbf{Status:} False.\newline
\textbf{Fact:} Male and female factors both contribute substantially. Difficulty conceiving can involve sperm production or delivery, ovulation, fallopian tubes, uterus, age-related fertility, sexual function, or combinations of factors. Male factors contribute to a substantial proportion of infertile couples, and no cause is found in some cases. Evaluation should consider both partners rather than assigning blame to one person. Timing, duration of trying, age, medical history, and pregnancy history guide testing. Treatments range from addressing reversible factors to assisted reproduction. Respectful joint assessment improves both accuracy and support.
\item[] \textbf{KID-086.} \textbf{Myth:} Put butter or toothpaste on a burn.\newline
\textbf{Status:} False.\newline
\textbf{Fact:} These can interfere with appropriate burn care; cool running water is recommended for minor thermal burns. Butter, oils, toothpaste, and household substances can trap heat, irritate tissue, contaminate the wound, and complicate assessment. For a small thermal burn, cool running water for about twenty minutes as soon as practical can reduce tissue injury. Remove jewelry or tight clothing before swelling develops, but do not pull away material stuck to skin. Cover loosely with a clean non-stick dressing. Large, deep, electrical, chemical, facial, genital, hand, joint, or circumferential burns need urgent medical care. Do not break blisters or apply ice.
\item[] \textbf{KID-087.} \textbf{Myth:} Put ice directly on a burn.\newline
\textbf{Status:} False.\newline
\textbf{Fact:} Extreme cold can further injure tissue. Ice causes intense vasoconstriction and can produce additional cold injury in already damaged skin. It may numb pain while worsening tissue damage and delaying appropriate evaluation. Cool running water is safer for many minor thermal burns. Remove constricting objects and cover the area after cooling. Do not use very cold water on a large burn if it risks hypothermia, especially in a child. Blistering, severe pain, white or charred skin, or burns in high-risk locations require medical assessment.
\item[] \textbf{KID-088.} \textbf{Myth:} Make someone vomit whenever they have swallowed poison.\newline
\textbf{Status:} False.\newline
\textbf{Fact:} Induced vomiting can make some poisonings substantially more dangerous. Vomiting can re-expose the mouth and esophagus to corrosive chemicals and can cause aspiration into the lungs. The correct response depends on the substance, dose, time, symptoms, and person affected. Do not give home antidotes or induce vomiting unless poison-control or emergency professionals specifically instruct it. Keep the container or product information available and call the appropriate poison service promptly. Seizure, unconsciousness, breathing difficulty, or severe symptoms requires emergency services. Do not delay care while searching for a household remedy.
\item[] \textbf{KID-089.} \textbf{Myth:} Sucking venom from a snakebite removes the poison.\newline
\textbf{Status:} False.\newline
\textbf{Fact:} It does not provide effective treatment and can cause injury. Venom spreads through tissue and circulation, and mouth suction cannot remove a clinically useful amount. Cutting, sucking, ice, chemicals, and electric shocks can worsen tissue injury or delay definitive care. Keep the person calm and still, immobilize the affected limb in a functional position, and seek emergency help. Do not apply an improvised tight tourniquet unless directed by trained professionals using a region-specific protocol. Note the snake from a safe distance rather than trying to capture it. Antivenom and supportive treatment are hospital therapies when indicated.
\item[] \textbf{KID-090.} \textbf{Myth:} Cutting a snakebite helps venom escape.\newline
\textbf{Status:} False.\newline
\textbf{Fact:} Cutting can cause bleeding, infection, and tissue damage. Incisions do not remove a clinically useful amount of venom and can injure nerves, vessels, and tendons. They increase bleeding and infection risk and may complicate surgery or wound healing. Keep the person still and seek emergency medical care. Do not apply suction, ice, chemicals, or electric shock. Remove rings or tight items before swelling increases, if this can be done safely. Antivenom and supportive treatment are provided according to the snake and clinical effects.
\item[] \textbf{KID-091.} \textbf{Myth:} A tight tourniquet should routinely be tied above every snakebite.\newline
\textbf{Status:} False.\newline
\textbf{Fact:} Improvised arterial tourniquets can cause serious tissue injury. A tight tourniquet can stop arterial blood flow and cause ischemia, nerve injury, compartment syndrome, or tissue death. Some specialized pressure-immobilization protocols are used for selected snakebites in specific regions, but they require trained instruction and are not the same as an improvised arterial tourniquet. Keep the limb still and obtain emergency guidance. Do not cut, suck, or massage the bite. Record the time and observe symptoms without trying to capture the snake. Local poison or emergency protocols should guide first aid.
\item[] \textbf{KID-092.} \textbf{Myth:} Someone who is fainting should be forced to drink water immediately.\newline
\textbf{Status:} False.\newline
\textbf{Fact:} Giving liquids to someone with impaired consciousness risks aspiration. A fainting person may not be able to swallow safely and can inhale liquid into the lungs. Lay the person flat and raise the legs when injury is not suspected, loosen restrictive clothing, and monitor breathing. Turn them on their side if vomiting or not fully alert. Do not force food, drink, medicines, or smelling substances into the mouth. Emergency care is needed for prolonged unconsciousness, injury, chest pain, abnormal breathing, seizure, pregnancy, or recurrent fainting. Once fully alert, fluids may be offered if appropriate while the cause is assessed.
\item[] \textbf{KID-093.} \textbf{Myth:} A person having a heart attack should drive themselves to hospital.\newline
\textbf{Status:} False.\newline
\textbf{Fact:} Emergency medical transport is safer when available. Heart-attack symptoms can worsen suddenly and may include chest pressure, breathlessness, sweating, nausea, or pain in the arm, back, jaw, or upper abdomen. Emergency responders can begin assessment and treatment en route and alert the hospital. Driving risks a crash and delays care if the person loses consciousness. Call emergency services, keep the person resting, and follow dispatcher instructions. Aspirin may be appropriate for some adults but should be taken only according to local emergency guidance and contraindications. Do not wait for symptoms to become dramatic.
\item[] \textbf{KID-094.} \textbf{Myth:} If someone can move an injured limb, it cannot be fractured.\newline
\textbf{Status:} False.\newline
\textbf{Fact:} Movement can still be possible with a fracture. Some fractures are stable, partial, or located where movement remains possible. Pain, swelling, deformity, bruising, tenderness, or loss of function may be present, but none is required in every case. Do not test the limb repeatedly or ask the person to walk on it. Support it in the position found, apply wrapped cold for swelling, and arrange medical assessment. Numbness, pale or blue skin, severe deformity, open wounds, or absent pulses is urgent. Imaging may be needed even when function seems relatively preserved.
\item[] \textbf{KID-095.} \textbf{Myth:} Drinking cold water causes a respiratory infection.\newline
\textbf{Status:} False.\newline
\textbf{Fact:} Infectious respiratory illnesses are caused by pathogens, not water temperature. Colds and influenza result from viruses transmitted through respiratory particles, contact, and contaminated hands or surfaces. Cold water may briefly feel uncomfortable or trigger cough in a susceptible person, but it does not create an infection. Hydration can support comfort during illness, and the preferred temperature is a matter of tolerance. Fever, breathlessness, persistent symptoms, or dehydration requires appropriate care. Prevention depends on vaccination where available, ventilation, hand hygiene, and reducing exposure. Avoiding cold drinks is not a substitute for infection prevention.
\item[] \textbf{KID-096.} \textbf{Myth:} Ice cream directly causes tonsillitis.\newline
\textbf{Status:} False.\newline
\textbf{Fact:} Tonsillitis is usually infectious; cold food does not itself create the infection. Tonsillitis is commonly caused by viruses and sometimes bacteria such as group A Streptococcus. Ice cream cannot generate those organisms in the tonsils. Cold foods may soothe a sore throat for some people and irritate it for others. Antibiotics are useful only for selected bacterial infections after appropriate assessment. Difficulty breathing, drooling, inability to swallow fluids, severe one-sided swelling, or dehydration needs urgent care. Food temperature is not a diagnostic test.
\item[] \textbf{KID-097.} \textbf{Myth:} Bathing immediately after eating is medically dangerous.\newline
\textbf{Status:} False.\newline
\textbf{Fact:} There is no evidence-based general prohibition. Digestion continues while a person bathes, and ordinary bathing does not divert blood in a way that predictably causes drowning or dangerous cramps. A large meal may make vigorous activity feel uncomfortable, so personal comfort can guide timing. Water safety depends on supervision, swimming skill, temperature, alcohol use, and the environment. Infants and children require close supervision regardless of meals. People with fainting disorders or cardiovascular disease should follow individualized advice. A universal post-meal bathing ban is unsupported.
\item[] \textbf{KID-098.} \textbf{Myth:} Eating rice at night inevitably causes obesity.\newline
\textbf{Status:} False.\newline
\textbf{Fact:} Overall energy intake, diet quality, and metabolism matter more than one food at one time. Rice is a carbohydrate food and can fit into a balanced diet at any time of day. Weight gain reflects sustained energy surplus and is influenced by portions, sleep, activity, appetite, medications, and biology. Night eating may be associated with excess intake in some routines, but rice itself does not become uniquely fattening after dark. Whole grains, vegetables, protein, and portion size affect satiety and diet quality. People with diabetes may need individualized carbohydrate timing. Fear of one staple can reduce dietary adequacy without improving health.
\item[] \textbf{KID-099.} \textbf{Myth:} Ghee is harmless in unlimited amounts because it is traditional.\newline
\textbf{Status:} False.\newline
\textbf{Fact:} Ghee is calorie-dense and rich in saturated fat; quantity matters. Traditional use does not make a food physiologically unlimited. Ghee provides energy and substantial saturated fat, which can raise LDL cholesterol in many people. Small amounts may fit within an overall dietary pattern, depending on health needs and what they replace. Excess intake can contribute to calorie surplus and cardiometabolic risk. Replacing some saturated fat with unsaturated oils, nuts, seeds, or fish is generally more favorable for cardiovascular health. People with lipid disorders should discuss portions with a clinician or dietitian.
\item[] \textbf{KID-100.} \textbf{Myth:} Honey becomes poisonous when added to warm water.\newline
\textbf{Status:} False.\newline
\textbf{Fact:} Ordinary culinary use of honey in warm water does not create a poison. Heating or mixing honey with warm water does not transform it into a clinically established poison. Very high temperatures can change flavor and degrade some heat-sensitive compounds, but that is not the same as generating a toxic substance in ordinary use. Honey remains a concentrated sugar and should be portioned accordingly. Honey must not be given to infants under one year because of botulism risk, regardless of water temperature. People with diabetes or allergies may need individualized precautions. The temperature and amount matter more than folklore about ``toxins.''
\item[] \textbf{KID-101.} \textbf{Myth:} Papaya is universally dangerous during pregnancy.\newline
\textbf{Status:} False.\newline
\textbf{Fact:} Normal consumption of ripe papaya is not established as a cause of miscarriage. Ripe papaya eaten as food is not equivalent to concentrated extracts or large amounts of unripe papaya latex. Evidence does not support a universal ban on ordinary ripe fruit in pregnancy. Nutritional advice should distinguish preparation, quantity, hygiene, allergy, and individual medical conditions. Pregnant people should avoid known unsafe foods and follow local foodsafety guidance, but unnecessary restriction can reduce dietary variety. Supplements or herbal papaya products may have different concentrations and evidence. Concerns about bleeding, pain, or pregnancy loss require obstetric care, not dietary blame.
\item[] \textbf{KID-102.} \textbf{Myth:} Turmeric can cure cancer.\newline
\textbf{Status:} False.\newline
\textbf{Fact:} Laboratory research on curcumin does not establish turmeric as a cancer cure in humans. Curcumin has shown biochemical effects in cells and animal models, but laboratory activity does not establish a safe, effective human cancer treatment. Turmeric food use and concentrated supplements differ in dose, absorption, interactions, and quality. Delaying surgery, chemotherapy, radiation, targeted therapy, or immunotherapy for an unproven cure can allow disease to progress. Supplements may affect bleeding, liver function, or drug metabolism. Patients should tell their oncology team about all products. Integrative approaches may support symptoms or wellbeing when evaluated alongside standard treatment, not instead of it.
\item[] \textbf{KID-103.} \textbf{Myth:} Ayurvedic or herbal products are automatically safe because they are natural.\newline
\textbf{Status:} False.\newline
\textbf{Fact:} Products can cause adverse effects, drug interactions, contamination, and in some cases heavy-metal toxicity. ``Natural'' describes origin, not safety, purity, dose, or effectiveness. Herbs can have pharmacologic effects, alter liver enzymes, increase bleeding, affect blood pressure or glucose, and interact with prescription medicines. Products may be contaminated, adulterated, mislabeled, or variable between batches. Some traditional preparations have contained harmful metals or other undeclared ingredients. Quality-control rules differ by country and are not the same as approval of efficacy. Discussing all products with a clinician or pharmacist helps prevent avoidable harm.

\section{Candidate Myths Requiring Provenance Review}

The audit identified the following common claims that were not represented elsewhere in the book. These entries are retained only as candidate myths pending entry-level provenance review; secondary myth lists are leads for investigation, not proof that a claim is widespread or that its correction is complete \cite{ref16,ref17,ref18,ref36,ref15,ref59,ref60,ref61}.
\item[] \textbf{CAN-001.} \textbf{Myth:} Cell phones cause cancer.\newline
\textbf{Status:} Uncertain.\newline
\textbf{Fact:} Current evidence has not established ordinary cell-phone use as a proven cause of cancer. Mobile phones emit radiofrequency energy, a form of non-ionizing radiation that differs from X-rays and does not have enough energy to directly break DNA bonds. Large studies have not established a consistent causal relationship between ordinary use and cancer, but research continues because exposure is widespread and technologies change. Risk estimates are complicated by recall bias, changing devices, and long latency periods. Using speaker mode or a headset can reduce local exposure if someone remains concerned. These precautions should not replace evidence-based cancer prevention and screening. The current conclusion is no proven routine causal link, with residual scientific uncertainty.
\item[] \textbf{CAN-002.} \textbf{Myth:} Power lines directly cause cancer.\newline
\textbf{Status:} Uncertain.\newline
\textbf{Fact:} A causal relationship has not been established; exposure and risk questions require careful interpretation of the evidence. Power lines produce extremely low-frequency electromagnetic fields, not ionizing radiation. Epidemiologic studies have explored a possible association with childhood leukemia, but findings are inconsistent and do not establish causation. Exposure falls rapidly with distance and varies by line, current, and environment. Mechanistic evidence for a carcinogenic effect at typical residential levels is limited. People should avoid unnecessary exposure to electrical hazards, but ordinary proximity is not a proven direct cancer cause. The claim is too absolute for the current evidence.
\item[] \textbf{CAN-003.} \textbf{Myth:} A positive attitude determines whether a person survives cancer.\newline
\textbf{Status:} False.\newline
\textbf{Fact:} Attitude may affect coping and wellbeing, but it does not determine tumor biology or replace treatment. Cancer behavior is governed by tumor genetics, stage, tissue, immune interactions, and response to treatment. Optimism, support, and psychological care can improve coping, communication, and quality of life. They do not reliably shrink tumors or determine survival independently of disease and treatment factors. Blaming a person for progression creates guilt and can undermine informed care. Patients should be offered emotional support without pressure to perform positivity. Evidence-based oncology remains central.
\item[] \textbf{CAN-004.} \textbf{Myth:} No family history means no cancer risk.\newline
\textbf{Status:} False.\newline
\textbf{Fact:} Many cancers arise without a known inherited family history. Most cancers are influenced by acquired cellular changes and environmental or lifestyle factors rather than a recognized inherited syndrome. Family history can increase risk, but its absence does not eliminate ordinary population risk. Families may also be small, unaware of diagnoses, or carry variants with incomplete expression. Screening recommendations are based on age and individual risk, not only relatives. New symptoms require assessment regardless of family history. Genetic counseling is useful when personal or family patterns suggest inherited risk.
\item[] \textbf{CAN-005.} \textbf{Myth:} Microwave-safe plastic containers always release cancer-causing substances.\newline
\textbf{Status:} False.\newline
\textbf{Fact:} Containers labeled for microwave use should be used as directed; unsuitable or damaged plastics should not be microwaved. Food-contact materials are evaluated for intended conditions of use, but safety depends on the container, temperature, time, and whether it is designed for microwave heating. Heating non-food-grade, damaged, or single-use containers can increase migration of chemicals or cause melting. Glass or ceramic labeled for heat use may be a practical alternative. Do not microwave containers that are sealed unless designed for it, because pressure can build. Avoid heating food in packaging not intended for that use. The claim that every microwave-safe plastic is automatically dangerous is unsupported \cite{ref6}.
\item[] \textbf{CAN-006.} \textbf{Myth:} Cancer treatment is always worse than the disease.\newline
\textbf{Status:} False.\newline
\textbf{Fact:} Treatment benefits and harms vary by cancer, stage, treatment, and patient; decisions require individualized discussion. Surgery, radiation, chemotherapy, immunotherapy, and targeted therapy can cause substantial side effects, but they may cure disease, extend life, prevent complications, or relieve symptoms. Some cancers require urgent treatment, while others can be monitored or treated less intensively. Benefit-harm decisions depend on goals, tumor biology, health, and available alternatives. Supportive care can prevent or reduce many treatment effects. Refusing all treatment based on a general slogan may remove beneficial options. Shared decision-making allows informed choices.
\item[] \textbf{CAN-007.} \textbf{Myth:} Cancer screening always causes more harm than benefit.\newline
\textbf{Status:} False.\newline
\textbf{Fact:} Screening has both benefits and harms, and recommendations target groups in which expected benefits outweigh those harms. Screening can detect disease early and reduce mortality for selected cancers and risk groups. It can also produce false positives, overdiagnosis, anxiety, radiation exposure, and unnecessary procedures. The balance depends on the test, age, risk, interval, and health-system context. Recommendations are therefore not ``screen everyone for everything.'' People should discuss personal risk, benefits, limitations, and alternatives. A screening test is not the same as diagnostic evaluation of symptoms.
\item[] \textbf{CAN-008.} \textbf{Myth:} A Pap smear detects ovarian cancer.\newline
\textbf{Status:} False.\newline
\textbf{Fact:} Cervical screening evaluates the cervix and is not a screening test for ovarian cancer. A Pap test samples cervical cells to detect precancerous or cancerous changes related mainly to HPV. It does not reliably examine the ovaries or detect ovarian cancer. Ovarian cancer can have nonspecific symptoms such as bloating, pelvic pain, early satiety, or urinary change, but symptoms alone are not diagnostic. People with concerning symptoms or inherited risk need appropriate clinical evaluation. Screening strategies differ from diagnostic tests. Keeping cervical screening appointments does not remove the need to report new pelvic symptoms.
\item[] \textbf{CAN-009.} \textbf{Myth:} Only older women develop ovarian cancer.\newline
\textbf{Status:} False.\newline
\textbf{Fact:} Risk varies by cancer type, age, and inherited factors, and younger women can also develop ovarian cancer. Ovarian cancer risk generally rises with age, but germ-cell, sex-cord, and epithelial tumors can occur in younger people. Inherited variants such as BRCA-related risk can change the age and pattern of disease. Persistent abdominal distension, pelvic pain, early satiety, or urinary symptoms should be assessed when new or recurrent. These symptoms usually have benign causes, so evaluation is not a diagnosis of cancer. Family history should be discussed to determine whether genetic counseling is appropriate. Age alone should not dismiss a concerning pattern.
\item[] \textbf{CAN-010.} \textbf{Myth:} Vaping is completely harmless.\newline
\textbf{Status:} False.\newline
\textbf{Fact:} Vaping can expose users to nicotine and other substances; it is not risk-free, especially for nonsmokers and young people. E-cigarette aerosols may contain nicotine, ultrafine particles, flavor chemicals, metals, and other substances. Nicotine can cause dependence and can affect the developing adolescent brain and pregnancy. Vaping may expose users to fewer toxic chemicals than combustible cigarettes in some comparisons, but lower risk is not no risk. Long-term effects remain under study. Nonsmokers should not start, and people who smoke should use evidence-based cessation support. Complete switching may be harm reduction for some adult smokers, not a health endorsement.
\item[] \textbf{CAN-011.} \textbf{Myth:} Nicotine causes most of the disease caused by cigarette smoking.\newline
\textbf{Status:} False.\newline
\textbf{Fact:} Nicotine is addictive, while many toxic chemicals in tobacco smoke cause much of the cancer and lung damage. Nicotine is the main driver of dependence and can affect cardiovascular function, pregnancy, and the developing brain. Combustible tobacco smoke also contains thousands of chemicals, including carcinogens and substances that injure lungs and blood vessels. Nicotine keeps people using tobacco, but it is not the sole explanation for smoke-related cancer and chronic lung disease. This distinction matters when comparing nicotine replacement with smoking. Nicotine medicines are substantially safer than cigarettes when used as directed, but they are not risk-free for everyone. Quitting all tobacco and nicotine remains the healthiest goal.
\item[] \textbf{CAN-012.} \textbf{Myth:} Vaping always causes popcorn lung.\newline
\textbf{Status:} False.\newline
\textbf{Fact:} The claim is not established; bronchiolitis obliterans was associated with specific occupational chemical exposures, not a universal effect of vaping. ``Popcorn lung'' refers to bronchiolitis obliterans, a rare obstructive lung disease linked historically to intense exposure to certain chemicals. There is no evidence that every person who vapes develops this condition. That does not make vaping safe: aerosols can irritate airways and contain nicotine, particles, metals, and other chemicals. Persistent cough, wheeze, chest pain, or breathlessness needs medical assessment. Avoiding exaggerated claims improves credibility of genuine risk communication. Nonsmokers should not begin vaping.
\item[] \textbf{CAN-013.} \textbf{Myth:} Switching completely from cigarettes to vaping provides no health benefit.\newline
\textbf{Status:} Uncertain.\newline
\textbf{Fact:} Complete switching can reduce exposure to many toxic chemicals compared with continued smoking, but vaping is not harmless and quitting all tobacco and nicotine is preferable when possible. For adults who smoke, complete substitution may reduce exposure to many combustion toxicants compared with continuing cigarettes. The benefit depends on actually stopping combustible tobacco; dual use may preserve substantial risk. Vaping still exposes users to nicotine and other substances, and long-term effects remain incompletely known. It is not appropriate as a reason for nonsmokers or youth to start. Evidence-based cessation medicines and counseling have established roles. The most health-protective endpoint is stopping cigarettes and, when feasible, all nicotine.
\item[] \textbf{CAN-014.} \textbf{Myth:} Everyone must sleep exactly eight hours.\newline
\textbf{Status:} False.\newline
\textbf{Fact:} Sleep needs vary by age and individual factors, although persistent insufficient sleep can harm health. Sleep recommendations are ranges that vary by age, development, health, and individual biology. Many adults function best around seven to nine hours, but an exact eight-hour requirement is not universal. Quality, regularity, timing, and daytime function also matter. Persistent short sleep is associated with cardiometabolic, mood, cognitive, and safety risks. Excessive sleep or unrefreshing sleep can also signal disease. A consistent schedule and evaluation of snoring, insomnia, or daytime sleepiness are more useful than chasing one number.
\item[] \textbf{CAN-015.} \textbf{Myth:} Extra weekend sleep completely repairs chronic sleep loss.\newline
\textbf{Status:} False.\newline
\textbf{Fact:} Additional sleep may help recovery, but it does not reliably erase the effects of ongoing sleep deprivation. Sleeping longer after a short night can reduce sleepiness and improve some performance measures. It does not reliably reverse all metabolic, cognitive, mood, or circadian effects of repeated restriction. Large shifts between weekday and weekend schedules can create social jet lag and make Monday sleep harder. Regular adequate sleep is more protective than repeated catch-up cycles. Persistent daytime sleepiness, snoring, or insomnia deserves evaluation. Sleep extension is helpful recovery, not a complete reset button.
\item[] \textbf{CAN-016.} \textbf{Myth:} Muscle turns into fat when a person stops exercising.\newline
\textbf{Status:} False.\newline
\textbf{Fact:} Muscle and fat are different tissues; reduced activity can cause muscle loss and fat gain separately. Muscle fibers do not transform chemically into adipose tissue. With inactivity, muscle size and strength can decline, while lower energy expenditure or continued intake may increase fat stores. These processes can happen together and create the appearance of one tissue changing into the other. Resistance training, adequate protein, activity, and gradual progression help preserve muscle. Illness, age, and medication may alter the response. Body-composition change is real, but it is not tissue conversion.
\item[] \textbf{CAN-017.} \textbf{Myth:} Everyone must have one bowel movement every day.\newline
\textbf{Status:} False.\newline
\textbf{Fact:} Healthy bowel frequency varies; persistent change, bleeding, severe pain, or unexplained weight loss requires evaluation. Normal frequency can range from several movements per day to several per week, depending on the person and diet. Constipation is defined more by hard stools, difficult passage, straining, or incomplete evacuation than by a calendar alone. Hydration, fiber, activity, medicines, and bowel disorders influence patterns. A sudden persistent change, blood, anemia, vomiting, severe pain, or weight loss needs assessment. Laxative use should be appropriate to the cause. Individual regularity matters more than a universal daily rule.
\item[] \textbf{CAN-018.} \textbf{Myth:} Organic food is always more nutritious than conventionally grown food.\newline
\textbf{Status:} False.\newline
\textbf{Fact:} Nutritional differences are not consistently large; diet quality depends more on the overall pattern and adequate intake of plant foods. Organic standards concern production methods and permitted inputs, not a guarantee of superior nutrient content. Nutrient levels vary with soil, variety, ripeness, storage, and processing in both systems. Organic foods can still be high in sugar, salt, or calories when highly processed. Conventional produce is nutritious and should not be avoided because it is not organic. Washing produce reduces some contamination risk, though it cannot remove every hazard. Affordability and adequate intake are important health considerations.
\item[] \textbf{CAN-019.} \textbf{Myth:} Red meat must always be avoided.\newline
\textbf{Status:} False.\newline
\textbf{Fact:} Risk depends on amount, processing, cooking, and the overall dietary pattern; processed meat and high intake deserve particular caution. Unprocessed red meat can provide protein, iron, zinc, and vitamin B12. Higher intake, especially processed meat, is associated with increased cardiometabolic and colorectal-cancer risk in population studies. Replacing some servings with legumes, fish, nuts, or poultry may improve dietary risk profiles. Cooking methods and portion size matter, and charred meat can contain harmful compounds. Individual choices should consider culture, access, nutritional needs, and medical conditions. ``Never'' is less evidence-based than moderation and overall pattern.
\item[] \textbf{CAN-020.} \textbf{Myth:} Coffee permanently stunts growth.\newline
\textbf{Status:} False.\newline
\textbf{Fact:} Coffee does not permanently stunt growth, although caffeine can affect sleep and may require caution in children and adolescents. There is no established evidence that ordinary coffee permanently stops linear growth. Caffeine can cause insomnia, anxiety, palpitations, reflux, and dependence, with greater sensitivity in children and adolescents. Poor sleep and excessive stimulant use can indirectly affect health and performance. Added sugar and cream can substantially change calorie content. Pregnancy and certain heart, anxiety, or medication conditions require individualized limits. The real concern is dose and context, not a universal growth-stunting effect.
\item[] \textbf{CAN-021.} \textbf{Myth:} Only girls develop eating disorders.\newline
\textbf{Status:} False.\newline
\textbf{Fact:} Eating disorders can affect people of every sex and gender. Eating disorders occur in boys, men, girls, women, and gender-diverse people. Some groups are underdiagnosed because stereotypes influence recognition and help-seeking. Symptoms can include restriction, binge eating, purging, compulsive exercise, fear of weight gain, or distress about eating and body shape. Medical complications can occur at any body size. Screening and assessment should focus on behavior, thoughts, physical signs, and function rather than gender assumptions. Early treatment improves the chance of recovery.
\item[] \textbf{CAN-022.} \textbf{Myth:} You can identify an eating disorder from a person's appearance.\newline
\textbf{Status:} False.\newline
\textbf{Fact:} Eating disorders occur across body sizes and cannot be reliably diagnosed by appearance. A person may have a serious eating disorder while appearing thin, average-weight, or larger-bodied. Weight can change, remain stable, or be obscured by fluid shifts and compensatory behaviors. Warning signs include secretive eating, rigid food rules, vomiting, laxative misuse, excessive exercise, fainting, or intense distress around food. Appearance-based reassurance can delay care and reinforce stigma. Assessment includes history, mental health, vital signs, laboratory tests, and physical complications. Concern should prompt supportive conversation and professional help.
\item[] \textbf{CAN-023.} \textbf{Myth:} Eating disorders are a lifestyle choice.\newline
\textbf{Status:} False.\newline
\textbf{Fact:} They are serious mental illnesses involving biological, psychological, and social factors. Eating disorders involve disturbed eating behavior and body-related thoughts that can impair health and function. Genetic vulnerability, temperament, trauma, social pressures, neurobiology, and environment may contribute. A person may feel ambivalent about change, but ambivalence is part of illness rather than proof of a casual choice. Malnutrition can affect the heart, bones, brain, hormones, and immune system. Treatment combines nutritional rehabilitation, medical monitoring, and evidence-based psychotherapy. Blame and coercive comments can worsen secrecy.
\item[] \textbf{CAN-024.} \textbf{Myth:} Eating disorders usually resolve without treatment.\newline
\textbf{Status:} False.\newline
\textbf{Fact:} Specialist assessment and treatment can be necessary, and early support improves the chance of recovery. Some symptoms fluctuate, but untreated eating disorders can become chronic and medically dangerous. Complications include electrolyte disturbance, arrhythmia, osteoporosis, infertility, gastrointestinal injury, and suicide risk. Treatment may involve primary care, mental-health professionals, dietitians, family support, and hospital care when medically unstable. Recovery is possible and does not require reaching a particular weight before help begins. Avoiding comments about appearance and offering practical support can facilitate care. Fainting, chest symptoms, severe dehydration, confusion, or suicidal thoughts require urgent assessment.

\section{Infection, Medication, Urinary, and Symptom Myths}

The following common claims were identified in additional Mayo Clinic, NHS, and Cleveland Clinic guidance \cite{ref37,ref38,ref39}.
\item[] \textbf{INF-001.} \textbf{Myth:} ``Starve a fever, feed a cold'' is a sound treatment rule.\newline
\textbf{Status:} False.\newline
\textbf{Fact:} Ill people should receive appropriate fluids and nutrition as tolerated; deliberately starving or forcing food is not a universal treatment. Fever and colds are caused by different illnesses, but neither requires deliberate starvation. Fluids help reduce dehydration, especially with fever, vomiting, diarrhea, or rapid breathing. Appetite may fall temporarily, and small, nutritious meals are reasonable as tolerated. Forcing food can worsen nausea, while refusing fluids can be dangerous. Infants, older adults, and people with chronic disease need closer attention to hydration and intake. Persistent fever, confusion, breathing difficulty, or inability to drink needs medical care.
\item[] \textbf{INF-002.} \textbf{Myth:} Milk or dairy products increase mucus during a cold.\newline
\textbf{Status:} False.\newline
\textbf{Fact:} Milk does not cause the body to produce extra phlegm, although its texture may briefly feel thick in the mouth or throat. The creamy sensation after drinking milk can be mistaken for increased mucus production. Controlled observations do not show that ordinary dairy intake increases respiratory secretions in most people. Individuals with milk allergy, lactose intolerance, or reflux may have separate symptoms. Dairy can provide calories, protein, and fluids during illness if tolerated. Avoiding it is unnecessary unless it causes a personal reaction or a clinician advises otherwise. Breathing difficulty or persistent cough requires assessment.
\item[] \textbf{INF-003.} \textbf{Myth:} Vitamin C prevents all colds.\newline
\textbf{Status:} False.\newline
\textbf{Fact:} Vitamin C is important for health, but routine supplementation does not guarantee that a person will avoid a cold. Vitamin C deficiency impairs health, but extra vitamin C does not create complete immunity to respiratory viruses. Regular supplementation may have modest effects on duration or severity in some populations, while starting it after symptoms may not help much. High doses can cause gastrointestinal symptoms and increase kidney-stone risk in susceptible people. Fruits and vegetables are generally good sources. Vaccination, ventilation, hand hygiene, and staying home when ill address transmission more directly. Supplements should not replace medical care.
\item[] \textbf{INF-004.} \textbf{Myth:} Cranberry juice cures an established urinary-tract infection.\newline
\textbf{Status:} False.\newline
\textbf{Fact:} Cranberry products may help prevent some recurrent infections, but they do not reliably clear an established UTI. Cranberry compounds may reduce bacterial attachment in some circumstances, but evidence for prevention is mixed and product-dependent. They do not substitute for testing and antibiotics when a bacterial UTI requires treatment. Delayed treatment can permit kidney infection, especially in pregnancy, children, older adults, or people with obstruction. Fever, flank pain, vomiting, or systemic illness requires prompt care. Sugary juice can add substantial calories and does not improve treatment. Prevention advice should be individualized.
\item[] \textbf{INF-005.} \textbf{Myth:} Drinking urine is therapeutic.\newline
\textbf{Status:} False.\newline
\textbf{Fact:} Urine is not an antiseptic or a treatment and can worsen dehydration or introduce contamination. Urine contains water plus waste products, electrolytes, and metabolites that the body is excreting. Re-ingesting it does not cleanse the body or treat infection, cancer, poisoning, or dehydration. Concentrated urine can worsen fluid and salt imbalance, especially when a person is already dehydrated. It may also be contaminated by bacteria, blood, medicines, or chemicals. Suspected poisoning or illness requires evidence-based emergency care. Do not substitute urine for safe drinking water or prescribed treatment.
\item[] \textbf{INF-006.} \textbf{Myth:} Healthy urine is completely sterile.\newline
\textbf{Status:} False.\newline
\textbf{Fact:} Modern testing shows that normal urine can contain small numbers of microorganisms. Sensitive sequencing and culture methods identify a urinary microbiome in some healthy people, especially in the lower urinary tract. This does not mean that every organism is causing infection or that a positive molecular signal requires antibiotics. Clinical UTI diagnosis combines symptoms, examination, urinalysis, and selected cultures. Contamination during collection can also affect results. Blood, pain, fever, foul odor, or urinary symptoms need appropriate assessment. ``Not sterile'' should not become a reason for unnecessary antimicrobial treatment.
\item[] \textbf{INF-007.} \textbf{Myth:} Antibacterial soap is always better than ordinary soap and water.\newline
\textbf{Status:} False.\newline
\textbf{Fact:} Proper handwashing with ordinary soap and running water is effective for routine hygiene. Soap removes oils, dirt, and microorganisms from skin when combined with friction and running water. For household handwashing, antibacterial additives have not shown a universal advantage over plain soap. Alcohol hand sanitizer is useful when soap and water are unavailable, except when hands are visibly dirty or certain pathogens are involved. Overuse of antimicrobial products can irritate skin and promote misleading confidence. Handwashing technique and duration matter more than marketing labels. Healthcare settings may use specialized products for specific purposes.
\item[] \textbf{INF-008.} \textbf{Myth:} Chicken soup cures a cold.\newline
\textbf{Status:} False.\newline
\textbf{Fact:} Warm fluids and soup may ease symptoms and support hydration, but they do not eradicate the virus causing a cold. Soup can provide fluid, salt, energy, and warmth when appetite is reduced. Steam and warm liquid may temporarily soothe congestion or throat discomfort. These supportive effects do not kill the respiratory virus or guarantee faster recovery. Viral colds generally improve with time, rest, hydration, and symptom-directed care. Persistent fever, shortness of breath, chest pain, or worsening illness needs evaluation. Chicken soup is a comfort food, not an antiviral medicine.
\item[] \textbf{INF-009.} \textbf{Myth:} Over-the-counter cough and cold medicines are automatically safe for young children.\newline
\textbf{Status:} False.\newline
\textbf{Fact:} Some products can cause serious adverse effects or dosing errors in young children; age-specific medical advice is important. Combination products can contain duplicate ingredients and cause overdose when labels are misunderstood. Some cough and cold ingredients have limited benefit and can produce sedation, agitation, fast heart rate, or breathing problems in young children. Never use an adult product or household spoon for dosing. Read active ingredients and follow age-specific instructions or professional advice. Saline, fluids, humidified air, and honey for children over one year may help selected symptoms. Breathing difficulty, dehydration, or unusual sleepiness requires urgent care.
\item[] \textbf{INF-010.} \textbf{Myth:} Zinc or other supplements are universally needed whenever someone has a cold.\newline
\textbf{Status:} False.\newline
\textbf{Fact:} Supplement benefits depend on the substance, dose, timing, deficiency status, and patient factors. Zinc lozenges may modestly shorten a cold for some adults when started early, but formulations and doses vary and side effects are possible. Intranasal zinc products have been linked to loss of smell and should be avoided. Vitamin supplements do not guarantee prevention or cure, especially when nutritional status is adequate. Excess zinc can cause nausea and copper deficiency, while other supplements interact with medicines. Food-first nutrition is usually sufficient. Persistent or severe symptoms require diagnosis rather than escalating supplements.
\item[] \textbf{INF-011.} \textbf{Myth:} Every fever should be treated until body temperature is completely normal.\newline
\textbf{Status:} False.\newline
\textbf{Fact:} Fever treatment is generally directed at discomfort and the underlying illness, with special attention to age, appearance, hydration, and cause. Fever is a regulated immune response and does not always need to be eliminated. Acetaminophen or ibuprofen may improve comfort when appropriate, but they do not cure the infection. Dosing errors and duplicate products can cause serious harm. Infants, immunocompromised people, and patients with concerning appearance need individualized assessment. Confusion, breathing difficulty, stiff neck, dehydration, seizure, or persistent high fever warrants medical care. Treating the person, not just the number, is the safer principle \cite{ref8}.
\item[] \textbf{INF-012.} \textbf{Myth:} Pulling a gray hair makes more gray hairs grow.\newline
\textbf{Status:} False.\newline
\textbf{Fact:} Plucking one hair does not cause neighboring follicles to turn gray, although the removed hair will usually grow back gray. Hair color is determined by pigment production within each follicle. Removing one hair does not send a signal that turns adjacent follicles gray. The same follicle will generally produce another gray hair because its pigment biology has not changed. Repeated plucking can injure follicles, cause ingrown hairs, or contribute to thinning in that area. Graying is influenced by genetics, age, and some medical or nutritional conditions. Sudden patchy changes deserve assessment rather than repeated plucking.

\section{Important High-Impact Myths}

The following high-impact myths are included because they can lead to preventable medication errors, delayed care, or avoidable harm.
\item[] \textbf{IMP-001.} \textbf{Myth:} The flu vaccine gives you influenza.\newline
\textbf{Status:} False.\newline
\textbf{Fact:} Flu vaccines cannot cause influenza illness, although a vaccinated person can still catch another respiratory infection or influenza not well matched by the vaccine \cite{ref40}. Injectable inactivated flu vaccines cannot reproduce influenza infection, and recombinant vaccines contain no live flu virus. The nasal live-attenuated vaccine uses weakened virus and is designed not to cause ordinary influenza in healthy eligible recipients. Soreness, fatigue, or low fever can occur as immune effects. A person can become infected before immunity develops, by a strain not covered well, or with another respiratory virus. Vaccination reduces risk and severity but is not perfect. Serious allergic reactions are rare and require appropriate medical advice.
\item[] \textbf{IMP-002.} \textbf{Myth:} A COVID-19 vaccine changes a person's DNA or causes infertility.\newline
\textbf{Status:} False.\newline
\textbf{Fact:} COVID-19 vaccines do not alter human DNA, and available evidence has not established that vaccination causes fertility problems \cite{ref41}. Vaccine instructions such as mRNA remain outside the cell nucleus and are broken down after use; they do not rewrite human DNA. Studies have not established a general reduction in fertility from vaccination. Temporary fever or menstrual-cycle changes can occur, while COVID-19 infection itself can affect health and pregnancy. Recommendations may differ for pregnancy, prior reactions, and changing products. People with specific fertility or pregnancy concerns should discuss them with a clinician. Misinformation should not delay protective care.
\item[] \textbf{IMP-003.} \textbf{Myth:} Generic medicines are weaker than brand-name medicines.\newline
\textbf{Status:} False.\newline
\textbf{Fact:} Approved generics contain the same active ingredient and must meet standards for strength, quality, and performance comparable to the brand product \cite{ref42}. Generic medicines must contain the same active ingredient and demonstrate comparable delivery to the body within regulatory standards. They may differ in color, shape, inactive ingredients, packaging, or price. Those differences do not ordinarily make the active treatment weaker. Narrow-therapeutic-index medicines and formulation changes may require careful clinical or pharmacy oversight. A patient with a new reaction should seek review rather than assume all generics fail. Use only medicines from reliable, regulated sources.
\item[] \textbf{IMP-004.} \textbf{Myth:} Paracetamol or acetaminophen is harmless at any dose.\newline
\textbf{Status:} False.\newline
\textbf{Fact:} Taking too much, or combining several products that contain it, can cause severe liver injury or death \cite{ref43}. Acetaminophen is effective and safe at recommended doses for many people, but overdose can overwhelm liver metabolism. Combination cold, pain, and prescription products may contain the same ingredient under different names. Alcohol use, fasting, liver disease, and repeated supratherapeutic dosing can increase risk. Do not exceed the label or prescribed limit, and check active ingredients before combining products. Suspected overdose requires immediate poison-control or emergency advice even if the person feels well. Early treatment can be lifesaving.
\item[] \textbf{IMP-005.} \textbf{Myth:} If a prescribed medicine makes you feel better, you can stop it immediately.\newline
\textbf{Status:} False.\newline
\textbf{Fact:} Some medicines must be tapered or continued for a specified course; stopping treatment should be discussed with the prescriber or pharmacist. Symptom improvement does not always mean the underlying disease has resolved. Abrupt withdrawal can cause rebound symptoms, seizures, adrenal crisis, blood-pressure changes, withdrawal effects, or relapse depending on the drug. Some antibiotics, steroids, antidepressants, antiseizure medicines, and cardiovascular drugs require specific plans. Other medicines are intended for symptom-limited use, so the correct duration differs. Read the instructions and ask before changing treatment. Urgent adverse effects should prompt immediate professional advice rather than unsupervised stopping.
\item[] \textbf{IMP-006.} \textbf{Myth:} Harder tooth-brushing cleans better.\newline
\textbf{Status:} False.\newline
\textbf{Fact:} Gentle brushing with a soft-bristled brush and fluoride toothpaste is preferred; forceful brushing can irritate gums and wear tooth surfaces \cite{ref44}. Plaque is removed by controlled technique at the gumline, not by damaging pressure. Hard scrubbing can cause gum recession, abrasion, sensitivity, and damage to restorations. A soft-bristled brush used twice daily with fluoride toothpaste is generally recommended. Replace a frayed brush and ask a dental professional to demonstrate technique when unsure. Brushing longer or harder does not compensate for missed areas or lack of interdental cleaning. Persistent bleeding or pain needs dental assessment.
\item[] \textbf{IMP-007.} \textbf{Myth:} Fluoride toothpaste is dangerous and should be avoided.\newline
\textbf{Status:} False.\newline
\textbf{Fact:} Used in the recommended amount and not swallowed, fluoride toothpaste helps prevent tooth decay \cite{ref45}. Fluoride strengthens enamel and helps reverse early mineral loss caused by bacterial acids. The appropriate amount is safe when brushed onto teeth and spat out. Young children need age-appropriate supervision because repeated swallowing of large amounts can cause fluorosis or toxicity. Fluoride concentration and local water guidance vary, so dental advice should be followed. Avoiding fluoride increases preventable decay risk for many people. Severe ingestion requires poison-control guidance.
\item[] \textbf{IMP-008.} \textbf{Myth:} Baby teeth do not matter because they eventually fall out.\newline
\textbf{Status:} False.\newline
\textbf{Fact:} Baby teeth support eating, speech, and normal development, and untreated decay can cause pain and infection. Primary teeth preserve space for permanent teeth and help children chew, speak, and develop oral habits. Decay can cause pain, sleep disruption, nutrition problems, infection, and missed school. Early loss may affect spacing and require dental management. Brushing with age-appropriate fluoride toothpaste, limiting frequent sugar exposure, and routine dental care reduce risk. A developing permanent tooth can be affected by infection in an overlying primary tooth. Temporary teeth still deserve permanent-quality care.
\item[] \textbf{IMP-009.} \textbf{Myth:} Tilt the head backward during a nosebleed.\newline
\textbf{Status:} False.\newline
\textbf{Fact:} Sit upright, lean slightly forward, and apply firm continuous pressure to the soft part of the nose; seek urgent care for heavy or persistent bleeding. Leaning backward allows blood to flow into the throat, causing swallowing, coughing, nausea, or aspiration. Sit forward, breathe through the mouth, and pinch the soft nose continuously for about ten to fifteen minutes without repeatedly checking. Spitting out blood is safer than swallowing it. Avoid forceful blowing immediately afterward. Emergency care is needed for heavy bleeding, faintness, breathing difficulty, anticoagulant use, major injury, or bleeding that does not stop. Recurrent episodes warrant evaluation.
\item[] \textbf{IMP-010.} \textbf{Myth:} People with mental illness are generally dangerous.\newline
\textbf{Status:} False.\newline
\textbf{Fact:} Mental illness does not by itself make a person violent; risk depends on individual circumstances, and stigma can discourage people from seeking care. Most people with mental illness are not violent and are more likely to experience violence or neglect than cause it. Risk is shaped by specific symptoms, substance use, history, crisis conditions, access to weapons, and social context---not diagnosis alone. Stigma can delay treatment, housing, employment, and help-seeking. Safety assessment should focus on current behavior and credible warning signs. Compassionate care and practical support improve outcomes. Immediate threats require emergency help based on behavior, not a label.
\item[] \textbf{IMP-011.} \textbf{Myth:} Antidepressants are addictive in the same way as recreational drugs.\newline
\textbf{Status:} False.\newline
\textbf{Fact:} They are not generally addictive, but some can cause discontinuation symptoms and should not be stopped abruptly without medical advice. Addiction involves compulsive use, craving, loss of control, and continued use despite harm; antidepressants do not generally produce that pattern. Some medicines can cause discontinuation symptoms such as dizziness, irritability, flu-like feelings, sleep disturbance, or electric-shock sensations when stopped suddenly. They may also have side effects, interactions, and delayed therapeutic benefits. Tapering is individualized and should be planned with the prescriber. Relapse after stopping is not proof of addiction. Seek urgent help for severe mood change, mania, or suicidal thoughts.
\item[] \textbf{IMP-012.} \textbf{Myth:} Therapy is only for people who are weak or unable to cope.\newline
\textbf{Status:} False.\newline
\textbf{Fact:} Psychotherapy is an evidence-based treatment and support that can help people with many different levels of distress. Therapy teaches skills and provides structured treatment for conditions such as depression, anxiety, trauma, addiction, pain, and relationship difficulties. Seeking help can reflect insight and active problem-solving rather than weakness. Different therapies have different goals, methods, and evidence, so fit and training matter. Therapy can be used alone or with medication and medical care. It is not a guarantee of instant improvement, and progress requires collaboration. Stigma should not prevent a person from accessing support.
\item[] \textbf{IMP-013.} \textbf{Myth:} Spicy food causes peptic ulcers.\newline
\textbf{Status:} False.\newline
\textbf{Fact:} Most peptic ulcers are associated with Helicobacter pylori infection or medicines such as nonsteroidal anti-inflammatory drugs; spicy food may worsen symptoms for some people but is not the usual cause. Ulcers are breaks in the stomach or duodenal lining caused mainly by H. pylori, NSAIDs, and other medical factors. Spicy foods can increase burning or discomfort in some people but do not usually create the ulcer. Treatment may require eradication antibiotics, acid suppression, stopping an offending medicine, or evaluation for complications. Black stools, vomiting blood, faintness, severe sudden pain, or anemia require urgent care. Avoiding personal triggers can improve comfort but does not replace diagnosis. Long-term NSAID use should be reviewed.
\item[] \textbf{IMP-014.} \textbf{Myth:} Emergency contraception is the same as an abortion.\newline
\textbf{Status:} False.\newline
\textbf{Fact:} Emergency contraception prevents or delays ovulation and does not terminate an established pregnancy. Emergency contraceptive pills work mainly before ovulation, and a copper intrauterine device prevents pregnancy through effects on sperm and fertilization. They do not end an established intrauterine pregnancy and are not the same as abortion medication. Effectiveness depends on method and timing after unprotected intercourse. They do not protect against STIs or future intercourse in the cycle. Severe pain, fainting, or unusual bleeding after use needs evaluation. Local clinical guidance determines available options and contraindications.
\item[] \textbf{IMP-015.} \textbf{Myth:} Exercise or sexual intercourse routinely causes miscarriage in an uncomplicated pregnancy.\newline
\textbf{Status:} False.\newline
\textbf{Fact:} Appropriate activity and sex are safe for many pregnancies, but individual restrictions may apply and should be discussed with the maternity-care professional. In an uncomplicated pregnancy, ordinary physical activity and sex are not generally causes of miscarriage. Miscarriage is most often related to chromosomal or other biologic factors outside a person's control. Restrictions may be needed for bleeding, placenta problems, cervical insufficiency, preterm labor risk, ruptured membranes, or specific medical conditions. Pain, bleeding, fluid leakage, or contractions after activity require prompt obstetric advice. Exercise should be adapted to fitness, symptoms, temperature, and fall risk. People should receive individualized guidance rather than blame.
\item[] \textbf{IMP-016.} \textbf{Myth:} A person who is choking should be given water.\newline
\textbf{Status:} False.\newline
\textbf{Fact:} Do not give food or drink during active choking; use appropriate first aid and call emergency services when needed. Water cannot reliably dislodge an object blocking the airway and may worsen aspiration. If the person can cough forcefully or speak, encourage coughing and observe closely. If they cannot breathe, speak, or cough effectively, use age-appropriate choking first aid and call emergency services. Do not blindly sweep the mouth or push food down. After a severe choking episode, medical evaluation may be needed even if the object seems to clear. Training in local first-aid protocols is valuable.
\item[] \textbf{IMP-017.} \textbf{Myth:} More medical screening is always better.\newline
\textbf{Status:} False.\newline
\textbf{Fact:} Tests in people without a clear indication can produce false positives, incidental findings, anxiety, and unnecessary procedures. Screening benefits selected people when early detection improves meaningful outcomes. Outside those situations, testing can detect harmless abnormalities, trigger invasive follow-up, and identify disease that would never have caused harm. The balance depends on age, risk, test accuracy, interval, treatment options, and personal values. A symptom-driven diagnostic test is different from population screening. People should ask what decision a test will change and what harms may follow. More information is not automatically better information.
\item[] \textbf{IMP-018.} \textbf{Myth:} A normal test result proves that serious disease is impossible.\newline
\textbf{Status:} False.\newline
\textbf{Fact:} Every test has limitations; results must be interpreted in light of symptoms, timing, pre-test probability, and follow-up needs. False-negative results occur when a disease is early, intermittent, outside the test's detection range, or measured under unsuitable conditions. A test's sensitivity, specificity, predictive value, and the person's prior risk determine what a result means. Persistent or worsening symptoms may require repeat testing, another test, or specialist evaluation. A normal result is reassuring but not an absolute guarantee. Conversely, an abnormal result may need confirmation before treatment. Clinical context remains essential.
\item[] \textbf{IMP-019.} \textbf{Myth:} A person with a concussion must be kept awake.\newline
\textbf{Status:} False.\newline
\textbf{Fact:} After appropriate assessment and in the absence of danger signs, uninterrupted sleep is generally acceptable; worsening symptoms require urgent medical care \cite{ref62}. Sleep supports recovery and does not by itself cause a dangerous loss of consciousness. The old rule came from concern about bleeding inside the skull, which is a different and much more serious condition. A person with a possible concussion still needs attention to danger signs such as repeated vomiting, worsening headache, seizure, confusion, weakness, unequal pupils, or difficulty waking. A clinician's advice takes priority, especially after a significant impact, in a young child, or when blood-thinning medicine is involved. The safe message is not ``ignore sleep''; it is ``obtain appropriate assessment and monitor for deterioration.''
\item[] \textbf{IMP-020.} \textbf{Myth:} Sugar makes children hyperactive.\newline
\textbf{Status:} False.\newline
\textbf{Fact:} Controlled trials and meta-analysis do not show sugar to be a general direct cause of hyperactive behavior, although individual diet and behavior are more complex \cite{ref7}. In blinded studies, children who consumed sugar did not consistently become more active or cognitively impaired than children given a placebo. Parents' expectations and the stimulating setting of parties can create a strong impression of a ``sugar high.'' This does not mean diet is irrelevant to health, sleep, weight, or a child's overall wellbeing. Nor does it rule out every individual response or every effect of food additives. Persistent inattention or hyperactivity should be assessed clinically rather than attributed automatically to sweets.
\item[] \textbf{IMP-021.} \textbf{Myth:} Postpartum depression is rare.\newline
\textbf{Status:} False.\newline
\textbf{Fact:} Depressive symptoms after childbirth are common and deserve screening, support, and treatment when present \cite{ref63}. Postpartum depression is a common medical and mental-health complication, not a personal failure. Symptoms can include persistent sadness, loss of interest, guilt, anxiety, sleep or appetite changes, and difficulty functioning. The ``baby blues'' usually improve, whereas more severe or persistent symptoms require professional assessment. Untreated illness can affect the parent's wellbeing, bonding, and ability to care for the infant. Thoughts of self-harm or harm to the baby require immediate emergency support.
\item[] \textbf{IMP-022.} \textbf{Myth:} Suicide affects only people with a mental illness.\newline
\textbf{Status:} False.\newline
\textbf{Fact:} Suicide risk can involve mental illness, substance use, physical illness, stress, relationships, access to lethal means, and other interacting factors \cite{ref64}. Suicide is rarely caused by one circumstance or one diagnosis. Mental illness can increase risk, but relationship conflict, substance use, chronic pain, financial crisis, isolation, trauma, and access to lethal means can also contribute. The absence of a known diagnosis does not make someone safe. Asking directly about suicidal thoughts does not create suicidal intent and can open a path to help. Any immediate danger requires local emergency services or a crisis line.
\item[] \textbf{IMP-023.} \textbf{Myth:} Bleeding after menopause is normal.\newline
\textbf{Status:} False.\newline
\textbf{Fact:} Any vaginal bleeding after menopause should be medically assessed, even if it occurs once or is only spotting \cite{ref65}. Postmenopausal bleeding can have benign causes, including thinning tissues, polyps, or hormone-related changes. It can also be an early sign of endometrial or other gynecological cancer. The amount of bleeding does not reliably distinguish harmless from serious causes. A single episode of spotting should therefore not be dismissed. Prompt assessment usually involves history, examination, and testing selected by a clinician.
\item[] \textbf{IMP-024.} \textbf{Myth:} Adults do not need routine vaccinations.\newline
\textbf{Status:} False.\newline
\textbf{Fact:} Adult vaccine recommendations vary by age, prior vaccination, pregnancy, health conditions, occupation, and travel; adults should review them periodically \cite{ref66}. Immunity from childhood vaccination may wane, and some vaccines are specifically recommended later in life. Adults may need influenza, COVID-19, tetanus-diphtheria-pertussis, shingles, pneumococcal, hepatitis, HPV, or travel-related vaccines depending on circumstances. Recommendations differ by country and are updated as disease patterns and vaccine evidence change. A vaccine history and health review are more reliable than age alone. ``Routine'' does not mean every adult receives every vaccine.
\item[] \textbf{IMP-025.} \textbf{Myth:} Tanning beds are safer than outdoor sunlight.\newline
\textbf{Status:} False.\newline
\textbf{Fact:} Tanning beds expose skin to ultraviolet radiation and are not a safe alternative to sun protection \cite{ref60}. Artificial tanning devices emit ultraviolet radiation that can damage cellular DNA. A tan is a response to skin injury, not evidence of improved health. Repeated exposure increases premature skin aging and skin-cancer risk. Avoiding indoor tanning is safer than trying to control the dose by shortening sessions. Protective clothing, shade, and appropriately used broad-spectrum sunscreen reduce exposure but do not make intentional tanning risk-free.
\item[] \textbf{IMP-026.} \textbf{Myth:} Oral or anal sex is automatically safe from sexually transmitted infections.\newline
\textbf{Status:} False.\newline
\textbf{Fact:} Several infections can be transmitted through oral or anal sex; appropriate barriers, testing, vaccination, and risk-reduction advice matter \cite{ref59}. Transmission risk depends on the infection, the body site, contact with mucosa or broken skin, and use of barriers. Gonorrhea, syphilis, herpes, HPV, hepatitis, and HIV can be relevant to sexual practices involving the mouth or anus. Some infections cause no symptoms, so appearance is not a reliable safety test. Condoms, dental dams, vaccination where available, testing, and mutual communication reduce risk. Individual advice should reflect partners, exposures, and local services.
\item[] \textbf{IMP-027.} \textbf{Myth:} Blue light from ordinary screens causes permanent eye damage.\newline
\textbf{Status:} False.\newline
\textbf{Fact:} Prolonged screen use can contribute to eyestrain and dry eye, but ordinary screen exposure has not been established as a cause of permanent retinal damage \cite{ref60}. Screen use commonly causes temporary discomfort through prolonged focusing and reduced blinking. Breaks, appropriate viewing distance, adequate lighting, and treatment of dry eye can reduce symptoms. Laboratory findings from very intense light exposures cannot be directly equated with normal phone or computer use. Children who sit unusually close or have persistent visual symptoms may need an eye examination. The greater established concern is visual fatigue, not routine screen-induced blindness.
\item[] \textbf{IMP-028.} \textbf{Myth:} All heel pain is plantar fasciitis.\newline
\textbf{Status:} False.\newline
\textbf{Fact:} Heel pain has multiple possible causes, including nerve problems, stress injury, and other disorders; persistent or severe pain may need assessment \cite{ref60}. Plantar fasciitis is common, but it is only one possible explanation. Stress fractures, nerve compression, tendon disorders, inflammatory disease, infection, and referred pain can produce similar symptoms. Location, timing, swelling, injury history, and neurological findings help distinguish causes. Severe pain, inability to bear weight, fever, numbness, or trauma increases the need for prompt assessment. Treatment should address the cause rather than assume every heel symptom is plantar fasciitis.
\item[] \textbf{IMP-029.} \textbf{Myth:} Heart failure means the heart has stopped beating.\newline
\textbf{Status:} False.\newline
\textbf{Fact:} Heart failure means the heart cannot pump or fill well enough to meet the body's needs; it is different from cardiac arrest. The heart usually continues beating, although breathlessness, fatigue, swelling, and exercise intolerance may occur. Cardiac arrest is an abrupt loss of effective heart pumping and requires immediate CPR and defibrillation when available. Confusing the terms can delay emergency action or misunderstand a chronic diagnosis. Heart failure has many causes and can often be treated or managed. Sudden severe breathlessness, chest pain, fainting, or collapse requires emergency care \cite{ref46}.
\item[] \textbf{IMP-030.} \textbf{Myth:} Young people cannot develop heart disease.\newline
\textbf{Status:} False.\newline
\textbf{Fact:} Cardiovascular risk can begin early through inherited cholesterol disorders, diabetes, hypertension, smoking, obesity, inflammation, and other factors. Atherosclerosis may develop for years before symptoms appear. Age is an important risk factor, but it does not provide immunity. Chest pain, unexplained breathlessness, fainting, or palpitations deserve assessment at any age. Prevention includes healthy activity, avoiding tobacco, managing blood pressure and lipids, and addressing conditions early. Risk should be assessed individually rather than dismissed because someone is young \cite{ref46}.
\item[] \textbf{IMP-031.} \textbf{Myth:} A family history of heart disease means there is nothing you can do to prevent it.\newline
\textbf{Status:} False.\newline
\textbf{Fact:} Family history can increase risk through genes, shared environment, and shared behaviors, but it is not a fixed outcome. Blood pressure, cholesterol, glucose, smoking, activity, sleep, and diet can often be improved. Earlier risk assessment may be especially important when relatives had premature cardiovascular disease. Preventive medicines may be appropriate for selected people after clinician review. A family history should prompt planning, not fatalism. Prevention cannot reduce risk to zero, but it can substantially change it \cite{ref46}.
\item[] \textbf{IMP-032.} \textbf{Myth:} Cholesterol testing can wait until middle age.\newline
\textbf{Status:} False.\newline
\textbf{Fact:} Cholesterol disorders can begin in childhood or young adulthood, particularly with inherited conditions. Early identification helps estimate lifetime cardiovascular risk and guides lifestyle or treatment decisions. Screening timing depends on age, family history, health conditions, and local recommendations. A person can feel completely well despite high LDL cholesterol. Testing is one part of risk assessment and does not replace attention to blood pressure, glucose, smoking, and activity. People with strong family histories or known conditions may need earlier testing \cite{ref46}.
\item[] \textbf{IMP-033.} \textbf{Myth:} Leg pain is always ordinary aging and cannot be related to the heart or circulation.\newline
\textbf{Status:} False.\newline
\textbf{Fact:} Exertional calf or thigh pain that improves with rest can be a symptom of peripheral artery disease. PAD reflects reduced blood flow in the limbs and is associated with increased risk of heart attack and stroke. Muscle, joint, nerve, and spinal conditions can cause similar pain, so the symptom is not diagnostic by itself. Cold feet, non-healing wounds, color change, or pain at rest increase concern. Evaluation may include pulses, ankle-brachial pressure, and cardiovascular risk assessment. Persistent or sudden severe limb pain needs prompt care \cite{ref46}.
\item[] \textbf{IMP-034.} \textbf{Myth:} A fast heartbeat automatically means a heart attack.\newline
\textbf{Status:} False.\newline
\textbf{Fact:} Heart rate normally rises with exercise, stress, fever, dehydration, pain, caffeine, and some medicines. Palpitations can also arise from harmless extra beats or from arrhythmias that need evaluation. A heart attack is diagnosed from symptoms, examination, ECG, and cardiac biomarkers rather than pulse rate alone. Persistent rapid heart rate with chest pain, fainting, breathlessness, or weakness is urgent. A new or recurrent palpitation pattern should be discussed with a clinician. Do not self-diagnose or ignore severe associated symptoms \cite{ref46}.
\item[] \textbf{IMP-035.} \textbf{Myth:} Exercise should be avoided after a heart attack.\newline
\textbf{Status:} False.\newline
\textbf{Fact:} Carefully planned activity and cardiac rehabilitation can improve recovery, function, and future cardiovascular health for many survivors. The intensity and timing must be individualized according to infarct size, procedures, rhythm, heart function, symptoms, and clinician advice. Unsupervised strenuous exertion immediately after an event is different from medically guided rehabilitation. Walking and gradual progression are often part of recovery. Chest pain, fainting, severe breathlessness, or sustained palpitations during activity requires urgent attention. Avoiding all movement can worsen deconditioning \cite{ref46}.
\item[] \textbf{IMP-036.} \textbf{Myth:} Diabetes medication completely removes the cardiovascular risk caused by diabetes.\newline
\textbf{Status:} False.\newline
\textbf{Fact:} Glucose-lowering treatment can reduce complications, but diabetes remains a cardiovascular risk condition even when glucose readings improve. Blood pressure, LDL cholesterol, kidney disease, smoking, weight, inflammation, and duration of diabetes also influence risk. Comprehensive care may include lipid and blood-pressure treatment, activity, nutrition, smoking cessation, and medicines with cardiovascular benefit when appropriate. A normal glucose value is not a complete cardiovascular assessment. Risk reduction requires ongoing follow-up. Medication should complement rather than replace broader prevention \cite{ref46}.
\item[] \textbf{IMP-037.} \textbf{Myth:} Oil pulling reverses cavities or cures gum disease.\newline
\textbf{Status:} False.\newline
\textbf{Fact:} Swishing oil may change mouth sensation, but it has not been shown to reverse established cavities or replace periodontal treatment. Dental caries involve demineralization and bacterial acid, while gum disease involves inflammation and tissue damage. Effective prevention includes fluoride toothpaste, cleaning between teeth, dietary management, and professional care. Delaying treatment can allow decay or periodontal destruction to progress without severe early pain. Oils can also be aspirated or cause gastrointestinal problems if swallowed. Use evidence-based dental care rather than a social-media substitute \cite{ref47}.
\item[] \textbf{IMP-038.} \textbf{Myth:} Charcoal toothpaste safely whitens teeth without harming them.\newline
\textbf{Status:} False.\newline
\textbf{Fact:} Charcoal products may remove surface stains temporarily but do not reliably whiten intrinsic tooth color. Abrasive products can wear enamel and expose more sensitive dentin with repeated use. Many charcoal toothpastes lack fluoride, which is important for cavity prevention. Black color and ``natural'' marketing do not establish safety or effectiveness. Professional whitening products and dental advice are safer ways to address discoloration. Persistent color change may reflect decay, trauma, medication, or enamel disease \cite{ref47}.
\item[] \textbf{IMP-039.} \textbf{Myth:} DIY braces, rubber bands, or home veneers are safe alternatives to dental treatment.\newline
\textbf{Status:} False.\newline
\textbf{Fact:} Unsupervised tooth movement can damage roots, bone, gums, the bite, and neighboring teeth. Rubber bands can migrate into gum tissue and cause loss of supporting structures. Home-applied veneers or filing can permanently remove enamel and conceal disease. Orthodontic treatment requires examination, imaging when indicated, controlled forces, and follow-up. A cheaper or faster appearance change may create costly irreversible injury. Cosmetic or alignment concerns should be assessed by a licensed dental professional \cite{ref47}.
\item[] \textbf{IMP-040.} \textbf{Myth:} Mouth taping is a safe treatment for snoring or sleep apnea.\newline
\textbf{Status:} Uncertain.\newline
\textbf{Fact:} Evidence for routine mouth taping is limited, and it may be unsafe for people with nasal obstruction, vomiting risk, lung disease, or untreated sleep apnea. Snoring can be a symptom of obstructive sleep apnea, which requires assessment rather than simply forcing nasal breathing. Taping may worsen anxiety, breathing difficulty, or delayed response to an airway problem. Some people may use it under professional guidance for a specific indication, but social-media use is not equivalent to clinical treatment. Persistent loud snoring, witnessed pauses, morning headaches, or daytime sleepiness merits evaluation. Do not use tape to avoid assessment \cite{ref47}.
\item[] \textbf{IMP-041.} \textbf{Myth:} ``Quick and easy'' weight-loss products are reliably safe and effective.\newline
\textbf{Status:} False.\newline
\textbf{Fact:} Products promising dramatic weight loss without sustained dietary or activity change are a common health-fraud warning sign. Supplements may contain undeclared drugs, stimulants, contaminants, or doses that have not been adequately tested. Rapid weight loss can also cause dehydration, gallstones, nutrient deficiency, and loss of lean tissue. A product sold as ``natural'' is not automatically safe or effective. People should be especially cautious of testimonials, secret ingredients, and guarantees. Weight-management treatment should be individualized and medically supervised when medication or surgery is considered \cite{ref48}.
\item[] \textbf{IMP-042.} \textbf{Myth:} Unapproved home tests can reliably diagnose serious disease.\newline
\textbf{Status:} False.\newline
\textbf{Fact:} A test result is useful only when the test has been appropriately validated, authorized, collected, and interpreted. Unapproved or fraudulent tests may produce false reassurance, false alarms, or direct people toward unnecessary products. A home result does not automatically establish a diagnosis or replace examination and confirmatory testing. Serious symptoms require clinical evaluation even after a ``normal'' home result. Consumers should verify regulatory status and understand what decision the test can legitimately support. Marketing claims are not validation \cite{ref48}.
\item[] \textbf{IMP-043.} \textbf{Myth:} Unapproved supplements are safer than prescription treatment because they are natural.\newline
\textbf{Status:} False.\newline
\textbf{Fact:} Supplements can contain pharmacologically active ingredients, contaminants, or undisclosed prescription drugs. They may interact with medicines, alter bleeding or blood pressure, injure the liver or kidneys, or delay effective treatment. ``Natural'' describes origin rather than purity, dose, or safety. Products claiming to cure cancer, diabetes, arthritis, sexual problems, or rapid weight loss deserve particular caution. Do not replace prescribed treatment without discussing the decision with a clinician. Report suspected adverse effects and retain the product label for assessment \cite{ref48}.
\item[] \textbf{IMP-044.} \textbf{Myth:} Anti-aging pills can reliably stop or reverse normal aging.\newline
\textbf{Status:} False.\newline
\textbf{Fact:} Aging is a complex biological process involving many organs and pathways, not a single deficiency that one pill can correct. No supplement has been established to stop human aging or restore youth. Products promising rapid rejuvenation may be ineffective, contaminated, or unsafe. Healthy activity, sleep, vaccination, nutrition, social connection, and management of treatable risk factors have stronger general support. Some medicines treat specific diseases but should not be marketed as universal anti-aging cures. Extraordinary claims require high-quality clinical evidence, not testimonials \cite{ref48}. \emph{(See also REP-049.)}
\item[] \textbf{IMP-045.} \textbf{Myth:} Taking cholesterol medicine means diet and lifestyle changes are no longer needed.\newline
\textbf{Status:} False.\newline
\textbf{Fact:} Cholesterol-lowering medicines reduce risk for selected patients but do not replace activity, nutrition, smoking cessation, sleep, or blood-pressure control. Lifestyle measures can improve LDL cholesterol, triglycerides, fitness, glucose, and overall cardiovascular risk. Medicines and lifestyle changes often work together rather than as alternatives. The correct plan depends on baseline risk, prior cardiovascular disease, side effects, and treatment goals. People should take medicines as prescribed and discuss barriers rather than stopping them. No single intervention removes all cardiovascular risk \cite{ref49}.
\item[] \textbf{IMP-046.} \textbf{Myth:} A food label showing ``no cholesterol'' means the food is heart-healthy.\newline
\textbf{Status:} False.\newline
\textbf{Fact:} A food can contain no dietary cholesterol and still be high in saturated fat, trans fat, added sugar, sodium, or calories. Dietary cholesterol is only one part of cardiovascular nutrition. Serving size and the overall dietary pattern matter more than one front-of-package claim. Foods such as some baked or highly processed products may use ``cholesterol-free'' marketing while remaining nutritionally poor. Reading saturated fat, trans fat, sodium, fiber, and ingredient information is more informative. A label claim is not a complete health assessment \cite{ref49}.
\item[] \textbf{IMP-047.} \textbf{Myth:} Switching from butter to any margarine automatically lowers cholesterol.\newline
\textbf{Status:} False.\newline
\textbf{Fact:} Margarine products differ greatly in saturated fat, trans fat, and formulation. Some older or solid products contained partially hydrogenated oils or substantial saturated fat. Replacing butter with a soft or liquid spread lower in saturated and trans fat may improve the lipid profile, but the label must be checked. Portion size and the rest of the diet also matter. Unsaturated oils, nuts, seeds, and fish are often useful replacements for saturated fats. ``Margarine'' alone is not a guarantee of cardiovascular benefit \cite{ref49}.
\item[] \textbf{IMP-048.} \textbf{Myth:} Children do not need to worry about cholesterol.\newline
\textbf{Status:} False.\newline
\textbf{Fact:} Children can have high cholesterol, especially with familial hypercholesterolemia, diabetes, obesity, or other risk factors. Inherited severe LDL elevation can cause premature cardiovascular disease if it is not recognized. Screening recommendations vary by age, family history, and country, but childhood is not an automatic exclusion. A child can look healthy and still have an important lipid disorder. Early diagnosis allows family assessment, lifestyle support, and treatment when indicated. Parents should discuss timing and interpretation with a pediatric clinician \cite{ref49}.
\item[] \textbf{IMP-049.} \textbf{Myth:} Sexual-enhancement supplements sold as natural products are reliably safe.\newline
\textbf{Status:} False.\newline
\textbf{Fact:} Products marketed for rapid sexual enhancement may contain undeclared prescription drug ingredients or stimulants. These substances can interact with nitrates and other medicines, dangerously lowering blood pressure or stressing the heart. ``Herbal,'' ``natural,'' and online availability do not establish purity, dose, or effectiveness. The FDA has repeatedly warned about contaminated sexual-enhancement products. People with erectile or sexual-function concerns should receive an evaluation because symptoms can reflect vascular, hormonal, medication, or psychological conditions. Avoid products promising immediate or guaranteed results \cite{ref48}.
\item[] \textbf{IMP-050.} \textbf{Myth:} Bodybuilding and sports-performance supplements are safe if they are sold as dietary supplements.\newline
\textbf{Status:} False.\newline
\textbf{Fact:} Some products marketed for muscle gain or performance contain undeclared anabolic steroids, steroid-like compounds, stimulants, or other active ingredients. Potential harms include liver injury, cardiovascular events, infertility, hormonal suppression, psychiatric symptoms, and drug interactions. Manufacturing and marketing as a supplement do not guarantee that every ingredient is accurately listed or tested for effectiveness. Adolescents and young adults may be particularly vulnerable to these claims. Food, training, recovery, and qualified sports-nutrition advice are safer foundations. Suspected adverse effects should be reported and medically assessed \cite{ref48}.
\item[] \textbf{IMP-051.} \textbf{Myth:} Chewing special ``fitness gum'' safely sculpts the jawline.\newline
\textbf{Status:} False.\newline
\textbf{Fact:} Chewing can exercise jaw muscles but cannot selectively reshape facial bone or remove facial fat in a predictable way. Excessive or very hard chewing can aggravate the temporomandibular joint, muscles, teeth, or existing jaw pain. Products marketed with guaranteed sculpting claims lack reliable clinical evidence. Jaw asymmetry or enlargement can also reflect dental, skeletal, or medical conditions. People with clicking, locking, headaches, or pain should stop aggravating activity and seek dental or medical advice. Cosmetic claims should not override joint health \cite{ref50}.
\item[] \textbf{IMP-052.} \textbf{Myth:} Lemon juice is a safe natural teeth-whitening treatment.\newline
\textbf{Status:} False.\newline
\textbf{Fact:} Lemon juice is acidic and can soften or erode enamel with repeated contact. It may make teeth look temporarily different while increasing sensitivity and susceptibility to wear. Acid erosion is permanent because enamel does not grow back. Homemade mixtures can also irritate or burn gums, especially when combined with abrasives or peroxide. Professional whitening products are formulated and used with attention to concentration and dental condition. Acidic ``natural'' does not mean harmless \cite{ref50}.
\item[] \textbf{IMP-053.} \textbf{Myth:} Swishing concentrated hydrogen peroxide is a safe way to whiten teeth at home.\newline
\textbf{Status:} False.\newline
\textbf{Fact:} Concentrated or improperly used peroxide can irritate or chemically burn oral tissues and increase tooth sensitivity. It may damage enamel or be harmful if swallowed. A product intended for oral use is not interchangeable with household or industrial peroxide. Discoloration can have causes that require diagnosis rather than bleaching. Dental professionals can identify decay, gum disease, restorations, and enamel defects before recommending whitening. Do not use internet recipes in place of regulated products and professional guidance \cite{ref50}.
\item[] \textbf{IMP-054.} \textbf{Myth:} Sparkling water is completely harmless to teeth.\newline
\textbf{Status:} False.\newline
\textbf{Fact:} Plain unsweetened sparkling water is generally less damaging than sugary or acidic soft drinks, but carbonation lowers pH and frequent exposure may contribute to enamel erosion in some circumstances. Flavored products may contain additional acids or sugar. Sipping acidic drinks over long periods prolongs contact with teeth. Drinking with meals, avoiding prolonged sipping, and choosing plain water can reduce exposure. Sparkling water is not equivalent to soda, but ``completely harmless'' is too absolute. People with enamel erosion or sensitivity should discuss beverage habits with a dentist.
\item[] \textbf{IMP-055.} \textbf{Myth:} ``Natural'' whitening or filing is a safe substitute for dental treatment.\newline
\textbf{Status:} False.\newline
\textbf{Fact:} Filing teeth, applying acids, or using unregulated whitening mixtures can permanently remove enamel and injure gums. Veneers and alignment procedures require examination, diagnosis, controlled materials, and professional follow-up. A cosmetic hack may conceal decay, cracks, bite problems, or gum disease while creating irreversible damage. Cheap or rapid results do not prove safety. Licensed dental care can distinguish surface staining from disease and recommend appropriate options. Do-it-yourself dentistry should not be used to avoid professional evaluation \cite{ref50}.
\item[] \textbf{IMP-056.} \textbf{Myth:} 5G mobile networks cause COVID-19.\newline
\textbf{Status:} False.\newline
\textbf{Fact:} COVID-19 is caused by infection with SARS-CoV-2, a virus that cannot travel through radio or mobile networks. 5G signals are a form of non-ionizing electromagnetic radiation and do not create viral particles. Outbreaks occurred in places without 5G, and the virus spreads through respiratory exposure and contaminated hands or surfaces. Network technology can affect communication but not the biological cause of an infection. Prevention and treatment should follow public-health and clinical guidance. Do not damage infrastructure or delay care because of this claim \cite{ref51}.
\item[] \textbf{IMP-057.} \textbf{Myth:} Hot weather, sunlight, or high temperatures prevent COVID-19.\newline
\textbf{Status:} False.\newline
\textbf{Fact:} SARS-CoV-2 transmission has occurred in warm and sunny regions. Weather can influence behavior, ventilation, and environmental survival, but heat and sunlight do not reliably prevent infection. Deliberate heat exposure can cause burns, dehydration, or heat illness. Vaccination, ventilation, staying home when ill, hand hygiene, and appropriate masks in high-risk settings are more relevant prevention measures. Fever or heat exposure is not a treatment for COVID-19. Symptoms requiring care should be assessed regardless of outdoor temperature \cite{ref51}.
\item[] \textbf{IMP-058.} \textbf{Myth:} Eating garlic prevents COVID-19.\newline
\textbf{Status:} False.\newline
\textbf{Fact:} Garlic can be a nutritious food and may have antimicrobial effects in laboratory settings, but it does not provide proven protection against SARS-CoV-2 infection. No food selectively creates immunity to a specific respiratory virus. Excessive raw garlic can irritate the mouth, stomach, or skin and can interact with some medicines. Prevention requires measures supported by clinical and public-health evidence. A person with symptoms should not delay testing or care while trying dietary remedies. Food can support general nutrition without being a COVID-19 cure \cite{ref51}.
\item[] \textbf{IMP-059.} \textbf{Myth:} Drinking alcohol prevents or cures COVID-19.\newline
\textbf{Status:} False.\newline
\textbf{Fact:} Alcohol consumed in the body does not disinfect the respiratory tract or blood. It can impair judgment, increase injury risk, worsen sleep, and harm the liver and immune system when consumed excessively. Drinking concentrated alcohol or applying it internally can cause poisoning and tissue injury. Alcohol-based hand sanitizer is useful on hands when formulated correctly, but it is not a beverage or internal treatment. COVID-19 prevention and care should follow medical guidance. Severe illness or breathing difficulty requires urgent assessment \cite{ref51}.
\item[] \textbf{IMP-060.} \textbf{Myth:} Hydroxychloroquine or chloroquine reliably prevents or cures COVID-19.\newline
\textbf{Status:} False.\newline
\textbf{Fact:} Early laboratory findings and small studies did not establish a reliable clinical benefit. Controlled evidence and public-health guidance do not support unsupervised use for prevention or routine treatment. These medicines can cause dangerous heart-rhythm effects, drug interactions, eye toxicity, and overdose. A medicine used for malaria or autoimmune disease is not automatically appropriate for a viral infection. Taking someone else's prescription can delay effective care and harm the patient. Treatment decisions should follow current clinical guidance and individual evaluation \cite{ref51}.
\item[] \textbf{IMP-061.} \textbf{Myth:} Antibiotics prevent or cure COVID-19.\newline
\textbf{Status:} False.\newline
\textbf{Fact:} COVID-19 is caused by a virus, while antibiotics act against susceptible bacteria. Antibiotics may be prescribed only when a bacterial co-infection or another bacterial diagnosis is present. Unnecessary use can cause allergy, diarrhea, drug interactions, and antimicrobial resistance. A patient should not pressure a clinician for antibiotics because of fever, cough, or a positive viral test. Severe or worsening illness needs medical evaluation for oxygenation and complications. Antibiotics are not a substitute for vaccination, supportive care, or indicated antiviral treatment \cite{ref51}.
\item[] \textbf{IMP-062.} \textbf{Myth:} UV lamps safely disinfect a person's hands or skin.\newline
\textbf{Status:} False.\newline
\textbf{Fact:} UV-C devices can injure skin and eyes, and consumer lamps are not a safe substitute for handwashing. Exposing skin to UV can cause burns and cumulative DNA damage without reliably disinfecting all surfaces. Some devices are designed for controlled environmental or equipment use, not direct human exposure. Soap and water or an appropriate alcohol hand rub are safer for hands. Never shine a germicidal lamp on the body to prevent infection. Eye pain, burns, or visual symptoms after exposure require medical attention \cite{ref51}.
\item[] \textbf{IMP-063.} \textbf{Myth:} A thermal scanner can diagnose COVID-19.\newline
\textbf{Status:} False.\newline
\textbf{Fact:} Thermal scanners estimate skin temperature and may identify fever, but fever is neither necessary nor specific for COVID-19. People can be infected without fever, and other illnesses or environmental conditions can raise temperature. A screening device cannot identify SARS-CoV-2 or replace diagnostic testing. False reassurance can lead to transmission, while false alarms can cause unnecessary concern. Symptoms, exposure, testing, and clinical assessment determine next steps. Devices should be used only for the limited purpose for which they were validated \cite{ref51}.
\item[] \textbf{IMP-064.} \textbf{Myth:} Daily aspirin is safe and beneficial for every adult to prevent a first heart attack or stroke.\newline
\textbf{Status:} False.\newline
\textbf{Fact:} Aspirin can reduce some cardiovascular events in selected high-risk adults, but it also increases gastrointestinal and intracranial bleeding. Primary prevention decisions depend on age, cardiovascular risk, bleeding risk, medicines, and personal preferences. Starting aspirin without advice is not appropriate for everyone. Recommendations for people who have already had a heart attack or stroke are different from primary prevention. Never stop prescribed aspirin without discussing it with the clinician who recommended it. The balance of benefit and harm must be individualized \cite{ref52}.
\item[] \textbf{IMP-065.} \textbf{Myth:} Statins are usually dangerous and their side effects outweigh their cardiovascular benefits.\newline
\textbf{Status:} False.\newline
\textbf{Fact:} Statins can cause adverse effects and interactions, but serious treatment-related muscle injury and severe liver injury are rare. Randomized trials show that many symptoms attributed to statins occur similarly with placebo. For patients with an appropriate indication, reducing LDL cholesterol lowers the risk of heart attack and stroke. Symptoms should be taken seriously and evaluated rather than dismissed or assumed to prove toxicity. Dose adjustment, switching agents, or rechallenge may be considered clinically. Patients should not stop treatment without discussing risks and alternatives \cite{ref53}.
\item[] \textbf{IMP-066.} \textbf{Myth:} Mercury or thimerosal in vaccines causes autism.\newline
\textbf{Status:} False.\newline
\textbf{Fact:} Thimerosal contains ethylmercury and has been used as a preservative in some multidose vaccines. Research has not shown that thimerosal causes autism. Most routine childhood vaccines in the United States have not contained thimerosal since 2001, while some multidose influenza vaccines may still contain it. Ethylmercury is processed differently from methylmercury and is cleared more rapidly. Local vaccine formulations and recommendations vary, so ingredient questions should be discussed with a clinician. The evidence does not support avoiding vaccines because of an autism claim \cite{ref2,ref54}.
\item[] \textbf{IMP-067.} \textbf{Myth:} Properly worn medical masks cause oxygen deficiency or carbon-dioxide poisoning.\newline
\textbf{Status:} False.\newline
\textbf{Fact:} Properly worn medical masks may feel uncomfortable but do not ordinarily cause carbon dioxide intoxication or oxygen deficiency in healthy users. Mask fit, moisture, heat, underlying lung disease, and the type of mask affect comfort and tolerance. A mask is not a substitute for emergency treatment when someone is severely breathless. People with specific medical conditions should seek individualized advice rather than relying on online claims. Disposable masks should be changed when damp or soiled. The safety question depends on the person and setting, not a universal claim that masks suffocate users \cite{ref51}.
\item[] \textbf{IMP-068.} \textbf{Myth:} Vitamin D, vitamin C, or zinc supplements cure COVID-19.\newline
\textbf{Status:} False.\newline
\textbf{Fact:} Adequate micronutrients support general health, but supplements have not been established as a cure for COVID-19. Benefits of correcting a deficiency are different from treating an established viral infection with megadoses. Excess vitamin D can cause hypercalcemia, while excess zinc can cause nausea and copper deficiency. Supplements can also interact with medicines and create false reassurance. Prevention and treatment should follow current public health and clinical guidance. Severe symptoms require medical evaluation rather than escalating supplements \cite{ref51}.
\item[] \textbf{IMP-069.} \textbf{Myth:} Swimming pools or ordinary water transmit COVID-19 between swimmers.\newline
\textbf{Status:} False.\newline
\textbf{Fact:} SARS-CoV-2 does not ordinarily spread through treated pool water. Transmission risk comes mainly from close contact with an infected person, particularly in crowded or poorly ventilated settings. Water safety does not remove the need for distancing, ventilation, hygiene, and staying home when ill. Untreated water may have other infectious hazards, so this claim should not be generalized to every water exposure. People should follow local pool and public health rules. Respiratory symptoms still warrant appropriate precautions \cite{ref51}.
\item[] \textbf{IMP-070.} \textbf{Myth:} MSG universally causes headaches and other symptoms.\newline
\textbf{Status:} False.\newline
\textbf{Fact:} FDA considers monosodium glutamate generally recognized as safe when used as an ingredient. Some sensitive people report short-lived symptoms after consuming large amounts without food, but controlled studies have not consistently reproduced a universal reaction. MSG is the sodium salt of glutamic acid, which also occurs naturally in foods such as tomatoes and cheese. People with a reproducible personal reaction can limit intake without claiming that MSG harms everyone. Symptoms such as persistent headache, palpitations, or allergy signs need an appropriate evaluation. ``Chinese restaurant syndrome'' is not a reliable diagnosis for every symptom after eating \cite{ref55}.
\item[] \textbf{IMP-071.} \textbf{Myth:} Irradiated food becomes radioactive.\newline
\textbf{Status:} False.\newline
\textbf{Fact:} Approved food irradiation exposes food to controlled radiation to reduce microorganisms and extend shelf life. It does not leave the food radioactive because the process does not make the food a radiation source. Irradiation is regulated for safety and is comparable in purpose to other preservation methods such as pasteurization. It may cause small changes in some nutrients or sensory properties, but it does not make food inherently poisonous. Food still needs appropriate storage and handling after treatment. The label describes a safety process, not radioactive contamination \cite{ref56}.
\item[] \textbf{IMP-072.} \textbf{Myth:} All genetically modified foods are inherently unsafe to eat.\newline
\textbf{Status:} False.\newline
\textbf{Fact:} Safety is assessed for each crop and trait rather than inferred from the breeding technique alone. FDA evaluates whether foods from bioengineered plants meet applicable safety standards and compares them with appropriate traditionally bred counterparts. Risks can involve allergens, toxins, nutrition, or environmental effects and must be evaluated case by case. ``GMO'' is not a single chemical or health outcome. People may choose foods for ethical, environmental, or production reasons without claiming that every GMO food is medically unsafe. Evidence-based assessment is more informative than a universal label judgment \cite{ref57}.
\item[] \textbf{IMP-073.} \textbf{Myth:} Long COVID is imaginary, purely psychological, or always permanent.\newline
\textbf{Status:} False.\newline
\textbf{Fact:} Long COVID, also called post-COVID condition, involves symptoms that persist or emerge after the initial infection and can affect physical, cognitive, and mental health. Symptoms and duration vary widely, and the condition may interfere with daily function. Psychological distress can coexist with symptoms but does not make the illness imaginary. Not every person recovers on the same schedule, and ``always permanent'' is as unsupported as ``always quickly cured.'' Assessment focuses on symptoms, function, exclusion of urgent conditions, rehabilitation, and individualized support. WHO identifies continuing misinformation about long COVID as a barrier to diagnosis and care \cite{ref58}.
\item[] \textbf{IMP-074.} \textbf{Myth:} Eating eggs causes heart disease.\newline
\textbf{Status:} False.\newline
\textbf{Fact:} Eggs are nutrient-dense foods that provide protein, vitamins, minerals, and choline. For most healthy adults, moderate egg consumption does not by itself establish a high risk of heart disease. Overall dietary patterns, saturated-fat intake, metabolic health, and existing cholesterol disorders matter more than treating one food as universally harmful. People with diabetes, familial hypercholesterolemia, or established cardiovascular disease may need individualized dietary advice. The evidence does not justify banning eggs for everyone, but it also does not make unlimited intake appropriate for every person \cite{ref18,ref19}.
\item[] \textbf{IMP-075.} \textbf{Myth:} Everyone needs a daily multivitamin to stay healthy.\newline
\textbf{Status:} False.\newline
\textbf{Fact:} Multivitamins can be useful when a person has a diagnosed deficiency, restricted diet, pregnancy-related need, malabsorption, or another clinical indication. They cannot replace the fiber, protein, and other beneficial components of a varied diet. A basic product is not automatically necessary for a well-nourished adult, and high-dose products can create toxicity or interact with medicines. Smokers and former smokers should be cautious with products containing high beta-carotene. Supplement decisions should be based on need rather than the idea that more vitamins provide universal insurance \cite{ref18,ref20}.
\item[] \textbf{IMP-076.} \textbf{Myth:} Breakfast is essential for weight loss, and skipping it automatically causes weight gain.\newline
\textbf{Status:} False.\newline
\textbf{Fact:} Meal timing does not create a universal weight-loss rule. Some people feel better and control hunger with breakfast, while others can skip it without automatically overeating later. Weight change depends chiefly on sustained energy intake, food quality, activity, sleep, medications, and individual biology. Observational associations between breakfast and lower weight can reflect other healthy behaviors rather than breakfast itself causing weight loss. A nutritious breakfast can be valuable, but it is not mandatory for every adult \cite{ref18}. \emph{(See also GEN-065.)}
\item[] \textbf{IMP-077.} \textbf{Myth:} Only alcohol causes liver disease.\newline
\textbf{Status:} False.\newline
\textbf{Fact:} Alcohol is one cause of liver injury, but liver disease also results from metabolic dysfunction-associated steatotic liver disease, viral hepatitis, autoimmune disease, inherited disorders, medicines, and toxic supplements. Excess body weight, diabetes, and other metabolic risks can cause fatty liver and progressive scarring without alcohol use. The cause must be established rather than inferred from drinking history alone. Avoiding alcohol is important for many patients, but it does not eliminate every liver risk. Evaluation may include history, examination, blood tests, and imaging \cite{ref21}.
\item[] \textbf{IMP-078.} \textbf{Myth:} If you have no symptoms, you cannot have liver disease.\newline
\textbf{Status:} False.\newline
\textbf{Fact:} Liver disease can remain clinically silent while fat, inflammation, or scarring develops. Symptoms such as jaundice, abdominal swelling, confusion, or severe fatigue may appear only after substantial injury. Normal feelings of well-being do not substitute for testing when risk factors are present. People with obesity, diabetes, hypertension, abnormal cholesterol, viralhepatitis risk, or significant alcohol exposure may need assessment. Screening and follow-up should be guided by a healthcare professional rather than symptoms alone \cite{ref21}.
\item[] \textbf{IMP-079.} \textbf{Myth:} Liver damage is always permanent and cannot improve.\newline
\textbf{Status:} False.\newline
\textbf{Fact:} The liver can recover when the cause of injury is identified and treated, especially before advanced scarring develops. Weight loss can improve metabolic fatty liver disease, stopping alcohol can improve alcohol-related injury, and antiviral treatment can control or cure some viral hepatitis. Some established cirrhosis and end-stage damage cannot be fully reversed, so ``the liver always regenerates'' would also be false. Prognosis depends on the cause, stage, ongoing injury, and treatment. Medical assessment is important because waiting for symptoms can reduce the opportunity for recovery \cite{ref21}.
\item[] \textbf{IMP-080.} \textbf{Myth:} All hepatitis is caused by a virus.\newline
\textbf{Status:} False.\newline
\textbf{Fact:} Hepatitis means inflammation of the liver and describes a process rather than one specific infection. Hepatitis viruses A, B, C, D, and E are important causes, but alcohol, metabolic fatty liver disease, medicines, toxins, and autoimmune disease can also inflame the liver. Different causes require different tests, prevention strategies, and treatments. A positive liverinflammation test does not identify the cause by itself. Clinical history and targeted evaluation are needed before treatment is selected \cite{ref21}.
\item[] \textbf{IMP-081.} \textbf{Myth:} A daily glass of red wine is good for the heart and is harmless.\newline
\textbf{Status:} False.\newline
\textbf{Fact:} Earlier observational studies helped create the impression that red wine directly protects the heart, but correlation does not establish that alcohol caused better outcomes. Alcohol can increase risks including cancer, hypertension, liver disease, dependence, injury, and cardiovascular problems. No one should start drinking for supposed health benefits. People who drink should discuss risk reduction and follow current low-risk guidance, recognizing that lower intake generally means lower alcohol-related risk. ``Natural,'' moderate, or traditional does not make alcohol risk-free \cite{ref22}.
\item[] \textbf{IMP-082.} \textbf{Myth:} Persistent pain is only in a person's head.\newline
\textbf{Status:} False.\newline
\textbf{Fact:} Pain is a real sensory and emotional experience produced by interactions among tissues, nerves, the spinal cord, brain, mood, sleep, and context. Psychological factors can amplify or reduce pain, but that does not make the symptom imaginary or blameworthy. Some pain has clear tissue injury, while chronic pain can persist after healing and involve nervous-system sensitization. A normal scan does not prove that pain is fabricated. Assessment should consider urgent causes, function, mental health, medicines, and physical or psychological therapies together \cite{ref23}.
\item[] \textbf{IMP-083.} \textbf{Myth:} Pain medication will fix the underlying problem by itself.\newline
\textbf{Status:} False.\newline
\textbf{Fact:} Analgesics can reduce suffering and support function, but they do not automatically diagnose or remove the cause of pain. Different causes require different approaches, such as rehabilitation, treatment of inflammation or infection, surgery, psychological support, or management of a chronic condition. Opioids and other medicines also have dose-related risks, interactions, and potential for dependence. The safest plan uses the smallest effective treatment for the appropriate duration while reassessing progress. New severe, unexplained, or neurologic pain needs medical evaluation rather than repeated self-medication \cite{ref23}.
\item[] \textbf{IMP-084.} \textbf{Myth:} Doctors will not work as hard to save a patient who is a registered organ donor.\newline
\textbf{Status:} False.\newline
\textbf{Fact:} The treating team's duty is to provide appropriate care and try to save the patient's life. Organ donation is considered only after death is determined according to accepted medical and legal standards. The clinicians treating a critically ill patient are separate from transplant allocation decisions. Registering as a donor does not lower the standard of emergency or intensive care. Donation systems use safeguards to protect patient care and public trust \cite{ref24}.
\item[] \textbf{IMP-085.} \textbf{Myth:} Older adults or people with medical conditions cannot donate organs.\newline
\textbf{Status:} False.\newline
\textbf{Fact:} There is no universal age cutoff for deceased organ donation. Transplant teams assess the condition and suitability of individual organs and tissues at the time of donation. A disease may prevent use of one organ without preventing use of another. Eligibility cannot be reliably decided by a person's age, diagnosis, or appearance in advance. People who wish to donate should register and allow medical professionals to make the final determination \cite{ref24}.
\item[] \textbf{IMP-086.} \textbf{Myth:} A donor's family must pay the medical costs of organ donation.\newline
\textbf{Status:} False.\newline
\textbf{Fact:} Families are not charged for the donation process itself. Medical care provided to attempt to save a patient's life is separate from the costs associated with recovering organs or tissues for transplantation. Billing and legal arrangements can vary by healthcare system, so families should ask the transplant organization or hospital for local details. Financial myths should not prevent people from discussing donation wishes. Clear registration and family communication can reduce uncertainty during a difficult time \cite{ref24}.
\item[] \textbf{IMP-087.} \textbf{Myth:} Duct tape is a proven and reliable cure for warts.\newline
\textbf{Status:} Uncertain.\newline
\textbf{Fact:} Duct-tape occlusion has been studied as a low-cost wart treatment, but trials have produced conflicting results. Some studies suggest benefit in selected situations, while others have not shown superiority to placebo or other care. The American Academy of Dermatology describes duct tape as a possible at-home option but recommends evidence-based treatment and medical review for stubborn warts. Skin irritation, bleeding, or damage can occur if treatment is misused. A possible adjunct should not be presented as a guaranteed cure or substitute for diagnosis \cite{ref25,ref26}.
\item[] \textbf{IMP-088.} \textbf{Myth:} Coffee, greasy food, more alcohol, or a supplement can reliably cure a hangover.\newline
\textbf{Status:} False.\newline
\textbf{Fact:} A hangover reflects the effects of alcohol and recovery takes time; there is no rapid, universally reliable cure. Water, rest, and bland food may ease selected symptoms but do not remove alcohol-related impairment or prevent all complications. More alcohol can prolong harm and may reinforce dependence. Pain medicines also require caution because alcohol can increase stomach, kidney, or liver risks depending on the drug. Persistent vomiting, confusion, breathing problems, seizures, or inability to wake someone require urgent medical help \cite{ref27}.
\item[] \textbf{IMP-089.} \textbf{Myth:} All artificial sweeteners are poisonous and cause cancer.\newline
\textbf{Status:} False.\newline
\textbf{Fact:} Artificial sweeteners are a diverse group of substances, so safety cannot be inferred from the category name alone. Regulatory agencies evaluate individual sweeteners, exposure levels, and acceptable daily intakes. Current FDA information reports no safety concern for aspartame when used under approved conditions, while cancer evidence remains subject to ongoing study. ``Possibly carcinogenic'' hazard classification is not the same as proof that normal exposure causes cancer. People may limit sweeteners for taste, tolerance, or dietary reasons without claiming that every approved sweetener is poisonous \cite{ref28,ref29}.
\item[] \textbf{IMP-090.} \textbf{Myth:} Talcum powder definitely causes cancer in everyone who uses it.\newline
\textbf{Status:} Uncertain.\newline
\textbf{Fact:} Talc and asbestos are different minerals, but talc deposits can occur near asbestos and contamination is an important safety concern. Studies of cosmetic talc and cancer have produced mixed findings, particularly for genital-area use and ovarian cancer. Current authorities do not support the absolute claim that every user will develop cancer. The possible risk depends on product composition, contamination, route of exposure, duration, and the cancer outcome studied. Avoiding inhalation of powder and choosing products with appropriate safety testing are reasonable precautions while evidence continues to be evaluated \cite{ref30,ref31,ref32}.
\item[] \textbf{IMP-091.} \textbf{Myth:} Hormonal contraception inevitably causes depression or is unsafe for anyone with a history of depression.\newline
\textbf{Status:} Uncertain.\newline
\textbf{Fact:} Mental-health responses to hormonal contraception are not uniform across methods or individuals. Current evidence supports the safety of hormonal contraception for many people with depressive disorders, while some users may notice mood changes and some methods may slightly alter depression risk. A personal temporal relationship deserves attention but does not prove that every hormonal method caused depression. Stopping contraception without a replacement plan can create an unintended pregnancy risk. Method choice should include mental-health history, medical eligibility, preferences, and follow-up if symptoms change \cite{ref30,ref33}.
\item[] \textbf{IMP-092.} \textbf{Myth:} Suicide is contagious like an infectious disease.\newline
\textbf{Status:} Uncertain.\newline
\textbf{Fact:} Suicide is not transmitted biologically like a virus or bacterial infection. However, exposure to suicide or suicidal behavior can increase risk among some vulnerable people through social influence, imitation, grief, and sensationalized reporting. The effect is not universal and suicide usually reflects multiple interacting risk and protective factors. Responsible reporting avoids details that may increase imitation and emphasizes help-seeking and recovery. The accurate term is social contagion or clustering, not ordinary infectious contagion \cite{ref30,ref34}.
\item[] \textbf{IMP-093.} \textbf{Myth:} Testosterone supplements reliably increase libido, energy, and health in every man.\newline
\textbf{Status:} False.\newline
\textbf{Fact:} Testosterone therapy is intended for men with compatible symptoms and repeatedly confirmed low testosterone, not as a universal wellness supplement. Benefits depend on the cause and severity of deficiency, and libido or energy can have many other causes. Treatment can suppress sperm production and requires assessment of contraindications, monitoring, and discussion of uncertain long-term risks. Over-the-counter ``testosterone boosters'' are not equivalent to prescribed therapy and may contain undeclared or unsafe ingredients. Diagnosis should precede treatment, and routine population-wide screening or supplementation is not recommended \cite{ref30,ref35}. \end{enumerate}
\chapter{Indian Myths}
This chapter focuses on medical beliefs particularly associated with parts of Indian society. These beliefs are not held by all Indians or Hindus; many are regional, caste- or community-specific, family traditions, or practices that predate or extend beyond Hinduism. Cultural significance and medical evidence are different questions, and respectful discussion should distinguish the two.

\newcommand{\indianentry}[5]{%
\item[] \textbf{#1.} \textbf{Myth:} #2\newline
\textbf{Status:} #3.\newline
\textbf{Fact:} #4 \cite{#5}
}
\begin{enumerate}[leftmargin=0pt,labelwidth=0pt,labelsep=0pt,itemindent=0pt]
\indianentry{IND-001}{Menstruating women are physically ``impure'' and can contaminate food.}{False}{Menstruation is a normal biological process, not medical uncleanliness. Menstrual status does not make food poisonous or contaminated; food safety depends on hygiene, water, storage, temperature, and microorganisms. Stigma can instead reduce access to washing facilities, pain relief, education, and care.}{ref67}
\indianentry{IND-002}{A menstruating woman touching pickles makes them spoil.}{False}{Menstruation has no effect on pickle preservation. Spoilage depends on microorganisms, moisture, salt, acidity, oxygen, temperature, and container hygiene. Cultural rules may be observed voluntarily, but they are not microbiology.}{ref67}
\indianentry{IND-003}{Menstruating women should not enter the kitchen because their presence affects food.}{False}{There is no biological mechanism by which menstruation changes food or cooking. The relevant risks are ordinary burns and food contamination, not menstrual status. Kitchen access should be based on safety and personal choice.}{ref67}
\indianentry{IND-004}{Menstruating girls should not wash their hair or bathe.}{False}{Bathing and washing hair during menstruation are safe and may support comfort and hygiene. Menstruation does not make ordinary bathing dangerous. Restrictions that prevent washing can increase discomfort and stigma.}{ref67}
\indianentry{IND-005}{Menstrual blood is ``dirty blood'' or toxins being expelled from the body.}{False}{Menstrual fluid contains blood, cervical mucus, vaginal secretions, and shed endometrial tissue. It is not a monthly detoxification process. Very heavy, prolonged, or unusually painful bleeding deserves medical assessment.}{ref67}
\indianentry{IND-006}{Girls should avoid curd, tamarind, lemon, or other sour foods during menstruation.}{False}{Sour foods do not inherently stop or disturb menstruation. Individual intolerance or reflux may justify avoiding a particular food, but there is no universal menstrual restriction. A balanced diet and adequate fluids remain appropriate.}{ref67}
\indianentry{IND-007}{Eating papaya during pregnancy necessarily causes miscarriage.}{False}{Normal dietary consumption of ripe papaya has not been shown to cause miscarriage. Unripe or semi-ripe papaya is a different exposure and is sometimes avoided as a precaution. Bleeding or pain in pregnancy requires clinical assessment, not a fruit-based remedy.}{ref13}
\indianentry{IND-008}{Eating saffron during pregnancy makes the baby fair.}{False}{Fetal skin pigmentation is determined mainly by inherited biology and melanin pathways, not by the color of food or spices. Saffron cannot make a baby fair. Large supplemental doses should not be self-prescribed during pregnancy.}{ref13}
\indianentry{IND-009}{Drinking coconut water during pregnancy makes the baby fair.}{False}{Coconut water does not determine fetal skin color. It can be a beverage but does not replace prenatal care or balanced nutrition. People with kidney disease or fluid or potassium restrictions need individualized advice.}{ref13}
\indianentry{IND-010}{Eating dark-colored foods during pregnancy can make the baby dark.}{False}{The color of maternal food does not transfer to fetal skin. Pigmentation depends on genetics and developmental biology. This claim can reinforce colorism without providing any health benefit.}{ref13}
\indianentry{IND-011}{Eating more ghee late in pregnancy lubricates the birth canal and produces an easier delivery.}{False}{Eating ghee does not lubricate the uterus, cervix, or birth canal. Labor depends on cervical change, contractions, fetal position, pelvic anatomy, and other factors. Extra ghee adds energy but cannot prevent obstructed labor or replace skilled maternity care.}{ref13}
\indianentry{IND-012}{Pregnant women must remain indoors during an eclipse to prevent birth defects.}{False}{Solar or lunar eclipses do not cause congenital abnormalities. Birth defects have genetic, developmental, infectious, nutritional, medication-related, and other causes. Pregnant people should follow ordinary safety advice, including safe solar viewing during a solar eclipse.}{ref13}
\indianentry{IND-013}{Using a knife or scissors during an eclipse can cause a cleft lip or deformity in the baby.}{False}{There is no biological connection between cutting an object during an eclipse and fetal facial development. Cleft lip and other congenital conditions arise through complex developmental pathways. Ordinary household activities do not imprint injuries on a fetus.}{ref13}
\indianentry{IND-014}{Cutting vegetables during an eclipse can give the baby corresponding cuts or deformities.}{False}{A cut vegetable cannot transmit a cut to an unborn baby. Fetal development is not controlled by visual resemblance, thought, or activity during an eclipse. Prenatal care and nutrition are medically relevant; ritual explanations cannot diagnose congenital conditions.}{ref13}
\indianentry{IND-015}{A high, low, round, or pointed abdomen predicts the baby's sex.}{False}{Abdominal shape varies with body build, muscle tone, pregnancy number, fetal position, and gestational age. It cannot reliably identify fetal sex. Visual guesses should not replace appropriate clinical testing or prenatal care.}{ref13}
\indianentry{IND-016}{Sweet versus salty pregnancy cravings indicate whether the fetus is a boy or girl.}{False}{Cravings do not reliably identify fetal sex. They may reflect appetite, habit, nausea patterns, or changing food preferences. Persistent cravings for non-food substances should be discussed with a clinician because they may signal pica or anemia.}{ref13}
\indianentry{IND-017}{A fetal heart rate above or below a particular number tells whether the baby is a boy or girl.}{False}{Fetal heart rate changes with gestational age, activity, sleep, medicines, and temporary physiology. A single number cannot determine fetal sex. Monitoring should be interpreted by maternity professionals as an assessment of wellbeing, not as a sex test.}{ref13}
\indianentry{IND-018}{Colostrum is old, dirty, or harmful milk and should be thrown away.}{False}{Colostrum is early breast milk produced after birth and contains nutrients and immune components suited to a newborn. Its yellow color does not mean it is spoiled. Discarding it can deprive a newborn of early feeding and protection.}{ref68}
\indianentry{IND-019}{A newborn should receive honey before breast milk.}{False}{Honey should not be given to infants younger than 12 months because it can contain spores that cause infant botulism. Breast milk or an appropriate medical substitute is safer. Honey provides no special cleansing or immunity benefit.}{ref69}
\indianentry{IND-020}{A newborn should receive ghutti as a traditional first feed.}{False}{Routine prelacteal feeds are unnecessary and can interfere with early breastfeeding. Homemade or commercial mixtures may be contaminated, incorrectly concentrated, allergenic, or pharmacologically active. Newborn feeding should follow maternity and pediatric guidance.}{ref68}
\indianentry{IND-021}{Janam-ghutti cleans a newborn's stomach or intestines.}{False}{A healthy newborn gastrointestinal tract does not require ritual cleansing. The intestine is adapted to receive milk, and colostrum supports normal feeding. Unregulated mixtures can introduce microbes, sugar, alcohol, medicines, or toxic substances.}{ref68}
\indianentry{IND-022}{Putting kajal or surma in a baby's eyes improves eyesight.}{False}{Kajal and surma do not improve visual development. Some preparations contain lead or other contaminants, and application can injure or infect the eye. Redness, discharge, swelling, or poor visual response needs pediatric or ophthalmic review.}{ref79}
\indianentry{IND-023}{Kajal protects babies from nazar and therefore prevents illness.}{False}{Kajal may have cultural or cosmetic meaning, but it has no demonstrated disease-preventive effect. It cannot prevent infection or developmental conditions. Contaminated products can create lead-exposure or eye-infection risks.}{ref79}
\indianentry{IND-024}{Applying oil into a newborn's nose and ears is necessary for health.}{False}{Routine oil instillation is not needed to keep a newborn's nose or ears healthy. Oil can irritate tissue, and inserting objects can injure the ear canal or airway. Discharge, fever, or breathing difficulty should not be treated with household oil.}{ref70}
\indianentry{IND-025}{Forceful traditional infant massage makes bones stronger or changes body shape.}{False}{Gentle touch may be soothing, but forceful manipulation does not reshape bones or reliably increase bone strength. Infants have vulnerable joints, nerves, skin, and developing bones. Excessive pressure can cause bruising, fractures, dislocation, or neurologic injury.}{ref70}
\indianentry{IND-026}{Pressing a newborn girl's swollen breasts removes ``witch's milk.''}{False}{Temporary breast enlargement can occur because of maternal hormones and usually resolves without squeezing. Pressing can damage tissue and introduce infection. Redness, warmth, fever, or pus needs prompt pediatric assessment.}{ref70}
\indianentry{IND-027}{Turmeric, ash, ghee, or oil applied to the umbilical stump helps it heal.}{False}{Clean, dry cord care is generally safer than applying household substances. Powders, ash, oil, ghee, or pastes can carry microorganisms and delay recognition of infection. Spreading redness, pus, foul odor, fever, or poor feeding requires urgent care.}{ref70}
\indianentry{IND-028}{A baby's head must be repeatedly shaved to make the hair grow thicker.}{False}{Shaving cuts the visible hair shaft but does not increase the number of follicles. Hair may look darker or coarser while growing back because the ends are blunt. Shaving also carries avoidable risks of skin cuts and infection.}{ref70}
\indianentry{IND-029}{Mundan permanently improves hair density.}{False}{Ritual shaving cannot create new follicles or permanently increase hair density. Apparent thickness after shaving is a visual effect caused by uniformly short shafts. Persistent hair loss or scalp inflammation should be evaluated medically.}{ref70}
\indianentry{IND-030}{A black thread, black dot, or amulet medically protects a child from disease caused by nazar.}{False}{These objects may have cultural meaning but have no demonstrated ability to prevent infection or other disease. They should never replace vaccination, safe water, nutrition, or medical care. Tight threads can injure skin and small objects can become choking hazards.}{ref70}
\indianentry{IND-031}{Illness following nazar lagna can be diagnosed from chillies, lemon, or salt.}{False}{The appearance or behavior of such objects cannot identify the cause of an illness. Symptoms require history, examination, and appropriate testing. Rituals may provide comfort but cannot distinguish infection, poisoning, dehydration, or an emergency.}{ref30}
\indianentry{IND-032}{Burning chillies can determine whether someone has the evil eye.}{False}{Smoke and smell are affected by moisture, oils, airflow, and the material itself. They cannot diagnose disease or supernatural causation. Burning material can irritate the eyes and lungs, especially in children and people with asthma.}{ref30}
\indianentry{IND-033}{Epileptic seizures are caused by possession, ghosts, or black magic.}{False}{Epileptic seizures result from abnormal electrical activity in the brain, although other conditions can look similar. A first seizure requires medical evaluation. During a seizure, protect the person from injury, do not restrain them, and never put objects in the mouth.}{ref71}
\indianentry{IND-034}{Making someone having a seizure smell shoes can revive them.}{False}{Smelling shoes or pungent substances does not stop seizure activity or restore consciousness. It may expose the person to irritants and distract bystanders from preventing injury. Time the seizure and seek emergency help when it is prolonged, repeated, or a first event.}{ref71}
\indianentry{IND-035}{Giving an iron object or key to a person during a seizure stops it.}{False}{Holding metal does not alter the brain's electrical activity. It can cause cuts, choking, or injury if forced into the hand or mouth. Clear hazards, cushion the head, and place the person on their side after convulsions allow it.}{ref71}
\indianentry{IND-036}{Tantric or exorcism rituals can cure epilepsy.}{False}{Rituals do not treat the abnormal electrical activity that causes epilepsy and should not replace neurologic assessment or prescribed medicine. Spiritual support may coexist with care if it is safe and voluntary. Restraint, injury, starvation, or treatment cessation is dangerous.}{ref71}
\indianentry{IND-037}{Mental illness results from possession, upari hawa, black magic, or ghosts.}{False}{Mental disorders arise through interacting biological, psychological, developmental, and social factors. Cultural beliefs may influence how distress is described, but possession is not a medical explanation for all psychiatric symptoms. Stigma and coercive rituals can worsen risk and delay treatment.}{ref30}
\indianentry{IND-038}{A jhaad-phoonk ritual can medically cure psychiatric disease.}{False}{Ritual practices are not substitutes for evidence-based psychiatric treatment. Depending on the condition, effective care may include psychotherapy, medication, social support, or emergency protection. Suicidal thoughts, severe confusion, or dangerous behavior require urgent help.}{ref30}
\indianentry{IND-039}{Cow urine cures cancer.}{False}{No reliable clinical evidence establishes cow urine as a cancer cure. Cancer requires diagnosis and treatment tailored to its type and stage. Delaying oncology care for an unproven remedy can be dangerous. Natural or traditional origin does not establish effectiveness, purity, or safety.}{ref48}
\indianentry{IND-040}{Cow urine cures diabetes.}{False}{Cow urine is not an established treatment for diabetes and cannot replace glucose monitoring, nutrition care, or prescribed medicines. Uncontrolled diabetes can damage the eyes, kidneys, nerves, heart, and blood vessels. Stopping treatment can cause severe hyperglycemia or ketoacidosis.}{ref48}
\indianentry{IND-041}{Cow urine can replace medicines because it is a natural disinfectant or medicine.}{False}{Natural origin does not establish dose, purity, efficacy, or safety. A product used outside a validated indication may contain contaminants, interact with medicines, or delay effective care. Disinfection of surfaces is not the same as treatment of an infection inside the body.}{ref48}
\indianentry{IND-042}{Cow dung applied to wounds promotes healing and prevents infection.}{False}{Cow dung can contaminate an open wound with bacteria, parasites, and spores, including organisms capable of causing tetanus. It does not provide sterile wound care. Rinse wounds with clean running water, control bleeding, and obtain medical advice for deep or contaminated wounds.}{ref48}
\indianentry{IND-043}{Cow dung applied to a newborn's umbilical stump helps healing.}{False}{Applying dung to a newborn's cord can introduce infection, including tetanus. Clean, dry cord care and skilled newborn care are safer. Redness, swelling, discharge, foul odor, fever, poor feeding, or lethargy requires urgent evaluation.}{ref48}
\indianentry{IND-044}{Ganga water cannot contain disease-causing organisms because it is sacred or self-purifying.}{False}{Religious significance does not guarantee microbiological safety. Untreated surface water can contain bacteria, viruses, parasites, and chemical contaminants. Clear or flowing water is not proof that it is safe, especially for infants and people with weakened immunity.}{ref48}
\indianentry{IND-045}{Panchagavya can cure major diseases.}{False}{Broad claims that panchagavya cures major diseases are not supported by robust clinical evidence. Preparations can vary in composition and may be contaminated. Serious illness requires diagnosis and therapies shown to improve outcomes.}{ref80}
\indianentry{IND-046}{Tulsi can prevent or cure virtually every infection.}{False}{Tulsi contains biologically active compounds, but laboratory activity is not the same as proven treatment in patients. Different infections require different diagnoses, medicines, doses, and durations. Tulsi cannot replace antibiotics, antivirals, vaccination, or emergency care when indicated.}{ref80}
\indianentry{IND-047}{Haldi doodh can cure bacterial or viral infections.}{False}{Turmeric milk may be consumed as food and can be comforting, but it is not a universal antimicrobial treatment. It cannot replace indicated antibiotics, antivirals, fluids, or monitoring. Concentrated products may interact with medicines or be unsafe in pregnancy.}{ref80}
\indianentry{IND-048}{Neem can cure virtually any skin disease.}{False}{Skin disorders have different causes, including infection, allergy, inflammation, autoimmune disease, and cancer. Neem is not a universal treatment and household preparations are not standardized or sterile. A spreading rash, painful lesion, pus, fever, or non-healing sore should be assessed.}{ref80}
\indianentry{IND-049}{Karela juice can cure diabetes and eliminate the need for medication.}{False}{Evidence does not establish karela as a replacement for proven diabetes treatment. Juice preparations vary in dose and can contribute to low blood glucose with medicines. Stopping prescribed therapy can produce dangerous hyperglycemia.}{ref81}
\indianentry{IND-050}{Jamun seeds can cure diabetes permanently.}{False}{There is insufficient clinical evidence that jamun seeds permanently cure diabetes. Diabetes is managed through individualized nutrition, activity, monitoring, and medicines when indicated. A lower glucose reading does not prove that the disease has been eliminated.}{ref81}
\indianentry{IND-051}{Jaggery is sugar-free and safe in unlimited amounts for diabetes.}{False}{Jaggery contains substantial sugars and can raise blood glucose. Traditional production does not make it carbohydrate-free or harmless in unlimited amounts. The relevant issue is the total carbohydrate pattern and portion.}{ref49}
\indianentry{IND-052}{Desi ghee can be eaten without limit because it is inherently heart-protective.}{False}{Ghee is energy-dense and contains substantial saturated fat. Health effects depend on quantity, replacement foods, total dietary pattern, and individual risk factors. Traditional preparation does not make unlimited intake heart-protective.}{ref49}
\indianentry{IND-053}{Rock salt or sendha namak is harmless for people with hypertension.}{False}{Rock salt still contains substantial sodium. Substituting one salt for another without reducing total sodium may not lower blood-pressure risk. People with hypertension, heart failure, or kidney disease should follow individualized dietary advice.}{ref49}
\indianentry{IND-054}{Black salt is essentially sodium-free.}{False}{Black salt contains sodium and is not a free-use substitute for people restricting sodium. Its flavor may allow some people to use less, but that depends on the amount consumed. Natural or mineral-rich does not mean sodium-free.}{ref49}
\indianentry{IND-055}{Water stored in a copper vessel cures numerous diseases.}{False}{Copper is an essential trace element, but copper-vessel water is not a general medicine. Evidence for broad treatment claims is inadequate, and excess copper exposure can cause harm. Water safety also depends on source contamination and storage hygiene.}{ref80}
\indianentry{IND-056}{Water kept overnight in copper cures digestive and metabolic disease.}{False}{Overnight storage does not transform water into a cure. Copper exposure depends on the vessel, water chemistry, contact time, and cleaning, while excessive exposure can be harmful. Ongoing symptoms need evaluation for their actual cause.}{ref80}
\indianentry{IND-057}{Repeatedly drinking hot water “melts” body fat.}{False}{Hot water does not physically melt stored adipose tissue. Body fat changes through sustained energy balance, influenced by biology and environment. Very hot drinks can burn the mouth and esophagus.}{ref5}
\indianentry{IND-058}{Ajwain water burns abdominal fat.}{False}{Ajwain water does not selectively metabolize abdominal adipose tissue. No drink can target fat loss from one body region. Weight management should use sustainable nutrition, activity, sleep, and medical support when needed.}{ref5}
\indianentry{IND-059}{Lemon-honey water detoxifies the liver.}{False}{The liver performs its own biochemical processing; lemon and honey do not wash toxins out of it. Honey is not suitable for infants under 12 months and adds sugar. Jaundice or other liver symptoms require medical assessment.}{ref5}
\indianentry{IND-060}{Eating curd at night produces harmful mucus in everyone.}{False}{There is no general physiological rule that curd causes pathological mucus at night. Some people may have reflux, lactose-related symptoms, or intolerance, which is different from a universal mucus effect. Persistent cough or breathing symptoms need evaluation.}{ref5}
\indianentry{IND-061}{Milk and fish together produce vitiligo.}{False}{This food combination does not cause vitiligo. Vitiligo involves loss or dysfunction of pigment-producing melanocytes and is commonly associated with autoimmune mechanisms. Unnecessary avoidance can impair nutrition and reinforce stigma.}{ref82}
\indianentry{IND-062}{Milk followed by sour food causes vitiligo.}{False}{Eating sour food after milk does not cause vitiligo. Vitiligo is not an incompatibility reaction between foods and is not caused by food colors mixing in the body. New or changing depigmented patches should be assessed.}{ref82}
\indianentry{IND-063}{Vitiligo is contagious.}{False}{Vitiligo is not an infection and does not spread by touch, food, or shared clothing. It can have autoimmune and genetic contributors. Stigma may cause isolation or delayed care, while dermatologic assessment can confirm the diagnosis.}{ref82}
\indianentry{IND-064}{Leprosy results from sins committed in a previous life.}{False}{Leprosy is an infectious disease caused by Mycobacterium leprae, not a punishment or moral failure. Early diagnosis and multidrug therapy can cure the infection and reduce disability. Stigma and delayed treatment harm patients and families.}{ref74}
\indianentry{IND-065}{Leprosy spreads simply by touching someone once.}{False}{Casual contact or a single touch is not the usual route of transmission. Prolonged, close exposure to an untreated person is more relevant, and effective treatment reduces infectiousness. Suspicious skin patches, numbness, or nerve symptoms should prompt evaluation without blame.}{ref74}
\indianentry{IND-066}{Snakebite can be cured by a mantra.}{False}{Chanting does not neutralize venom. Suspected snakebite requires rapid transport, limb immobilization, and assessment for antivenom and supportive care. Do not cut, suck, burn, or apply chemicals to the bite.}{ref72}
\indianentry{IND-067}{A snake-stone or nagmani can pull venom from a snakebite.}{False}{A stone cannot extract venom that has entered tissue and the bloodstream. Attaching or cutting around it can injure the skin and waste time. Keep the person still and arrange urgent transport for professional assessment.}{ref72}
\indianentry{IND-068}{An ojha, bhopa, or tantrik can remove snake poison before hospital treatment.}{False}{Ritual intervention cannot remove systemic venom. Delay increases the risk of paralysis, bleeding, tissue injury, kidney failure, and death. The victim should be transported promptly without unnecessary walking or a tight tourniquet.}{ref72}
\indianentry{IND-069}{Turmeric, lime, or herbal paste neutralizes snake venom.}{False}{These substances do not neutralize venom circulating through the body. They can irritate the wound, introduce infection, and delay antivenom evaluation. Keep the person calm and still and obtain emergency medical care.}{ref72}
\indianentry{IND-070}{A dog bite can be treated with chilli powder, turmeric, oil, or another traditional application.}{False}{Household substances do not prevent rabies or reliably disinfect a bite. Wash the wound immediately with soap and running water and obtain urgent assessment for rabies prophylaxis, tetanus, antibiotics, and wound care. Rabies is preventable after exposure but almost always fatal after symptoms begin.}{ref73}
\indianentry{IND-071}{A person bitten by a dog should avoid particular foods while receiving rabies vaccination.}{False}{Standard rabies vaccination does not require traditional food restrictions. A normal balanced diet is usually appropriate unless a clinician identifies another issue. The critical steps are wound washing, timely vaccination, and rabies immunoglobulin when indicated.}{ref73}
\indianentry{IND-072}{The old ``14 injections in the stomach'' are still required after dog bites.}{False}{Modern rabies post-exposure regimens do not use the historical abdominal injection schedule. The schedule and route depend on the vaccine, exposure, prior vaccination, immune status, and national guidance. Some exposures also require rabies immunoglobulin.}{ref73}
\indianentry{IND-073}{A child with measles should not be bathed because the rash will go back inside the body.}{False}{Bathing does not force a measles rash inside the body. Gentle washing can support comfort and hygiene if the child is stable and kept warm. Measles can cause pneumonia, dehydration, encephalitis, and other complications, so medical and public-health advice matters.}{ref75}
\indianentry{IND-074}{Measles is a visitation of Sheetala Mata and should not be medically treated.}{False}{Measles is caused by measles virus and can produce serious complications. Cultural interpretation may be meaningful, but it does not replace diagnosis, supportive care, vitamin A when recommended, or treatment of complications. Breathing difficulty, dehydration, altered consciousness, or persistent fever requires urgent care.}{ref75}
\indianentry{IND-075}{Chickenpox is simply a manifestation of a goddess, so medicines should always be avoided.}{False}{Chickenpox is caused by varicella-zoster virus. Many cases are self-limited, but complications occur particularly in adults, pregnancy, infants, and people with weakened immunity. Medical assessment may identify a need for antiviral treatment, symptom control, or urgent care.}{ref76}
\indianentry{IND-076}{Smallpox was caused by divine displeasure.}{False}{Smallpox was caused by variola virus, not divine displeasure. It was eradicated through surveillance, isolation, and vaccination. The history demonstrates the value of evidence-based public health for infectious disease.}{ref77}
\indianentry{IND-077}{Food restrictions alone can cure jaundice.}{False}{Jaundice is a sign of excess bilirubin, not a single disease. Causes include hepatitis, medication injury, obstruction, blood disorders, and liver disease. Treatment depends on the cause and may require tests, medicines, procedures, or hospital care.}{ref78}
\indianentry{IND-078}{A person with jaundice should eat only bland or yellow-free food and avoid normal nutrition.}{False}{Food color does not determine bilirubin or cure jaundice. Unless a clinician gives a specific restriction, people generally need adequate energy, protein, and fluids appropriate to their condition. Yellow eyes or skin, dark urine, pale stool, fever, pain, confusion, or bleeding requires assessment.}{ref78}
\indianentry{IND-079}{Sugarcane juice cures jaundice.}{False}{Sugarcane juice does not treat the diseases that can cause jaundice. Unpasteurized juice may carry contamination, and large amounts can be harmful for people with diabetes or vomiting. It is not a substitute for diagnosis or treatment.}{ref78}
\indianentry{IND-080}{A traditional healer can remove jaundice by transferring it into water, leaves, or another object.}{False}{Changing an object cannot measure or remove bilirubin from the body. Temporary symptom changes do not establish that liver or blood disease has resolved. Jaundice should be assessed by a qualified clinician with appropriate examination and testing.}{ref78}
\end{enumerate}


\chapter{Marketing-Created Myths}
Some health myths were not inherited from folklore but manufactured by advertising. This chapter documents claims that were created, amplified, or sustained primarily by commercial marketing, together with what regulators, courts, or published research found. The examples draw on public records, including Federal Trade Commission actions, Food and Drug Administration advisories, court rulings, and peer-reviewed historical analyses. Naming a company reflects the documentary record and is not commentary on any organization's current products. Marketing can spread misinformation without breaking any law, and several campaigns described here predate modern advertising standards. Recognizing the commercial origin of a belief is often the first step toward correcting it.
\begin{enumerate}[leftmargin=0pt,labelwidth=0pt,labelsep=0pt,itemindent=0pt]
\item[] \textbf{MKT-001.} \textbf{Myth:} Adhesive foot pads pull toxins, heavy metals, and parasites out of the body through the soles of the feet overnight.\newline
\textbf{Status:} False.\newline
\textbf{Fact:} Overnight foot pads darken through moisture and heat, not by extracting toxins. Testing cited in federal litigation found the same discoloration whether pads were worn or merely exposed to steam. The marketers of a widely advertised brand were charged by the Federal Trade Commission with falsely claiming scientific proof, and a court later barred the sellers from marketing supplements, foods, drugs, and devices \cite{ref83,ref84}. Sweat glands do not excrete heavy metals in meaningful amounts. Fatigue, hypertension, or joint pain deserves medical evaluation, not adhesive pads. Money spent on detox pads cannot treat an underlying condition. \emph{(See also GEN-060.)}
\item[] \textbf{MKT-002.} \textbf{Myth:} An effervescent cold remedy invented by a schoolteacher prevents colds.\newline
\textbf{Status:} False.\newline
\textbf{Fact:} The inventor story anchored the advertising, but no adequate clinical trial showed that the product prevents colds. The makers settled a false-advertising lawsuit for roughly 23 million dollars in refunds and agreed to add disclaimers after testing failed to support the prevention claims \cite{ref85}. High doses of vitamins do not block respiratory viruses. Prevention measures with real evidence include vaccination, hand hygiene, and adequate sleep. A pleasant-taking remedy is still not a vaccine. Dramatic origin stories are marketing assets, not pharmacology.
\item[] \textbf{MKT-003.} \textbf{Myth:} A daily antiseptic mouthwash can replace flossing and prevent colds and sore throats.\newline
\textbf{Status:} False.\newline
\textbf{Fact:} Mid-twentieth-century advertising claimed this mouthwash cured colds and made flossing unnecessary until regulators forced the claims to stop. Current dental guidance is unchanged: mouthwash may freshen breath or deliver fluoride, but it does not remove plaque between teeth \cite{ref86}. Colds are viral infections unaffected by gargling. Interdental cleaning remains the evidence-based way to address plaque between teeth. A household habit built by advertising is not a clinical recommendation. Familiarity is not proof of benefit.
\item[] \textbf{MKT-004.} \textbf{Myth:} A sugary breakfast cereal improves children's attentiveness by nearly 20 percent compared with other mornings.\newline
\textbf{Status:} False.\newline
\textbf{Fact:} The Federal Trade Commission found that the study behind this claim did not show what the advertisement implied, and the company signed a consent order restricting such advertising \cite{ref87}. Eating any balanced breakfast helps concentration mainly when a child would otherwise eat nothing. Added sugar does not sharpen attention. Diet quality across the day matters more than any single branded meal. Labels deserve more trust than slogans. Parents can compare cereals by fiber and sugar content rather than by performance promises.
\item[] \textbf{MKT-005.} \textbf{Myth:} A cocoa breakfast cereal supports the immune system during flu season.\newline
\textbf{Status:} False.\newline
\textbf{Fact:} The immunity labeling appeared on store shelves and was withdrawn amid criticism shortly afterward \cite{ref88}. Vitamins added to sweetened foods do not tune immune defenses, and no cereal prevents influenza. Immune protection depends on vaccination, sleep, overall nutrition, and exposure patterns. Fortification can correct deficiencies, but it is not a shield against seasonal illness. Medical language borrowed for flu-season marketing deserves particular skepticism. Immunization remains the measure with demonstrated protection.
\item[] \textbf{MKT-006.} \textbf{Myth:} Pomegranate juice slows prostate cancer and treats erectile dysfunction.\newline
\textbf{Status:} False.\newline
\textbf{Fact:} The juice maker funded studies and advertised disease-treatment claims that courts upheld the Federal Trade Commission in rejecting as unsupported by competent and reliable evidence \cite{ref89}. Pomegranate juice contains antioxidants, but drinking it is not cancer therapy. Men with prostate concerns need urological assessment rather than grocery-aisle treatment. Company-funded research requires independent confirmation before clinical conclusions are drawn. Juice also delivers substantial sugar without the fiber of whole fruit. Antioxidant marketing is not a substitute for oncologic care.
\item[] \textbf{MKT-007.} \textbf{Myth:} A probiotic yogurt with special bacterial strains relieves digestive irregularity better than ordinary yogurt.\newline
\textbf{Status:} False.\newline
\textbf{Fact:} European regulators reviewed the evidence and declined to approve the strain-specific digestion claims, while a United States class action produced refunds and label changes \cite{ref90,ref91}. Yogurt can be a nutritious food regardless of branding. Strain-specific benefits require strain-specific evidence, which is rarely as strong as packaging suggests. Persistent bowel symptoms warrant medical assessment rather than repeat purchases. A price premium does not establish clinical superiority. Fermented foods are foods, not regulated medicines.
\item[] \textbf{MKT-008.} \textbf{Myth:} Bracelets with hologram chips or ion technology improve balance, strength, flexibility, or pain.\newline
\textbf{Status:} False.\newline
\textbf{Fact:} The leading hologram seller publicly acknowledged it had no credible scientific basis for its claims and faced refunds, while a separate ionized-bracelet marketer received a multimillion-dollar judgment after a court described its demonstrations as producing a placebo-like effect \cite{ref92,ref93}. Live demonstrations rely on applied-kinesiology tricks that collapse under blinded conditions. Static wrist accessories have no mechanism for altering neuromuscular performance. Chronic pain requires proper diagnosis. Training, rehabilitation, and medical care cannot be replaced by jewelry. Testimonials are not measurements.
\item[] \textbf{MKT-009.} \textbf{Myth:} Curved-sole toning shoes sculpt legs and buttocks during ordinary walking.\newline
\textbf{Status:} False.\newline
\textbf{Fact:} Two major brands refunded tens of millions of dollars after regulators concluded that muscle-toning and calorie-burning claims were unsubstantiated \cite{ref94,ref95}. Walking is genuinely beneficial exercise, but shoe geometry does not multiply its effect. Unstable soles can increase fall risk for some users, particularly older adults. Fitness improvements come from progressive activity, not footwear engineering. Comfort and fit remain the sound reasons for choosing shoes. Advertising photographs of toned models are casting decisions, not data.
\item[] \textbf{MKT-010.} \textbf{Myth:} Intranasal zinc cold gels and swabs are proven, safe ways to shorten colds.\newline
\textbf{Status:} False.\newline
\textbf{Fact:} Regulators warned consumers to stop using three intranasal zinc products after receiving more than 130 reports of smell loss, some apparently permanent, and advised the manufacturer that the products could not be marketed without approval \cite{ref96}. Anosmia removes protection against gas leaks, smoke, and spoiled food. Zinc lozenges are a different exposure from liquid zinc applied inside the nose near olfactory tissue. Homeopathic labeling did not make the products risk-free. Cold remedies should be selected on evidence and safety history. Intranasal zinc specifically should be avoided.
\item[] \textbf{MKT-011.} \textbf{Myth:} Essential oils can treat cancer, autism, and serious infections.\newline
\textbf{Status:} False.\newline
\textbf{Fact:} Distributors of two major oil companies received warning letters demanding removal of unlawful disease-treatment claims \cite{ref97}. Aromatherapy may help some people relax, but concentrated plant oils are not chemotherapy, antibiotics, or neurodevelopmental therapy. Substituting oils for effective care can delay diagnosis and allow disease to progress. Essential oils can irritate skin, poison if swallowed, and interact with medicines. Serious illness requires qualified clinicians. Wellness vocabulary does not turn fragrance into pharmacology.
\item[] \textbf{MKT-012.} \textbf{Myth:} Tiny homeopathic pellets prevent influenza.\newline
\textbf{Status:} False.\newline
\textbf{Fact:} A national health research council reviewed the evidence and concluded that homeopathy is not efficacious beyond placebo, and the maker of a heavily advertised flu pellet paid millions to settle claims about it \cite{ref98,ref99}. At classic homeopathic dilutions, preparations contain no measurable active ingredient. Influenza vaccination remains the preventive tool with demonstrated efficacy. Relying on sugar pellets during flu season risks both illness and transmission to others. Long shelf life and gentle image are not evidence. Prevention claims require prevention trials.
\item[] \textbf{MKT-013.} \textbf{Myth:} Doctors endorsed a cigarette brand as their own favorite, proving it was the sensible choice.\newline
\textbf{Status:} False.\newline
\textbf{Fact:} The famous slogan came from surveys asking physicians which cigarette they smoked, not whether any cigarette was healthy, and the same industry sold filter brands that contained asbestos \cite{ref100}. No tobacco product ever received a genuine medical endorsement. The campaign borrowed professional authority to reassure smokers as lung cancer rates climbed. Such advertising is now taught as a textbook case of misleading implication. White coats in commercials are costumes. Professional imagery is not evidence.
\item[] \textbf{MKT-014.} \textbf{Myth:} Light, mild, and low-tar cigarettes are a safer alternative for smokers.\newline
\textbf{Status:} False.\newline
\textbf{Fact:} Machine-measured tar yields fell, but smokers compensated by inhaling more deeply and covering filter ventilation holes, and major reviews found no meaningful reduction in disease risk \cite{ref101,ref102}. Regulators later banned descriptors such as light and mild precisely because they create false reassurance. The only strategy clearly shown to reduce tobacco risk is cessation. Switching brands preserves nicotine addiction and most harm. Reduced-yield marketing delayed quitting for many users. Industry safety comparisons deserve no benefit of the doubt.
\item[] \textbf{MKT-015.} \textbf{Myth:} Cigarettes marketed as additive-free or natural are less harmful.\newline
\textbf{Status:} False.\newline
\textbf{Fact:} Litigation forced packs and advertisements to carry statements clarifying that additive-free does not mean safer \cite{ref103}. Tobacco harm comes overwhelmingly from combustion products, not from added flavorings alone. Removing additives does not remove tar, carbon monoxide, or carcinogens. Organic and natural are marketing categories here, not toxicological findings. Smokers who switch for perceived safety usually continue smoking rather than quit. The required disclaimer itself concedes there is no safety difference.
\item[] \textbf{MKT-016.} \textbf{Myth:} A base tan protects the skin, and indoor tanning is a safe, necessary source of vitamin D.\newline
\textbf{Status:} False.\newline
\textbf{Fact:} Ultraviolet tanning devices are classified among human carcinogens, and a trade association settlement prohibited false health claims about indoor tanning \cite{ref104,ref105}. A base tan offers sun protection equivalent to a low factor at best while adding cumulative DNA damage. Ordinary casual sunlight meets vitamin D needs for most people, and supplements cover the rest safely. Tanning before a vacation compounds injury rather than preventing it. Melanoma risk rises with indoor tanning, especially when exposure starts young. Golden skin is radiation damage, whatever the brochure says.
\item[] \textbf{MKT-017.} \textbf{Myth:} Dietary fat is the principal dietary cause of heart disease, so sugary foods are heart-safe.\newline
\textbf{Status:} False.\newline
\textbf{Fact:} Internal documents show that a sugar trade group funded a prominent review that minimized evidence linking sucrose to heart disease and selected the scientists and journal without disclosure \cite{ref107}. Contemporary research associates sugary drinks with weight gain, type 2 diabetes, and cardiovascular risk. Fat, carbohydrate quality, and total energy balance all matter; no single villain explains heart disease. Disclosure rules strengthened partly because of such episodes. Industry-shaped consensus deserves scrutiny in every decade. Nutrition advice should always ask who paid for it.
\item[] \textbf{MKT-018.} \textbf{Myth:} The hearty bacon-and-eggs breakfast is the scientifically ideal way to start the day.\newline
\textbf{Status:} False.\newline
\textbf{Fact:} The all-American morning meal was deliberately popularized by a public-relations pioneer hired by a bacon packing company, complete with physician surveys arranged to endorse a heavy breakfast \cite{ref108}. Processed meat is now classified as a human carcinogen, and diets high in saturated fat are associated with cardiovascular risk. A substantial breakfast suits some people, but no science crowns bacon as optimal. Meal patterns matter more than any single iconic plate. One of the most trusted food routines in the United States began as copywriting. Tradition can be manufactured. \emph{(See also GEN-065.)}
\item[] \textbf{MKT-019.} \textbf{Myth:} A glass of fruit juice at breakfast is nutritionally equivalent to eating fruit.\newline
\textbf{Status:} False.\newline
\textbf{Fact:} Decades of citrus promotion helped make juice a breakfast default, but pediatric guidance now recommends limiting juice because it delivers sugar without the fiber and satiety of whole fruit \cite{ref109}. Whole fruit is preferable for children and adults. Frequent juice contributes to dental caries and excess calorie intake. A small serving with a meal is reasonable, but juice is not a fruit replacement. Beverage marketing reshaped the breakfast table more than nutrition science did. Squeezing fruit removes much of what makes fruit healthy.
\item[] \textbf{MKT-020.} \textbf{Myth:} Vitamin-fortified sweetened waters are a healthy way to meet daily nutrient needs.\newline
\textbf{Status:} False.\newline
\textbf{Fact:} Litigation over fortified-water advertising produced label changes and greater disclosure of sugar content \cite{ref106}. Added vitamins do not offset a sugary drink's calories, and people eating ordinary varied diets rarely need supplementation. Plain water remains the appropriate everyday beverage. Fortified drinks create a health halo that can encourage overconsumption. Nutrition panels carry more information than front-of-package promises. Vitamins are nutrients, not redemption credits for sugar.
\item[] \textbf{MKT-021.} \textbf{Myth:} Antibacterial soaps protect households better than plain soap and water.\newline
\textbf{Status:} False.\newline
\textbf{Fact:} Regulators banned numerous active ingredients, including triclosan, from consumer antiseptic washes after manufacturers failed to demonstrate greater safety or effectiveness than ordinary soap \cite{ref110}. Proper handwashing technique with plain soap removes pathogens effectively. Some antibacterial agents raised concerns about hormonal effects and antimicrobial resistance. Alcohol-based sanitizers are a distinct and useful tool when washing is unavailable. Decades of premium pricing purchased anxiety relief, not extra protection. The word antibacterial on a bottle is a claim, not a measurement.
\item[] \textbf{MKT-022.} \textbf{Myth:} Alkaline or ionized water neutralizes body acidity and prevents disease.\newline
\textbf{Status:} False.\newline
\textbf{Fact:} Blood pH is tightly regulated near neutrality, and systematic reviews find no meaningful human evidence for the proposed benefits of alkaline water \cite{ref111}. Diet cannot appreciably shift blood pH; urine changes reflect normal kidney function. Gastric acid neutralizes alkaline water within minutes of drinking. Hydration needs are met by ordinary water at ordinary prices. Expensive machines solve a problem that physiology says does not exist. Multilevel marketing incentives explain the enthusiasm better than biochemistry does.
\item[] \textbf{MKT-023.} \textbf{Myth:} Cellulite is a disease that special creams can eliminate.\newline
\textbf{Status:} False.\newline
\textbf{Fact:} Dermatology literature describes cellulite as normal subcutaneous architecture, common in most postadolescent women, and notes that the condition was substantially invented as a cosmetic concern by mid-twentieth-century media \cite{ref112}. Topical creams cannot restructure the fibrous bands beneath skin, and no preparation reliably removes the dimpled appearance. Massage and devices produce temporary changes at best. Products promising elimination monetize a normal body feature. New, painful, or rapidly changing skin changes belong with a clinician. A cosmetics counter is not a diagnostic clinic.
\item[] \textbf{MKT-024.} \textbf{Myth:} Collagen creams rebuild youthful skin by restoring collagen from a jar.\newline
\textbf{Status:} False.\newline
\textbf{Fact:} Topical collagen molecules are generally too large to penetrate the epidermis and reach the dermal layer where structural collagen actually resides \cite{ref113}. Collagen-containing formulas can moisturize and temporarily smooth surface texture. Interventions with stronger evidence include sun protection, retinoid treatment, and procedures performed by qualified professionals. Delivery physics matters more than a hero molecule printed on the label. Skincare marketing frequently borrows laboratory language without laboratory results. Healthy skepticism is a better companion to a mirror than hope in a jar.
\end{enumerate}

\chapter{Alternative Medicine Myths}
Traditional systems such as homeopathy, Ayurveda, and traditional Chinese medicine carry deep cultural history, and many practices offer comfort, community, and physical activity. This chapter evaluates specific health claims against clinical evidence and regulatory findings; it does not dismiss the cultures from which these practices come. Three recurring problems appear across the examples. First, plausible traditions can still rest on incorrect biology. Second, natural origin does not guarantee safety, because some preparations contain potent pharmacological agents, toxic metals, or undeclared drugs. Third, the greatest harm usually comes not from the practice itself but from replacing effective care with an unproven substitute. Readers who value these traditions can apply the same standard used throughout this book: examine the specific claim, the evidence behind it, and what regulators found.
\section{Homeopathy}
\begin{enumerate}[leftmargin=0pt,labelwidth=0pt,labelsep=0pt,itemindent=0pt]
\item[] \textbf{ALT-001.} \textbf{Myth:} Substances that cause symptoms in healthy people will cure the same symptoms in sick people.\newline
\textbf{Status:} False.\newline
\textbf{Fact:} This principle, formulated around 1796, was a guess about biology that subsequent science has not supported. Symptoms are responses to illness, not toxins to be matched and canceled. Comprehensive national assessments have concluded that homeopathy is not efficacious beyond placebo for any condition studied \cite{ref98}. Homeopathy also differs fundamentally from herbal medicine, because classic homeopathic preparations contain little or no active substance at all. Symptom matching sounds elegant, but elegance is not evidence. Treatments earn adoption through controlled trials, not through eighteenth-century intuition.
\item[] \textbf{ALT-002.} \textbf{Myth:} The more a homeopathic remedy is diluted, the stronger it becomes, because water remembers the original substance.\newline
\textbf{Status:} False.\newline
\textbf{Fact:} Common pharmacy potencies involve dilutions far beyond the point where any molecule of the original substance remains; a 30C preparation corresponds to roughly one part in ten raised to the sixtieth power. The memory-of-water idea contradicts basic chemistry, because solution structure persists for only fractions of a second. High-quality trials show no benefit regardless of dilution level \cite{ref98}. If dilution truly increased potency, ordinary tap water, diluted countless times over millennia, would be the strongest medicine on earth. Dilution removes pharmacology rather than concentrating it.
\item[] \textbf{ALT-003.} \textbf{Myth:} Homeopathic nosodes can replace childhood vaccines.\newline
\textbf{Status:} False.\newline
\textbf{Fact:} Nosodes are diluted preparations made from diseased tissue or discharge, and they generate no measurable antibody response. European science academies have publicly warned against substituting homeopathic products for vaccination, and national reviews reached the same conclusion \cite{ref98,ref114}. Vaccine-preventable diseases return quickly when coverage falls, as measles resurgences demonstrate. Parents offered nosodes instead of routine immunization are being sold false protection. The stakes are highest for infants too young to be fully vaccinated. Immunization schedules exist because they work; nosodes do not. \emph{(See also MKT-012.)}
\item[] \textbf{ALT-004.} \textbf{Myth:} Homeopathy must work because it helps babies and animals, who cannot experience placebo effects.\newline
\textbf{Status:} False.\newline
\textbf{Fact:} The argument confuses the patient's expectations with those of everyone observing the patient. A parent's rating of an infant's colic, or an owner's impression of a pet's arthritis, is strongly shaped by hope and by the natural course of the condition. Blinded trials in infants and veterinary populations show the same pattern seen elsewhere: no reliable effect beyond placebo \cite{ref98}. Improvement after a remedy usually reflects time, growth, or attention. Observers can be treated even when patients cannot. Anecdotes from loving observers are the least reliable evidence precisely because caring distorts judgment.
\end{enumerate}
\section{Ayurvedic and Traditional Indian Claims}
\begin{enumerate}[leftmargin=0pt,labelwidth=0pt,labelsep=0pt,itemindent=0pt]
\item[] \textbf{ALT-005.} \textbf{Myth:} Metal-based Ayurvedic preparations are safe because traditional processing neutralizes their toxicity.\newline
\textbf{Status:} False.\newline
\textbf{Fact:} Some preparations deliberately incorporate lead, mercury, or arsenic, and laboratory analyses have repeatedly detected toxic metals in a substantial fraction of products, whether manufactured in the United States or India \cite{ref115,ref116}. Documented cases of lead poisoning, particularly in children given tonics for development or digestion, trace directly to such products. Processing standards vary widely between producers, and no certification guarantees a safe blood lead level. Heavy metals accumulate silently, damaging nerves, kidneys, and development before symptoms appear. Traditional use demonstrates cultural continuity, not toxicological clearance. Anyone using these products, especially for children, deserves blood testing and medical advice.
\item[] \textbf{ALT-006.} \textbf{Myth:} Herbal products are automatically safe because they are natural, and cannot interact with medicines.\newline
\textbf{Status:} False.\newline
\textbf{Fact:} Plants are chemistry, and some of medicine's most potent drugs came from them. St John's wort induces liver enzymes that can render contraceptives, transplant medicines, and anticoagulants less effective, and several herbs increase bleeding risk during surgery \cite{ref117}. Regulatory reviews document contamination, misidentification, and undeclared ingredients in commercial herbal products. Natural origin describes provenance, not a safety profile. Patients should disclose every supplement before operations and new prescriptions. \emph{(See also KID-103.)}
\item[] \textbf{ALT-007.} \textbf{Myth:} Drinking one's own urine cures infections, allergies, or chronic disease.\newline
\textbf{Status:} False.\newline
\textbf{Fact:} Urine is the end product of filtering blood, containing urea, salts, hormones, and any drugs the body is clearing. Reintroducing it adds concentrated waste back to the kidneys while providing no established therapeutic effect. Small amounts are unlikely to cause acute poisoning, but it is neither sterile nor nutritious, and bacterial growth begins quickly outside the bladder. Claims that urine therapy treats cancer, infertility, or viral illness lack clinical trial support. The practice persists through testimonials rather than evidence. Wasting time on it delays diagnosis of treatable disease. \emph{(See also IND-039.)}
\end{enumerate}
\section{Meditation and Yoga}
\begin{enumerate}[leftmargin=0pt,labelwidth=0pt,labelsep=0pt,itemindent=0pt]
\item[] \textbf{ALT-008.} \textbf{Myth:} Yoga alone can cure diabetes, thyroid disease, or cancer without medication.\newline
\textbf{Status:} False.\newline
\textbf{Fact:} Yoga genuinely improves flexibility, balance, stress measures, and quality of life, and comparisons with other exercise show overlapping benefits for general health \cite{ref118}. Those benefits do not extend to eliminating autoimmune thyroiditis, destroying tumors, or replacing glucose-lowering therapy. Stopping insulin or thyroid hormone for a movement practice can cause rapid, life-threatening deterioration. The most defensible role for yoga is alongside medical care, where it may improve adherence and wellbeing. Ancient disciplines deserve accurate credit, not inflated claims. Stretching is not chemotherapy.
\item[] \textbf{ALT-009.} \textbf{Myth:} Meditation done properly means emptying the mind and stopping all thoughts.\newline
\textbf{Status:} False.\newline
\textbf{Fact:} Minds produce thoughts the way hearts produce beats; every tradition's own instruction manuals describe noticing thoughts and gently returning attention, not achieving blankness. Struggling to suppress thoughts is the most common reason beginners quit, convinced they are failing at something everyone else finds easy. Wandering attention followed by redirection is the exercise itself, analogous to a repetition in strength training. There is no failing state of meditation other than never practicing. Expecting instant silence sets an impossible standard. A busy mind meditating is normal minds at work.
\item[] \textbf{ALT-010.} \textbf{Myth:} Mindfulness meditation can replace antidepressants and psychiatric treatment.\newline
\textbf{Status:} False.\newline
\textbf{Fact:} Structured mindfulness programs produce small-to-moderate improvements in anxiety, depression, and stress measures in meta-analysis, and one specific program reduces relapse in recurrent depression when added to care \cite{ref136}. Those findings do not support swapping meditation for treatment of moderate-to-severe illness, psychosis, or acute suicidality. Stopping prescribed medication abruptly risks relapse and withdrawal effects. Meditation works best as a component of care, not as its replacement. People considering changes should involve their prescriber first. Breathing exercises complement treatment; they do not substitute for it.
\item[] \textbf{ALT-011.} \textbf{Myth:} Meditation has no side effects and is beneficial for absolutely everyone.\newline
\textbf{Status:} False.\newline
\textbf{Fact:} Systematic study documents meditation-related adverse experiences, including anxiety, dissociation, sleep disturbance, and reactivation of traumatic memories, particularly during intensive retreats \cite{ref134}. Pooled estimates suggest a meaningful minority of practitioners experience at least transient difficulties, and people with trauma histories or psychotic disorders face higher risk. This does not make ordinary practice dangerous for most people, any more than exercise being good for most people erases injuries. Screening, graduated intensity, and qualified instruction reduce risk. Universal prescriptions deserve suspicion in every domain of health. Even breathing practices have contraindications.
\item[] \textbf{ALT-012.} \textbf{Myth:} Meditation apps rapidly rewire the brain and deliver guaranteed transformation within weeks.\newline
\textbf{Status:} False.\newline
\textbf{Fact:} Marketing borrows neuroscience language far beyond what studies show. Large reviews find modest psychological benefits that shrink in the best-controlled comparisons, and dramatic structural brain-change claims largely fail replication \cite{ref136}. Consistent practice supports skills like attention regulation and stress tolerance gradually, similar to fitness gains. Anyone promising transformation in seven days is selling urgency, not evidence. Free recordings work about as well as premium subscriptions for most outcomes. Slower, realistic expectations predict persistence better than hype.
\item[] \textbf{ALT-013.} \textbf{Myth:} Hot yoga detoxifies the body through profuse sweating.\newline
\textbf{Status:} False.\newline
\textbf{Fact:} Sweat consists mainly of water and electrolytes, with only trace amounts of some compounds; the organs that clear genuine toxins are the liver and kidneys, working continuously regardless of room temperature. Heated rooms raise heart rate and perceived effort, which some people enjoy, but they also increase dehydration, heat exhaustion, and fainting risk. Weight lost in a hot class returns with the first glass of water. Flexibility feels greater in heat, which can tempt overstretching. Detoxification claims belong to marketing, not physiology. \emph{(See also GEN-088 and KID-036.)}
\item[] \textbf{ALT-014.} \textbf{Myth:} Yoga reliably cures chronic low back pain.\newline
\textbf{Status:} Uncertain.\newline
\textbf{Fact:} Randomized trials comparing yoga with stretching or self-care found modest improvements in back-related function and symptoms lasting months, placing yoga among reasonable active options alongside exercise therapy \cite{ref135}. Cures are another matter: relapses occur, effects are partial, and results depend on style, instructor training, and adaptation to the individual. Aggressive postures can worsen certain spine conditions. For people who enjoy it, supervised gentle yoga is a defensible part of a recovery plan. It is one tool among several, not a guaranteed remedy. Claiming cure oversells what trials actually measured.
\item[] \textbf{ALT-015.} \textbf{Myth:} Breathing exercises cure asthma and boost immunity to infection.\newline
\textbf{Status:} False.\newline
\textbf{Fact:} Breathing retraining may modestly improve symptom scores and quality-of-life measures in asthma, but it does not reverse airway inflammation, and stopping inhalers in favor of breathwork has led to severe attacks. Intense hyperventilation styles cause dizziness, tingling, and fainting by lowering carbon dioxide sharply. Immune boosting through respiration patterns lacks credible human evidence. Controlled breathing is a legitimate adjunct for relaxation and panic-type symptoms. Asthma management rests on inhaled therapy and an action plan. Breath is powerful medicine for the nervous system, not for airways.
\item[] \textbf{ALT-016.} \textbf{Myth:} Advanced yoga poses are safe for anyone flexible enough, and pain during practice means progress.\newline
\textbf{Status:} False.\newline
\textbf{Fact:} Reviews of adverse events document strains, ligament injuries, disc problems, and rare arterial dissections associated with extreme neck movements and forceful postures \cite{ref137}. Hypermobility allows reaching impressive positions without the muscular control that protects joints, accumulating silent damage. Pain signals tissue stress, not spiritual breakthrough. Props, modifications, and honest instruction exist precisely because bodies differ. Progress in yoga is usually measured in steadier breathing and fewer injuries, not deeper contortion. No posture is worth a shoulder reconstruction.
\end{enumerate}
\section{Traditional Chinese Medicine}
\begin{enumerate}[leftmargin=0pt,labelwidth=0pt,labelsep=0pt,itemindent=0pt]
\item[] \textbf{ALT-017.} \textbf{Myth:} Aristolochia-containing Chinese herbs are safe traditional remedies for weight loss and joint pain.\newline
\textbf{Status:} False.\newline
\textbf{Fact:} Aristolochic acid is one of the most consequential herbal toxins ever identified. A Belgian slimming clinic's patients developed rapidly progressive kidney failure requiring dialysis or transplantation, and follow-up found urothelial abnormalities in nearly half of those examined \cite{ref119}. International cancer authorities classify aristolochic acid as a human carcinogen, and related compounds cause endemic kidney disease in agricultural regions of Asia and Europe \cite{ref120}. Many countries restrict these herbs, yet substitutions and online sales continue. Kidney damage from this herb family is permanent. Herbal labels do not reliably reveal aristolochia content.
\item[] \textbf{ALT-018.} \textbf{Myth:} Herbal weight-loss capsules contain only plant ingredients.\newline
\textbf{Status:} False.\newline
\textbf{Fact:} Regulators have repeatedly found undeclared prescription drugs, including withdrawn appetite suppressants associated with heart valve damage, hidden inside products marketed as natural slimming teas and capsules \cite{ref121}. Adulteration exists because the drugs work briefly, creating satisfied early customers before the harms appear. Consumers cannot distinguish a clean product from an adulterated one by taste, smell, or packaging. Rapid unexplained weight loss from an herbal pill should prompt suspicion, not celebration. Purchasing regulated treatments through licensed channels remains the safer path. Natural labeling is a marketing decision, not a formulation guarantee.
\item[] \textbf{ALT-019.} \textbf{Myth:} Rhino horn and tiger bone have proven medicinal value that justifies their use.\newline
\textbf{Status:} False.\newline
\textbf{Fact:} Rhino horn is dense keratin, biochemically similar to compressed fingernail and hair, with no credible pharmacological activity for fever, cancer, or any other condition \cite{ref122}. Tiger bone offers no demonstrated advantage over ordinary analgesics. Demand for these products drives organized poaching that has pushed wild populations toward extinction, creating a conservation catastrophe without any offsetting health benefit. Legal herbal and pharmaceutical alternatives exist for every claimed indication. Using endangered-animal products harms species and ecosystems while treating nothing. Medicine advances by abandoning remedies that never worked.
\item[] \textbf{ALT-020.} \textbf{Myth:} Acupuncture can cure infertility, allergies, asthma, or cancer.\newline
\textbf{Status:} False.\newline
\textbf{Fact:} Systematic reviews of acupuncture alongside fertility treatment found no meaningful improvement in clinical pregnancy or live-birth outcomes \cite{ref123}. Evidence for curing allergic disease, asthma, or malignancy is similarly lacking, and no needling protocol replaces oncologic or respiratory therapy. Needles are generally safe when sterile, and sessions may feel relaxing, which explains part of their appeal. Comfort measures are legitimate, but they should be labeled honestly. Disease modification requires disease-modifying evidence. Paying for cure claims that reviews reject wastes money needed for effective care.
\item[] \textbf{ALT-021.} \textbf{Myth:} Acupuncture provides meaningful, lasting relief for chronic back, neck, and joint pain.\newline
\textbf{Status:} Uncertain.\newline
\textbf{Fact:} A large individual-patient-data meta-analysis found small but statistically significant short-term advantages of acupuncture over both sham needling and usual care for chronic back, neck, shoulder, and osteoarthritis pain \cite{ref124}. Effect sizes were modest, differences between true and sham needling were smaller than differences from no treatment, and durability beyond months is unclear. Sham-controlled designs struggle here because patients recognize penetrating needles. For people seeking non-drug options, a trial course alongside standard care is defensible. Expecting transformation is not supported. Relief that is modest, partial, and temporary should be described exactly that way.
\item[] \textbf{ALT-022.} \textbf{Myth:} Cupping draws toxins out of the body through the skin.\newline
\textbf{Status:} False.\newline
\textbf{Fact:} The dark circles left by cupping are bruises created when suction ruptures small superficial vessels, and they fade as the body repairs them. Sweat glands cannot excrete heavy metals or metabolic waste in meaningful quantities, so no suction device extracts toxins through intact skin. Some athletes and patients report temporary relief or relaxation, which may reflect local blood flow changes and expectation. Measurable toxin removal has never been demonstrated in properly designed studies. Bruising is evidence of capillary injury, not of purification. \emph{(See also GEN-060.)}
\end{enumerate}
\section{Energy, Manual, and Crystal Practices}
\begin{enumerate}[leftmargin=0pt,labelwidth=0pt,labelsep=0pt,itemindent=0pt]
\item[] \textbf{ALT-023.} \textbf{Myth:} Practitioners of energy therapies can feel and manipulate human energy fields to heal organs.\newline
\textbf{Status:} False.\newline
\textbf{Fact:} No instrument has detected the proposed biofield, and practitioners' claims collapse under simple testing. In a study published in a leading medical journal, experienced therapeutic-touch practitioners, blindfolded, identified which of their hands hovered above the researcher's hand barely more often than coin-flip odds would predict \cite{ref125}. Organs cannot be healed by adjusting a field that has never been shown to exist. Gentle presence and attention can be calming, and that benefit belongs to human contact rather than to mysticism. Patients with serious illness deserve treatments whose mechanisms survive scrutiny. Feeling energy is not the same as measuring it.
\item[] \textbf{ALT-024.} \textbf{Myth:} Spinal manipulation cures infant colic, ear infections, and asthma.\newline
\textbf{Status:} False.\newline
\textbf{Fact:} Controlled evidence does not support manipulation as treatment for these conditions, and infants' spines and necks are structurally vulnerable to forces that adults tolerate \cite{ref126}. Colic resolves with maturation in most babies, which makes any intervention look successful to exhausted parents. Ear infections and asthma require medical assessment, and delayed treatment carries real consequences. Case reports describe serious injuries following pediatric spinal manipulation. A practitioner trained in adult spines is not thereby qualified to diagnose infant illness. Fussy babies need pediatric evaluation, not high-velocity adjustments.
\item[] \textbf{ALT-025.} \textbf{Myth:} Pressing specific zones of the feet treats the corresponding internal organs.\newline
\textbf{Status:} False.\newline
\textbf{Fact:} Reflexology maps have no anatomical basis; no nerve wiring connects a patch of sole to the liver, kidneys, or uterus in the way the charts depict. Reviews of randomized trials find inconsistent results that do not establish organ-specific effects beyond relaxation and pleasant interaction \cite{ref127}. Foot massage can lower distress and feels good, and there is nothing wrong with enjoying it as comfort care. Problems arise when charts are used to diagnose disease or replace endocrine, cardiac, or cancer treatment. A soothing touch is not a diagnostic instrument. Maps drawn by intuition remain intuition.
\item[] \textbf{ALT-026.} \textbf{Myth:} Crystals heal disease through their energy vibrations.\newline
\textbf{Status:} False.\newline
\textbf{Fact:} Crystals have stable atomic lattices, not healing frequencies, and no experiment has shown disease modification from proximity to quartz, amethyst, or any other stone. Believers often report warmth or tingling, yet blinded designs erase these sensations, indicating expectation rather than emission. Collecting beautiful minerals is harmless and even delightful. Harm appears only when stones stand in for insulin, antibiotics, or psychiatric care. Placebos borrowed from geology still leave pathology untouched. Enjoy crystals as objects; consult clinicians as clinicians.
\end{enumerate}
\section{Unproven Cancer and Detoxification Therapies}
\begin{enumerate}[leftmargin=0pt,labelwidth=0pt,labelsep=0pt,itemindent=0pt]
\item[] \textbf{ALT-027.} \textbf{Myth:} Apricot kernels or laetrile, marketed as vitamin B17, cure cancer safely.\newline
\textbf{Status:} False.\newline
\textbf{Fact:} Laetrile is a chemically modified plant compound that releases cyanide during metabolism, and a rigorous clinical trial found no anticancer activity whatsoever \cite{ref128}. Poisoning and deaths from cyanide toxicity occurred among people self-treating with kernels and illicit preparations, particularly children who swallow handfuls. Calling laetrile a vitamin was a marketing device; no nutritional deficiency corresponds to it. The remedy spread through clinics beyond regulatory oversight during the nineteen-seventies and failed wherever measured. Cancer centers offer therapies selected by survival data, not by smuggled promises. Bitter kernels are poison, not nutrition.
\item[] \textbf{ALT-028.} \textbf{Myth:} Intensive juicing regimes combined with coffee enemas cure cancer by detoxifying the liver.\newline
\textbf{Status:} False.\newline
\textbf{Fact:} This protocol claims that frequent coffee enemas open bile ducts so the liver can expel tumors, an assertion never supported by controlled evidence \cite{ref129}. Complications documented in followers include severe electrolyte imbalance, bowel inflammation, infection, and deaths linked to the enemas themselves. Prolonged restrictive regimens cause protein and calorie deficiency precisely when cancer patients need nutrition most. The liver already performs its own detoxification continuously without assistance. Choosing this program instead of oncologic care converts a treatable disease into a fatal one. Vegetables help health; enemas do not direct them to tumors.
\item[] \textbf{ALT-029.} \textbf{Myth:} Black salve safely burns out skin cancer while sparing healthy tissue.\newline
\textbf{Status:} False.\newline
\textbf{Fact:} Black salve contains sanguinarine and related caustic agents that destroy whatever tissue they contact, producing deep wounds called eschars indiscriminately. Photographs submitted to regulators show craters, cartilage loss, and disfigurement, frequently with cancer remaining underneath and spreading afterward \cite{ref130}. Regulators list these products among fraudulent cancer cures and have warned sellers repeatedly. Skin lesions require diagnosis first, because benign and malignant conditions can look alike. Destruction without diagnosis is not treatment. Real skin-cancer surgery preserves healthy tissue and confirms clear margins under a microscope.
\item[] \textbf{ALT-030.} \textbf{Myth:} Chelation therapy cures autism by removing hidden mercury and other metals.\newline
\textbf{Status:} False.\newline
\textbf{Fact:} Autism is a neurodevelopmental condition, not metal poisoning, and trials of chelation for autistic children show no benefit. Chelating agents strip calcium and essential minerals along with trace metals, and therapy has caused fatal hypocalcemia, including the death of a child receiving it for unproven indications \cite{ref131}. Regulators warn against over-the-counter chelation products marketed for autism, heart disease, and other conditions. Legitimate chelation is reserved for confirmed heavy-metal poisoning, dosed and monitored medically. Subjecting children to dangerous infusions for a disproven theory is harm, not hope. Behavioral and educational therapies remain the supported interventions.
\item[] \textbf{ALT-031.} \textbf{Myth:} Colloidal silver strengthens immunity and cures infections, including serious ones.\newline
\textbf{Status:} False.\newline
\textbf{Fact:} Silver solutions have no demonstrated ability to prevent or treat infection in humans when ingested, and regulators decades ago barred over-the-counter silver products from making such claims \cite{ref132}. Swallowed silver deposits permanently in skin and eyes, producing argyria, an irreversible blue-gray discoloration visible for life. Silver can also block absorption of common medicines, including thyroid hormone and certain antibiotics. Topical medical silver dressings are a different, clinically validated application managed by professionals. Drinking metal particles is not immune support. The color change is permanent proof of misplaced trust.
\item[] \textbf{ALT-032.} \textbf{Myth:} Ear candles vacuum wax and toxins out of the ear canal safely.\newline
\textbf{Status:} False.\newline
\textbf{Fact:} Measurements show that burning hollow candles create no meaningful negative pressure, and trials found the waxy residue inside used candles comes from the candle itself, not the ear \cite{ref133}. Burns to the ear canal, face, and hair, and candle wax occluding the canal, are documented injuries. Impacted cerumen responds to softening drops, irrigation, or manual removal by a clinician. Toxins are not removed from the head by combustion beside it. The technique fails physics, anatomy, and safety simultaneously. Cotton swabs pushed deep into ears cause more harm than the wax they chase.
\end{enumerate}

\chapter{Asian Food and Medicine Myths}
Across Asia, food and medicine have shared a single vocabulary for millennia, and many remedies began as reasonable observations in their own time. This chapter examines specific food-as-medicine claims from China, Japan, Korea, and Southeast Asia, together with everyday beliefs that carry health consequences. India receives its own dedicated treatment elsewhere in this book. Three patterns recur. Foods praised as near-miraculous usually contain ordinary nutrients that the body digests like any other meal. Some traditional preparations carry genuine pharmacological activity along with genuine toxicity, which is precisely why regulated pharmaceutical versions exist. And the most serious harm rarely comes from enjoying a traditional food; it comes from substituting it for vaccination, anticoagulation, insulin, or cancer care. Enjoying cuisine and evaluating claims are compatible activities.
\section{Chinese Food as Medicine}
\begin{enumerate}[leftmargin=0pt,labelwidth=0pt,labelsep=0pt,itemindent=0pt]
\item[] \textbf{ASM-001.} \textbf{Myth:} Bird's nest soup builds immunity, beautifies skin, and justifies prices above silver.\newline
\textbf{Status:} False.\newline
\textbf{Fact:} Edible bird's nest is swiftlet saliva, built mostly from glycoproteins with some sialic acid, and laboratory analysis shows its proteins are digested like any other food protein \cite{ref138}. Rigorous human trials do not support advertised benefits for skin, immunity, or longevity. The trade has documented adulteration with bleaching agents, gelatin, and sugar to inflate weight and appearance. As an occasional delicacy it is harmless for most people; as medicine priced above precious metals it is theater. Expensive molecules are still digested into ordinary amino acids. Skin glows no better on soup than on sleep.
\item[] \textbf{ASM-002.} \textbf{Myth:} Shark fin soup promotes virility, longevity, and protection against cancer.\newline
\textbf{Status:} False.\newline
\textbf{Fact:} Shark fin is cartilage, whose collagen offers no advantage over the protein in an omelet, and shark-cartilage supplements already failed controlled cancer trials. Sharks accumulate mercury efficiently, and regulatory advisories list shark among the species pregnant people and children should avoid entirely \cite{ref139}. The fin itself is nearly tasteless, borrowing flavor from broth while carrying the contamination. Finning has devastated shark populations across oceans for a dish with no nutritional distinction. No virility pathway connects cartilage to hormones. The safest description of shark fin soup is expensive risk.
\item[] \textbf{ASM-003.} \textbf{Myth:} Bear bile treats fever and liver disease, and no substitute exists.\newline
\textbf{Status:} False.\newline
\textbf{Fact:} Bear bile contains ursodeoxycholic acid, a genuine medication for certain liver and gallstone conditions, but pharmaceutical companies synthesize the identical molecule cheaply and cleanly, making bear farms medically obsolete \cite{ref140}. Farmed bile varies wildly in purity and carries bacterial contamination and impurities absent from regulated tablets. Demand keeps bears in extraction cages for decades, a welfare catastrophe with no therapeutic justification. The claim that tradition requires the animal product fails pharmacology twice over. Same molecule, sterile factory, no bear. Every medical indication for bile acid therapy can be met without harming one.
\item[] \textbf{ASM-004.} \textbf{Myth:} Ginseng restores energy, immunity, fertility, and youth, working as a general panacea.\newline
\textbf{Status:} False.\newline
\textbf{Fact:} Systematic reviews of randomized trials find most claimed benefits supported only by small, low-quality studies, with possible modest effects on fatigue-like symptoms and nothing approaching panacea \cite{ref141}. Ginseng also interacts with medicines, including warfarin, and reports of insomnia and blood-pressure effects exist. Quality varies enormously between products, with some commercially sold ginseng containing little of the labeled species. Enjoying ginseng tea is harmless for most adults. Treating it as a cure-all delays real diagnosis of fatigue causes, which include thyroid disease, anemia, sleep apnea, and depression. Adaptogen marketing is not evidence.
\item[] \textbf{ASM-005.} \textbf{Myth:} Cordyceps and reishi mushrooms cure cancer, kidney disease, and immune disorders.\newline
\textbf{Status:} False.\newline
\textbf{Fact:} Laboratory experiments showing effects on isolated cells do not translate into cures for humans swallowing capsules, and no controlled trial has demonstrated either mushroom treating malignancy or organ failure \cite{ref142}. Wild cordyceps is a parasite that kills caterpillars; commercial products are fermented powders of variable composition sold in an industry with minimal quality enforcement. Immune-active chemistry in a dish does not equal immunotherapy. Patients who replace oncologic care with mushroom regimens return with advanced disease. Mushrooms deserve credit as food and as a source of drug leads, not as finished cures. A lead is where medicine starts, not where it ends.
\item[] \textbf{ASM-006.} \textbf{Myth:} Eating an animal organ strengthens the corresponding human organ, so brains feed brains and hearts feed hearts.\newline
\textbf{Status:} False.\newline
\textbf{Fact:} Digestion dismantles all tissue, whatever its shape, into amino acids, fats, vitamins, and minerals that the body redistributes according to need. Eating liver does not route nutrition to your liver by recognition, and brain tissue confers no cognitive benefit beyond ordinary protein. Organ meats are nutrient-dense foods, sometimes rich in iron or B vitamins, which explains part of their reputation. They can also concentrate heavy metals and contaminants acquired by the source animal. The doctrine of similar shapes treating similar parts predates knowledge of metabolism entirely. Symmetry feels intuitive, but intestines cannot read intentions.
\item[] \textbf{ASM-007.} \textbf{Myth:} Snake wine and raw snake blood convey vitality and treat joint disease safely.\newline
\textbf{Status:} False.\newline
\textbf{Fact:} Raw reptile blood and tissue transmit salmonella and parasitic infections, including sparganosis, a larval tapeworm disease documented after consuming raw snake and flesh exposures. Alcohol does not sterilize reliably at drinking strengths, and venom inside sealed wine is a lesser concern than the microbes. Handlers face envenomation during preparation, an occupational injury with no compensating benefit. No trial shows vitality or joint improvement beyond alcohol's familiar warmth. The spectacle of the serpent does the persuasive work. Vitality purchased this way arrives with infectious-disease interest charges.
\end{enumerate}
\section{Japanese Claims}
\begin{enumerate}[leftmargin=0pt,labelwidth=0pt,labelsep=0pt,itemindent=0pt]
\item[] \textbf{ASM-008.} \textbf{Myth:} Nattokinase dissolves blood clots and can replace prescription anticoagulants.\newline
\textbf{Status:} False.\newline
\textbf{Fact:} Nattokinase shows fibrin-dissolving activity in test tubes, and small short-term trials suggest minor blood-pressure effects, neither of which establishes stroke prevention in people \cite{ref143}. Swallowed enzymes face destruction by stomach acid and digestive proteases, limiting what reaches circulation intact. Natto is also among the richest vitamin K sources, directly opposing warfarin, so patients on anticoagulation face contradictory forces if they self-manage. Stopping prescribed anticoagulants for fermented soybeans invites the clots the medication prevents. Natto remains a nutritious traditional breakfast. It is not warfarin, and pretending otherwise has sent people to stroke units.
\item[] \textbf{ASM-009.} \textbf{Myth:} Drinking several glasses of water immediately upon waking cures diabetes, arthritis, constipation, and cancer.\newline
\textbf{Status:} False.\newline
\textbf{Fact:} Morning hydration is pleasant and harmless in normal amounts, but no fluid schedule reverses autoimmune, degenerative, or malignant disease. The kidneys regulate blood volume continuously regardless of when water arrives. Protocols urging liters before breakfast risk hyponatremia in extreme adherence, a genuinely dangerous sodium dilution. Constipation may improve with overall fluid intake, which explains the practice's kernel of truth. Cures attributed to the routine reflect time, placebo, and concurrent treatments. Water sustains life; it does not select targets.
\item[] \textbf{ASM-010.} \textbf{Myth:} Eating the marketed Okinawan superfoods replicates Okinawan centenarian lifespans.\newline
\textbf{Status:} False.\newline
\textbf{Fact:} Okinawa's extraordinary longevity was real, and researchers attribute it to a combination of genetics, caloric moderation, active lifestyles, social cohesion, and healthcare rather than to any single sweet potato, tofu dish, or seaweed \cite{ref144}. Younger Okinawans who adopted Westernized diets saw life expectancy advantages shrink, which is difficult to reconcile with a magic-food theory. Diet books cherry-pick the appealing elements while omitting wartime scarcity and daily labor. No imported purple potato conveys community structure. Longevity recipes sell because mortality resists advertising scrutiny. Eat well by all means; expect no century in exchange.
\end{enumerate}
\section{Korean Beliefs}
\begin{enumerate}[leftmargin=0pt,labelwidth=0pt,labelsep=0pt,itemindent=0pt]
\item[] \textbf{ASM-011.} \textbf{Myth:} Kimchi prevented SARS in 2003 and protects against respiratory viruses including COVID-19.\newline
\textbf{Status:} False.\newline
\textbf{Fact:} Fermented vegetables are a nutritious food group, and meta-analysis of probiotic trials shows at most a modest reduction in respiratory infection frequency and duration \cite{ref145}. Nothing in that evidence touches prevention of specific viral epidemics, and SARS-era sales surges followed unverified speculation, not trials. Kimchi's capsaicin and garlic content irritate rather than disinfect mucous membranes. South Korea fought COVID-19 with testing, tracing, and vaccination, not with cabbage. Loving a national dish needs no medical exaggeration. Fermentation preserves vegetables; it does not immunize anyone.
\item[] \textbf{ASM-012.} \textbf{Myth:} Sleeping in a closed room with an electric fan running can kill a healthy person overnight.\newline
\textbf{Status:} False.\newline
\textbf{Fact:} Fan death survives despite no verified forensic case report of a healthy person dying this way anywhere in the world literature. Proposed mechanisms fail basic physiology: fans neither create vacuums nor consume oxygen, and moving air cools rather than chills a resting body. Officialdom long reinforced belief instead of dispelling it, with a government consumer safety alert in 2006 listing fan ``asphyxiation'' among Korea's most common summer injuries \cite{ref146}. Media reports still attribute unexplained overnight deaths to running fans when autopsies point to undiagnosed heart disease or alcohol. Heat waves kill through ambient temperature, and fans remain a recommended low-energy cooling measure in ordinary conditions. A fan is a wind machine, not an assassin.
\item[] \textbf{ASM-013.} \textbf{Myth:} Dog-meat tonics build stamina and male vigor beyond ordinary nutrition.\newline
\textbf{Status:} False.\newline
\textbf{Fact:} Dog meat provides protein comparable to other meats, with no demonstrated effect on strength, stamina, or virility. The trade carries documented infectious risks, including rabies exposure during transport and slaughter, cholera linkage, and trichinellosis, since dogs receive no veterinary inspection chain. Farming conditions concentrate disease spillover potential precisely when zoonotic vigilance matters globally. Any stamina attributed to tonic soups reflects calories, salt, warmth, and expectation. Nutritionally there is nothing here that poultry lacks. The costs are measured in suffering, outbreak risk, and zero pharmacology.
\end{enumerate}
\section{Southeast Asian Remedies}
\begin{enumerate}[leftmargin=0pt,labelwidth=0pt,labelsep=0pt,itemindent=0pt]
\item[] \textbf{ASM-014.} \textbf{Myth:} Tongkat ali root raises testosterone dramatically and cures male infertility.\newline
\textbf{Status:} False.\newline
\textbf{Fact:} A few small trials suggest possible modest effects on stress measures, testosterone levels, or semen parameters, but study sizes, durations, and quality fall far short of supporting cure claims \cite{ref150}. Products marketed across the region have been adulterated with sildenafil analogs, meaning some effects attributed to roots were pharmaceuticals all along. Dosing and standardization vary so widely that consumers cannot compare products. Infertility deserves endocrine evaluation, because treatable causes exist. Herbal energy drinks are beverages, not reproductive medicine. Hope bottled in bark extract should at least be honest about its contents.
\item[] \textbf{ASM-015.} \textbf{Myth:} Kratom is a safe natural herb for energy and pain that causes no addiction.\newline
\textbf{Status:} False.\newline
\textbf{Fact:} Kratom's active compounds bind opioid receptors, producing dose-dependent stimulation at low amounts and opioid-like analgesia and sedation at higher ones \cite{ref147}. Dependence, withdrawal syndromes resembling opioid withdrawal, and use disorder are documented, and regulators have recorded adverse events and deaths predominantly in combination with other substances. Natural origin changes nothing about receptor pharmacology. People using kratom for pain or opioid problems deserve medical assessment, not vending-machine mitragynine. Herbal does not mean harmless, and opioid-family herbs least of all. The label says supplement; the receptors say otherwise.
\item[] \textbf{ASM-016.} \textbf{Myth:} Snakehead fish soup accelerates surgical wound healing far beyond ordinary nutrition.\newline
\textbf{Status:} False.\newline
\textbf{Fact:} Snakehead fish is lean, protein-rich food, and adequate protein genuinely supports wound repair, which is the grain of truth beneath the legend. Claims of specific super-healing rest on preliminary arginine studies and tradition rather than comparative trials against equally protein-rich meals. Chicken breast, eggs, lentils, and fish generally provide the same substrates. Postoperative recovery depends on nutrition, glycemic control, smoking cessation, and wound care together. Enjoying the regional delicacy after surgery is perfectly sensible. Expecting miracles from one fish undervalues everything else healing requires.
\item[] \textbf{ASM-017.} \textbf{Myth:} Eating durian together with alcohol is a lethal combination.\newline
\textbf{Status:} Uncertain.\newline
\textbf{Fact:} Durian contains sulfur compounds that can inhibit aldehyde dehydrogenase, the enzyme clearing acetaldehyde, and a small controlled study showed measurable interference with alcohol metabolism. Whether this rises to danger is unresolved: verified deaths from the pairing are anecdotal and unconfirmed, while discomfort, flushing, nausea, and hangover intensification are plausibly real. Regional warnings likely arose from genuine bad nights amplified by storytelling. People taking disulfiram-like drugs or with alcohol metabolism disorders have better reasons than folklore to separate the two. Deadly appears overstated; unpleasant appears well supported. Moderation covers both possibilities.
\item[] \textbf{ASM-018.} \textbf{Myth:} Coconut water can replace intravenous fluids in emergencies.\newline
\textbf{Status:} False.\newline
\textbf{Fact:} Wartime anecdotes describe coconut water used briefly as improvised intravenous fluid when nothing else existed, and those stories became modern marketing mythology. Coconut water's electrolyte profile differs fundamentally from medical solutions, carrying high potassium that risks dangerous heart rhythm disturbances if infused. It is not sterile, and infusion through the shell invites infection. As an oral rehydration and sports drink it performs admirably. The intravenous legend persists because it sounds like resourceful genius. Real saline bags exist precisely because improvised ones occasionally stop hearts.
\item[] \textbf{ASM-019.} \textbf{Myth:} Bitter melon lowers glucose so effectively that diabetes medicines become unnecessary.\newline
\textbf{Status:} False.\newline
\textbf{Fact:} Small trials show bitter melon preparations produce mild glucose reductions that fall short of metformin, let alone insulin, and formulations vary too much for reliable dosing. Substituting it for prescribed therapy has contributed to diabetic ketoacidosis in reported cases. As a vegetable within a balanced diet it is a fine choice for anyone, including people with diabetes counting its carbohydrates. The bitterness that suggests potency to folklore has no relationship to glycemic power. Vegetables complement treatment plans; they do not graduate into them. Uncontrolled diabetes damages quietly while patients sip alternatives.
\end{enumerate}
\section{Betel and Daily Practices}
\begin{enumerate}[leftmargin=0pt,labelwidth=0pt,labelsep=0pt,itemindent=0pt]
\item[] \textbf{ASM-020.} \textbf{Myth:} Betel nut chewed without tobacco is harmless, aids digestion, and cleans the teeth.\newline
\textbf{Status:} False.\newline
\textbf{Fact:} Areca nut alone is classified as a human carcinogen, independent of tobacco, driving oral cancers and oral submucous fibrosis, a progressive scarring that freezes the jaw open \cite{ref148}. Hundreds of millions of chewers across South Asia, Southeast Asia, Taiwan, and Pacific communities give the habit enormous public-health reach. The mild euphoria and salivation users cite as digestive aid come from arecoline's cholinergic stimulation, not from dental hygiene. Fibrosis makes chewing, speaking, and eating progressively harder. Red-stained smiles advertise cumulative DNA damage. No tobacco-free version of a Group 1 carcinogen becomes safe by subtraction.
\item[] \textbf{ASM-021.} \textbf{Myth:} New mothers must avoid showering and stay intensely warm for a month after childbirth to prevent lifelong illness.\newline
\textbf{Status:} False.\newline
\textbf{Fact:} Confinement traditions offer valued rest and family support, but their extreme rules contradict postnatal care guidance emphasizing hygiene, gentle mobility, balanced nutrition, and attention to warning signs \cite{ref149}. Forbidding bathing increases infection risk at the very time lochia and perineal wounds require cleanliness. Overheating practices have caused maternal and newborn heatstroke deaths, particularly in summer. Rest, hydration, and help with the baby are the beneficial core worth keeping. The month-long prohibition of water belongs to pre-plumbing eras. Mothers recover faster clean, fed, and mobile, not sweltering.
\item[] \textbf{ASM-022.} \textbf{Myth:} Gua sha scraping releases toxins and illness from the body through the skin.\newline
\textbf{Status:} False.\newline
\textbf{Fact:} The red or purple marks left by coin-and-oil scraping are petechiae and bruising caused by ruptured capillaries, identical in origin to the marks from cupping. Sweat glands do not export toxins on demand, and no measurement has shown waste products leaving through scraped skin. Many recipients report warmth and relaxation afterward, consistent with local blood flow and attentive touch. Those sensations are real comfort, honestly available without diagnostic pretense. Bruise color charts correlating shades with organ sickness lack physiological basis. \emph{(See also ALT-022.)}
\item[] \textbf{ASM-023.} \textbf{Myth:} Drinking cold beverages weakens the body and causes disease, so warm water must be preferred at all times.\newline
\textbf{Status:} False.\newline
\textbf{Fact:} Ingested liquids reach body temperature within minutes through ordinary thermoregulation, and cold drinks cause no respiratory, digestive, or menstrual disease in controlled study. Preferences shaped by climate and tradition are legitimate personal choices, and warm water feels soothing during illness. The pathology attributed to cold ingestion belongs instead to viruses, allergens, and reflux triggers. Ice water may briefly worsen migraine in susceptible people and cramp sensitive guts, which is tolerance, not toxicity. Hydration science cares about volume, not thermometer readings. \emph{(See also IND-057.)}
\end{enumerate}

\chapter{Myths Created or Spread by Health Professionals}
Medicine's credibility rests on professionals policing themselves, and the stories in this chapter exist because they mostly did: journalists, editors, courts, licensing boards, and rival scientists exposed each one. This chapter is not an argument against doctors. It documents a narrower phenomenon, cases where credentialed practitioners created, endorsed, or profited from health claims the evidence did not support, usually because money flowed toward the claim. The documentation here comes from retractions, court judgments, government investigations, and peer-reviewed critiques rather than rumor. Two lessons recur. Expertise does not immunize against incentive, and every safeguard described in these entries, from peer review to licensure, worked precisely when someone insisted on evidence.
\section{Research Fraud and Retracted Findings}
\begin{enumerate}[leftmargin=0pt,labelwidth=0pt,labelsep=0pt,itemindent=0pt]
\item[] \textbf{PRO-001.} \textbf{Myth:} The MMR vaccine causes autism, and the doctor who discovered this was silenced for telling the truth.\newline
\textbf{Status:} False.\newline
\textbf{Fact:} The 1998 paper suggesting an MMR-autism link was investigated by journalist Brian Deer, whose reporting triggered a General Medical Council inquiry that found children underwent invasive procedures without ethical approval and data were misrepresented. The journal retracted the paper in 2010, and the British Medical Journal's editor concluded in 2011 that the work was not error but deliberate fraud \cite{ref151}. Before publication, its lead author received hundreds of thousands of pounds from lawyers preparing litigation against vaccine makers and held a patent on a competing measles product. His medical license was revoked. Silencing never occurred; exposure did.
\item[] \textbf{PRO-002.} \textbf{Myth:} Once a study passes peer review, its findings are reliable science that doctors can act on.\newline
\textbf{Status:} False.\newline
\textbf{Fact:} Peer review screens for obvious flaws; it cannot verify data, reproduce experiments, or detect fabricated statistics, which surface only through replication attempts and whistleblowers. Retractions have grown steadily as oversight improves, and analysis of citation records shows retracted papers continue to be cited as valid for years after withdrawal \cite{ref152}. Some retractions involve outright fabrication by credentialed researchers whose careers depended on publication counts. The practical rule is calibration, not cynicism: single studies are leads, treatment standards rest on replicated evidence, and ``peer reviewed'' describes a filter, not a guarantee.
\end{enumerate}
\section{Invented Diagnoses and Disease Expansion}
\begin{enumerate}[leftmargin=0pt,labelwidth=0pt,labelsep=0pt,itemindent=0pt]
\item[] \textbf{PRO-003.} \textbf{Myth:} Osteopenia is a disease that most people diagnosed with it need medication to treat.\newline
\textbf{Status:} False.\newline
\textbf{Fact:} Osteopenia is not a disease but a statistical zone, bone density between two arbitrary cutoffs chosen partly at industry-supported consensus meetings as scanning machines multiplied in the 1990s. Most postmenopausal women labeled osteopenic will never fracture, and reviewers have argued that wholesale conversion of measurement into diagnosis drives overtreatment \cite{ref153}. Commentators call this pattern selling sickness, the expansion of diagnostic boundaries until healthy people become patients \cite{ref154}. Genuine high fracture risk deserves evaluation and often treatment. A number slightly below average is not itself an illness, and scanners profit from saying otherwise.
\item[] \textbf{PRO-004.} \textbf{Myth:} An epidemic of male menopause means millions of middle-aged men need lifelong testosterone gels.\newline
\textbf{Status:} False.\newline
\textbf{Fact:} Testosterone declines gradually with age in most men, and fatigue, low mood, and reduced libido have many common causes that advertising recast as one deficiency. Analysts documented how branded ``Low T'' awareness campaigns taught men to self-diagnose from symptom checklists and demand prescriptions \cite{ref155}. Cardiovascular safety signals later prompted regulators to require label warnings and restrict marketing. True hypogonadism exists, is diagnosed by repeated morning blood tests, and deserves treatment. A questionnaire on a manufacturer-funded website is not endocrinology, and lifelong hormone exposure is not a lifestyle accessory.
\item[] \textbf{PRO-005.} \textbf{Myth:} Chronic stress exhausts the adrenal glands, producing ``adrenal fatigue'' that saliva tests confirm and glandular supplements fix.\newline
\textbf{Status:} False.\newline
\textbf{Fact:} A systematic review set out to test the concept and found no consistent evidence that the condition exists as described \cite{ref156}. Endocrinology recognizes adrenal insufficiency, a rare and dangerous failure diagnosable by blood testing, and it looks nothing like the tiredness attributed to fatigue. Direct-to-consumer salivary cortisol panels interpret unvalidated reference ranges, and the same clinics sell the supplements that treat their diagnoses. Real causes of fatigue, sleep apnea, depression, anemia, thyroid disease, respond to actual medicine. Glands worn out by modern life make intuitive fiction, not physiology.
\item[] \textbf{PRO-006.} \textbf{Myth:} A persistently low body temperature proves hidden thyroid dysfunction that only specialized timed-release T3 therapy can correct.\newline
\textbf{Status:} False.\newline
\textbf{Fact:} ``Wilson's temperature syndrome'' was named for and trademarked by the physician who sells certification seminars and proprietary dosing protocols for it, a business model visible in the branding. Major thyroid organizations state that the syndrome lacks scientific support, that body temperature varies normally across people and hours, and that thyroid hormone should be prescribed for tested disease, not thermometer readings \cite{ref157}. Driving metabolism with unmonitored T3 risks palpitations, tremor, bone loss, and arrhythmia. Hypothyroidism is real and easily confirmed by blood tests. Trademarks are for protecting markets, not validating mechanisms.
\item[] \textbf{PRO-007.} \textbf{Myth:} Vague symptoms from bloating to brain fog signal systemic yeast overgrowth that stool tests reveal and antifungal courses plus restrictive diets cure.\newline
\textbf{Status:} False.\newline
\textbf{Fact:} The concept was introduced by a physician in the orthomolecular literature of the 1980s and spread through practitioner networks despite rejection by gastroenterology and infectious disease societies. Commercial stool and blood ``candida'' panels use thresholds with no validation, diagnosing a condition whose serious medical namesake, invasive candidiasis, is a bloodstream infection of hospitalized patients, not a wellness complaint. Elimination diets and repeat antifungal courses generate clinic revenue while genuine causes of symptoms go unevaluated \cite{ref158}. Yeast exists; the profitable version of it lives mainly on lab reports.
\item[] \textbf{PRO-008.} \textbf{Myth:} ``Leaky gut syndrome'' is an established diagnosis behind most chronic disease, confirmed by commercial permeability panels.\newline
\textbf{Status:} False.\newline
\textbf{Fact:} Intestinal permeability is a genuine, measurable phenomenon that research links to conditions including celiac disease and critical illness, and increased permeability may contribute to some disease processes \cite{ref159}. What lacks support is the syndromized marketing version: a standalone diagnosis covering dozens of complaints, detected by panels with unvalidated interpretation, and remedied by supplement protocols sold alongside the test. Gastroenterologists diagnose the diseases where permeability matters using standard criteria. Patients deserve the real science without the subscription. A mechanism in a journal is not yet a diagnosis in a clinic, whatever the checkout page promises.
\end{enumerate}
\section{Profitable Fears in Dentistry}
\begin{enumerate}[leftmargin=0pt,labelwidth=0pt,labelsep=0pt,itemindent=0pt]
\item[] \textbf{PRO-009.} \textbf{Myth:} Dentistry's focal infection era shows the profession protecting patients from hidden poisons.\newline
\textbf{Status:} False.\newline
\textbf{Fact:} In the early twentieth century, the theory that infected tooth sites caused arthritis, heart disease, and mental illness led dentists and physicians to extract millions of teeth, many perfectly restorable, sometimes clearing entire mouths on suspicion alone. Edentulism surged before controlled evidence collapsed the doctrine by mid-century. The episode is now studied as a cautionary case of plausible mechanism outrunning proof, and its logic resurfaces wherever clinics revive it to sell full-mouth extractions and implants \cite{ref160}. Historical confidence harmed real patients on an industrial scale. Tradition is not evidence of safety, even in retrospect.
\item[] \textbf{PRO-010.} \textbf{Myth:} Root-canalled teeth breed toxins that cause cancer, arthritis, and chronic fatigue, so they should be extracted and jaw cavitations surgically cleaned.\newline
\textbf{Status:} False.\newline
\textbf{Fact:} Large epidemiological studies find no such associations, and endodontic organizations reviewing the evidence affirm that properly treated root canals do not cause systemic disease \cite{ref160}. The alternative protocol, extraction plus ``cavitation'' curettage, ozone, and implant packages, bills thousands per mouth while removing functional teeth. Modern endodontics seals and sterilizes canals with documented success rates. Patients frightened by photographs of extracted teeth are buying surgery, not health. Keeping a natural tooth remains the evidence-based default. Fear of decay should not be redirected into fear of its successful treatment.
\item[] \textbf{PRO-011.} \textbf{Myth:} Silver amalgam fillings continuously poison the body, so everyone should have them drilled out and replaced.\newline
\textbf{Status:} False.\newline
\textbf{Fact:} Regulators reviewing decades of data conclude amalgam is safe for the general population at typical exposure, releasing trace vapor well below thresholds of harm \cite{ref161}. Mass removal is the one scenario that measurably increases exposure, since drilling heats fillings and briefly raises inhaled vapor. Replacement generates fees many times those of watchful waiting, which explains part of the enthusiasm among ``biological'' practices. A small allergy subset legitimately needs alternatives, and alternatives exist. For everyone else, the filling that has quietly served for twenty years outperforms the invoice that removes it.
\end{enumerate}
\section{The Opioid Era's Teachings}
\begin{enumerate}[leftmargin=0pt,labelwidth=0pt,labelsep=0pt,itemindent=0pt]
\item[] \textbf{PRO-012.} \textbf{Myth:} Pain is a vital sign, and compassionate medicine requires scoring it toward zero, chiefly with opioids.\newline
\textbf{Status:} False.\newline
\textbf{Fact:} Vital signs measure physiology objectively; pain scales record a subjective report. The ``fifth vital sign'' framing, promoted through advocacy groups and educational campaigns with manufacturer funding, helped embed pain scoring into accreditation standards in 2001, reshaping prescribing expectations across hospitals \cite{ref162}. Documented investigations of opioid marketing describe how the teaching eased prescriber reluctance while dependence data were minimized. Pain relief matters and multidisciplinary care achieves it; a numerical target tied to escalation is not care but throughput. The slogan outlived its sponsors' honesty.
\item[] \textbf{PRO-013.} \textbf{Myth:} When opioid patients show drug-seeking behavior, it proves undertreated pain, so doses should be escalated.\newline
\textbf{Status:} False.\newline
\textbf{Fact:} Pseudoaddiction originated in a 1989 palliative-care observation about terminally ill cancer patients, where resolving pain resolved the behaviors \cite{ref163}. Industry training materials generalized the concept far beyond hospice, teaching community prescribers to read warning signs of dependence as evidence for higher doses \cite{ref162}. Applied broadly, the heuristic instructed clinicians to override exactly the signals that deserved scrutiny. The original authors described an iatrogenic syndrome; its commercial second life made dependence genuinely iatrogenic at scale. Context transforms the same observation from reassurance into red flag.
\item[] \textbf{PRO-014.} \textbf{Myth:} Fewer than one percent of patients treated with narcotics become addicted, so long-term opioid therapy carries negligible risk.\newline
\textbf{Status:} False.\newline
\textbf{Fact:} The statistic traces to a single five-sentence 1980 letter to a medical journal describing hospitalized inpatients receiving short courses, mostly single doses \cite{ref164}. Over three decades it accumulated hundreds of scholarly citations as if it were a rigorous study, appearing in textbooks and promotional material. Later analyses traced the distortion, noting the letter never examined outpatient chronic therapy, where dependence risk is substantially higher, and its own authors expressed alarm at the citation chain. One paragraph, amplified by incentive, shaped a generation of prescribing. Citation counts measure circulation, not validity.
\end{enumerate}
\section{Clinic Economies and Unproven Protocols}
\begin{enumerate}[leftmargin=0pt,labelwidth=0pt,labelsep=0pt,itemindent=0pt]
\item[] \textbf{PRO-015.} \textbf{Myth:} Lingering symptoms after Lyme disease mean persistent infection requiring months or years of intravenous antibiotics.\newline
\textbf{Status:} False.\newline
\textbf{Fact:} Post-treatment symptoms are real, but randomized trials repeatedly show prolonged antibiotic courses provide no lasting benefit while causing line infections, Clostridioides difficile colitis, and drug reactions, which is why current guidelines advise against them \cite{ref165}. Cash-pay ``Lyme literate'' clinics bill extended regimens unsupported by microbiology, and court records include judgments against physicians who billed insurers for such care. Fatigue and joint pain after any infection deserve evaluation, not reflexive antimicrobials. The spirochete is fragile; the billing model is durable.
\item[] \textbf{PRO-016.} \textbf{Myth:} Physician-administered stem cell injections can already rebuild joints, organs, and youthfulness.\newline
\textbf{Status:} False.\newline
\textbf{Fact:} Stem cell medicine has genuine triumphs, bone marrow transplantation chief among them, and promising orthopedic trials underway. What it does not have is evidence that fat or marrow injections reverse osteoarthritis, neurological disease, or aging, which has not stopped physician-affiliated clinics from charging thousands per injection. Federal regulators have issued warning letters and pursued enforcement against clinics processing cells in ways that create unapproved drugs with contamination and misprocessing risks \cite{ref166}. Paying cash for one's own cells does not convert experiment into therapy. Hope is the product; the syringe is packaging.
\item[] \textbf{PRO-017.} \textbf{Myth:} Autism is heavy metal poisoning that chelation therapy can remove.\newline
\textbf{Status:} False.\newline
\textbf{Fact:} A systematic review of the protocol found no evidence of benefit for autistic features \cite{ref167}, while chelation drugs carry documented dangers including fatal hypocalcemia when administered incorrectly. The hypothesis linking vaccines or environment to removable metals failed specific testing, yet conference networks marketed multi-thousand-dollar chelation packages to families, and licensing authorities have disciplined physicians for the practice. Autistic children deserve developmental support, speech therapy, and family education. Infusing a child with a metal-binding drug on behalf of a disproven theory is not medicine taking chances; it is commerce taking advantage.
\item[] \textbf{PRO-018.} \textbf{Myth:} Telehealth physicians selling ivermectin subscription plans proved the drug treats COVID-19.\newline
\textbf{Status:} False.\newline
\textbf{Fact:} Randomized controlled trials, including large community-based studies, found no meaningful benefit of ivermectin against COVID-19, and health regulators advised against its use outside trials \cite{ref168}. During the pandemic, some ventures charged recurring consultation fees specifically to prescribe it, shipping veterinary-strength formulations with dosing errors, and medical boards disciplined licensed prescribers for substandard care. A prescription pad is not evidence, and willingness to charge for a tablet measures neither efficacy nor courage. Antiparasitic fame faded; the subscription revenue model, unfortunately, travels.
\item[] \textbf{PRO-019.} \textbf{Myth:} Annual whole-body CT scans let healthy people catch disease early and buy peace of mind.\newline
\textbf{Status:} False.\newline
\textbf{Fact:} Screening asymptomatic populations with whole-body CT exposes them to meaningful radiation while generating incidentalomas, harmless findings that trigger further imaging, biopsies, surgeries, and anxiety, with net harm demonstrated in assessments of the technology \cite{ref169}. Professional radiology organizations oppose the practice, yet physician-owned imaging centers marketed it directly to consumers because the scanner pays for itself one worried customer at a time. Targeted, guideline-based screening saves lives; panoramic scanning of the healthy manufactures patients. Ownership of the machine should never substitute for indication for the scan.
\end{enumerate}
\section{Surgical Overuse and Cash-Pay Practices}
\begin{enumerate}[leftmargin=0pt,labelwidth=0pt,labelsep=0pt,itemindent=0pt]
\item[] \textbf{PRO-020.} \textbf{Myth:} Hysterectomy is automatically the right answer for fibroids or heavy bleeding once childbearing ends.\newline
\textbf{Status:} False.\newline
\textbf{Fact:} When a hospital system audited its records, researchers found a substantial share of women undergoing hysterectomy for benign conditions had no documented trial of alternatives such as hormonal management or uterine artery embolization first \cite{ref170}. Surgery is sometimes clearly correct, and shared decision-making with a second opinion meaningfully changes plans. Organs removed cannot be reconsidered. The operation's routine familiarity should raise vigilance rather than lower it, because convenience for the schedule is not benefit for the patient. Ask what was tried before asking what will be taken.
\item[] \textbf{PRO-021.} \textbf{Myth:} Planned cesarean delivery without medical indication is just as safe as labor and simply more convenient.\newline
\textbf{Status:} False.\newline
\textbf{Fact:} The World Health Organization's assessment found that population cesarean rates beyond roughly ten to fifteen percent show no associated mortality benefit, while maternal risks, hemorrhage, infection, and abnormal placentation in future pregnancies, accumulate with each procedure, and newborns face higher rates of transient breathing difficulty \cite{ref171}. Documented drivers of non-medical cesareans include scheduling predictability and reimbursement patterns, benefits that accrue to calendars and balance sheets. Cesarean birth is indispensable when indicated. Convenience is not an indication, whichever side of the scalpel it favors.
\item[] \textbf{PRO-022.} \textbf{Myth:} Spinal fusion is the definitive cure for chronic nonspecific back pain.\newline
\textbf{Status:} False.\newline
\textbf{Fact:} Fusion volumes climbed steeply while regional variation widened far beyond any explanation by patient mix, a signature of supply-sensitive demand \cite{ref172}. Trials in mixed chronic pain populations show fusion generally performs no better than structured conservative care, and analyses have linked surgeon industry payments to higher operative rates, though fusion remains genuinely appropriate for instability, deformity, and specific structural disease. Back pain is usually managed, not machined away. Hardware fixes anatomy that may not be the pain source. The screw's warranty is longer than the evidence supporting its insertion.
\item[] \textbf{PRO-023.} \textbf{Myth:} Intravenous vitamin cocktails at clinician-run drip bars detoxify, energize, and immunity-proof healthy customers.\newline
\textbf{Status:} False.\newline
\textbf{Fact:} Healthy people absorb oral vitamins efficiently, and water-soluble excess simply leaves in urine, so infusion adds expense rather than benefit. What infusion adds clinically is catheter-related infection risk, thrombosis, fluid overload in susceptible hearts, and allergic reaction, hazards familiar to hospitals that reserve lines for necessity. The wellness drip model converts sterile technique into a premium priced beverage service, staffed by credentialed hands that lend legitimacy the evidence cannot. Rehydration therapy belongs to dehydration. \emph{(See also GEN-060.)}
\end{enumerate}
\backmatter
\begingroup
\sloppy
\raggedright
\renewcommand{\bibname}{Bibliography}
\phantomsection
\addcontentsline{toc}{chapter}{Bibliography}
\begin{thebibliography}{999}
\bibitem{ref1}Madsen KM, et al. A population-based study of measles, mumps, and rubella vaccination and autism. New England Journal of Medicine. 2002;347:1477--1482. DOI: \url{https://doi.org/10.1056/NEJMoa021134}.
\bibitem{ref2}Taylor LE, Swerdfeger AL, Eslick GD. Vaccines are not associated with autism: an evidence-based meta-analysis of case-control and cohort studies. Vaccine. 2014;32:3623-- 3629. DOI: \url{https://doi.org/10.1016/j.vaccine.2014.04.085}.
\bibitem{ref3}World Health Organization. Antimicrobial resistance: fact sheet. Updated 2023. \url{https://www.who.int/news-room/fact-sheets/detail/antimicrobial-resistance}.
\bibitem{ref4}Centers for Disease Control and Prevention. About common cold. Updated 2026. \url{https://www.cdc.gov/common-cold/about/index.html}.
\bibitem{ref5}National Center for Complementary and Integrative Health. ``Detoxes'' and ``cleanses'': what you need to know. \url{https://www.nccih.nih.gov/health/detoxes-and-cleanses-what-you-need-to-know}.
\bibitem{ref6}U.S. Food and Drug Administration. Microwave ovens. \url{https://www.fda.gov/radiation-emitting-products/resources-you-radiation-emitting-products/microwave-ovens}.
\bibitem{ref7}Wolraich ML, Wilson DB, White JW. The effect of sugar on behavior or cognition in children: a meta-analysis. JAMA. 1995;274:1617--1621. PMID: \url{https://pubmed.ncbi.nlm.nih.gov/7474248/}.
\bibitem{ref10}American Heart Association. Here's the latest on dietary cholesterol and how it fits in with a healthy diet. 2023. AHA article.
\bibitem{ref13}University of Utah Health. The top 10 women's health myths. Published 2021. \url{https://healthcare.utah.edu/womens-health/patient-education}.
\bibitem{ref16}National Cancer Institute. Common cancer myths and misconceptions. \url{https://www.cancer.gov/about-cancer/causes-prevention/risk/myths}.
\bibitem{ref18}WebMD. 10 health myths debunked. Medically reviewed 2026. \url{https://www.webmd.com/balance/ss/slideshow-10-health-myths-debunked}.
\bibitem{ref21}Mayo Clinic Press. Myths and facts about liver health. 2025. \url{https://mcpress.mayoclinic.org/living-well/myths-and-facts-about-liver-health/}.
\bibitem{ref22}Mayo Clinic Press. Common myths and misconceptions about alcohol use. 2024. Mayo Clinic article.
\bibitem{ref23}WebMD. 5 common myths about pain and pain relief. Medically reviewed 2026. \url{https://www.webmd.com/pain-management/5-common-myths-about-pain-and-pain-relief}.
\bibitem{ref24}Mayo Clinic. Organ donation: don't let these myths confuse you. 2026. \url{https://www.mayoclinic.org/healthy-lifestyle/consumer-health/in-depth/organ-donation/art-20047529}.
\bibitem{ref25}Kramer MS. Believe It or Not: The History, Culture, and Science Behind Health Beliefs and Practices. Springer; 2023. \url{https://link.springer.com/book/10.1007/978-3-031-46022-7}.
\bibitem{ref27}Mayo Clinic. Hangovers: diagnosis and treatment. Updated 2024. \url{https://www.mayoclinic.org/diseases-conditions/hangovers/diagnosis-treatment/drc-20373015}.
\bibitem{ref28}U.S. Food and Drug Administration. Aspartame and other sweeteners in food. Updated 2024. \url{https://www.fda.gov/food/food-additives-petitions/aspartame-and-other-sweeteners-food}.
\bibitem{ref30}Yasmin S. Viral BS: Medical Myths and Why We Fall for Them. Johns Hopkins University Press; 2021. \url{https://search.worldcat.org/title/1153340618}.
\bibitem{ref37}Mayo Clinic. Myths about catching a cold and cold symptoms. \url{https://newsnetwork.mayoclinic.org/discussion/mayo-clinic-q-and-a-myths-about-catching-a-cold/}.
\bibitem{ref40}Centers for Disease Control and Prevention. Misconceptions about seasonal flu and flu vaccines. Updated 2024. \url{https://www.cdc.gov/flu/prevention/misconceptions.html}.
\bibitem{ref41}Centers for Disease Control and Prevention. Myths and facts about COVID-19 vaccines. \url{https://stacks.cdc.gov/view/cdc/110056}.
\bibitem{ref42}U.S. Food and Drug Administration. Generic drugs: overview and basics. Updated 2025. \url{https://www.fda.gov/drugs/generic-drugs/overview-basics}.
\bibitem{ref43}U.S. Food and Drug Administration. Acetaminophen: safe use and overdose risks. Updated 2026. \url{https://www.fda.gov/drugs/safe-use-over-counter-pain-relievers-and-fever-reducers/acetaminophen}.
\bibitem{ref44}American Dental Association. Toothbrushes. \url{https://www.ada.org/resources/ada-library/oral-health-topics/toothbrushes}.
\bibitem{ref45}American Dental Association. Toothpastes. Updated 2025. \url{https://www.ada.org/resources/ada-library/oral-health-topics/toothpastes}.
\bibitem{ref46}American Heart Association. Top 10 myths about cardiovascular disease. Reviewed 2024. \url{https://www.heart.org/en/health-topics/consumer-healthcare/what-is-cardiovascular-disease/top-10-myths-about-cardiovascular-disease}.
\bibitem{ref47}American Dental Association. Online dental misinformation. Dental Sound Bites, 2025. \url{https://www.ada.org/publications/dental-sound-bites/season-6/online-dental-misinformation}.
\bibitem{ref48}U.S. Food and Drug Administration. Fraudulent products and health-fraud warnings. Updated 2025. \url{https://www.fda.gov/drugs/medication-health-fraud/fraudulent-products}.
\bibitem{ref49}American Heart Association. Cholesterol: myths vs. facts. \url{https://recipes.heart.org/-/media/Files/Health-Topics/Cholesterol/Cholesterol-Myths-vs-Facts-English.pdf}.
\bibitem{ref50}American Dental Association. Debunking dental trends. 2025. \url{https://adanews.ada.org/ada-news/2025/january/debunking-dental-trends/}.
\bibitem{ref51}World Health Organization Regional Office for Africa. COVID-19 mythbusters. \url{https://www.afro.who.int/health-topics/coronavirus-covid-19/mythbusters}.
\bibitem{ref52}U.S. Preventive Services Task Force. Aspirin use to prevent cardiovascular disease: preventive medication. Updated 2022. \url{https://www.uspreventiveservicestaskforce.org/uspstf/recommendation/aspirin-to-prevent-cardiovascular-disease-preventive-medication}.
\bibitem{ref53}American Heart Association. Statin safety and associated adverse events. 2018. \url{https://professional.heart.org/en/science-news/statin-safety-and-associated-adverse-events/top-things-to-know}.
\bibitem{ref54}Centers for Disease Control and Prevention. Thimerosal and vaccines. Updated 2025. \url{https://www.cdc.gov/vaccine-safety/about/thimerosal.html}.
\bibitem{ref55}U.S. Food and Drug Administration. Questions and answers on monosodium glutamate (MSG). \url{https://www.fda.gov/food/food-additives-petitions/questions-and-answers-monosodium-glutamate-msg}.
\bibitem{ref56}U.S. Food and Drug Administration. Food irradiation: what you need to know. Updated 2024. \url{https://www.fda.gov/food/buy-store-serve-safe-food/food-irradiation-what-you-need-know}.
\bibitem{ref57}U.S. Food and Drug Administration. Agricultural biotechnology. \url{https://www.fda.gov/food/consumers/agricultural-biotechnology}.
\bibitem{ref58}World Health Organization Regional Office for Europe. Myth-busters: debunking long COVID myths and misconceptions. 2026. \url{https://www.who.int/europe/news-room/events/item/2026/01/30/default-calendar/myth-busters--debunking-long-covid-myths-and-misconceptions}.
\bibitem{ref59}Andersen CH. ``55 Health Myths That Need to Die.'' Reader's Digest, republished by The Centers for Advanced Orthopaedics. Published 2017. \url{https://www.cfaortho.com/media/news/2017/02/55-health-myths-that-need-to-die-reader-s-digest}.
\bibitem{ref60}Koutsky J. The biggest health myths of all time. Medically reviewed by Ashley Maltz, M.D. HealthCentral. Updated 2021. \url{https://www.healthcentral.com/article/top-health-myths-of-all-time}.
\bibitem{ref62}Centers for Disease Control and Prevention. Concussion questions and answers: what parents need to know. Updated 2024. \url{https://www.cdc.gov/heads-up/media/pdfs/2024/07/concussion_qa-508.pdf}.
\bibitem{ref63}Centers for Disease Control and Prevention. Symptoms of depression among women. Updated 2024. \url{https://www.cdc.gov/reproductive-health/depression/index.html}.
\bibitem{ref64}Centers for Disease Control and Prevention. Risk and protective factors for suicide. Updated 2026. \url{https://www.cdc.gov/suicide/risk-factors/}.
\bibitem{ref65}National Health Service. Postmenopausal bleeding. Updated 2023. \url{https://www.nhs.uk/symptoms/post-menopausal-bleeding/}.
\bibitem{ref8}Sullivan JE, Farrar HC. Fever and antipyretic use in children. Pediatrics. 2011;127:580--587. DOI: \url{https://doi.org/10.1542/peds.2010-3852}.
\bibitem{ref9}Centers for Disease Control and Prevention. Human papillomavirus vaccine safety. \url{https://www.cdc.gov/vaccine-safety/vaccines/hpv.html}.
\bibitem{ref11}University of Arkansas for Medical Sciences Health. Medical myths. \url{https://uamshealth.com/medical-myths/}.
\bibitem{ref12}Osmosis. Top 10 medical myths: what your patients should know. \url{https://www.osmosis.org/blog/top-10-medical-myths-what-your-patients-should-know}.
\bibitem{ref14}Mount Elizabeth Hospitals. 10 more common health myths and misconceptions debunked. Updated 2022. \url{https://www.mountelizabeth.com.sg/health-plus/article/10-health-myths-debunked}.
\bibitem{ref15}Max Healthcare. Cancer myths men believe: hidden dangers and health risks. Published 2025. \url{https://www.maxhealthcare.in/blogs/cancer-myths-men-believe}.
\bibitem{ref17}Healthline. Common health myths debunked. Updated 2026. \url{https://www.healthline.com/health/common-health-myths}.
\bibitem{ref19}MedicineNet. Are 2 eggs a day healthy? Medically reviewed 2026. \url{https://www.medicinenet.com/are_2_eggs_a_day_healthy/article.htm}.
\bibitem{ref20}National Institutes of Health, Office of Dietary Supplements. Multivitamin/mineral supplements: consumer fact sheet. \url{https://ods.od.nih.gov/factsheets/MVMS-Consumer/}.
\bibitem{ref26}American Academy of Dermatology. Warts: dermatologists' tips for at-home treatment. \url{https://www.aad.org/public/diseases/a-z/warts-self-care}.
\bibitem{ref29}National Cancer Institute. Artificial sweeteners and cancer. \url{https://www.cancer.gov/about-cancer/causes-prevention/risk/diet/artificial-sweeteners-fact-sheet}.
\bibitem{ref31}U.S. Food and Drug Administration. Talc. Updated 2025. \url{https://www.fda.gov/cosmetics/cosmetic-ingredients/talc}.
\bibitem{ref32}American Cancer Society. Talcum powder and cancer. \url{https://www.cancer.org/cancer/risk-prevention/chemicals/talcum-powder-and-cancer.html}.
\bibitem{ref33}World Health Organization. Contraceptive use: a catalyst for women's health and socioeconomic empowerment. 2025. WHO evidence brief.
\bibitem{ref34}Centers for Disease Control and Prevention. Guidelines for reporting a suicide cluster. Updated 2026. \url{https://www.cdc.gov/suicide/prevention/cluster.html}.
\bibitem{ref35}Endocrine Society. Testosterone therapy in men with hypogonadism: guideline resources. \url{https://www.endocrine.org/clinical-practice-guidelines/testosterone-therapy}.
\bibitem{ref36}Gleneagles Hospital. 10 popular health myths debunked. \url{https://www.gleneagles.com.sg/health-plus/article/12-popular-health-myths-debunked}.
\bibitem{ref38}Cleveland Clinic. Can cranberry stop your UTIs? \url{https://health.clevelandclinic.org/can-cranberry-juice-stop-uti}.
\bibitem{ref39}Cleveland Clinic. Is urine sterile? The truth behind the myth. \url{https://health.clevelandclinic.org/is-urine-sterile}.
\bibitem{ref61}Rowan K, Nixon Pompa R. 25 medical myths that just won't go away. Live Science. Published 2016. \url{https://www.livescience.com/36100-10-medical-myths.html}.
\bibitem{ref66}Centers for Disease Control and Prevention. What vaccines are recommended for you: adult vaccines. Updated 2025. \url{https://www.cdc.gov/vaccines-adults/recommended-vaccines/index.html}.
\bibitem{ref67}World Health Organization. Menstrual health. Updated 2026. \url{https://www.who.int/news-room/fact-sheets/detail/menstrual-health}.
\bibitem{ref68}World Health Organization. Support for mothers to initiate and establish breastfeeding after childbirth. Updated 2023. \url{https://www.who.int/tools/elena/interventions/breastfeeding-support}.
\bibitem{ref69}Centers for Disease Control and Prevention. Foods and drinks to avoid or limit: honey before 12 months. Updated 2026. \url{https://www.cdc.gov/infant-toddler-nutrition/foods-and-drinks/foods-and-drinks-to-avoid-or-limit.html}.
\bibitem{ref70}World Health Organization. Caring for a newborn. Updated 2022. \url{https://www.who.int/tools/your-life-your-health/life-phase/newborns-and-children-under-5-years/caring-for-newborns}.
\bibitem{ref71}World Health Organization. Epilepsy. Updated 2024. \url{https://www.who.int/news-room/fact-sheets/detail/epilepsy}.
\bibitem{ref72}World Health Organization. Snakebite envenoming. Updated 2023. \url{https://www.who.int/news-room/fact-sheets/detail/snakebite-envenoming}.
\bibitem{ref73}World Health Organization. Animal bites. Updated 2024. \url{https://www.who.int/news-room/fact-sheets/detail/animal-bites}.
\bibitem{ref74}World Health Organization. Leprosy. Updated 2026. \url{https://www.who.int/news-room/fact-sheets/detail/leprosy}.
\bibitem{ref75}World Health Organization. Measles. Updated 2026. \url{https://www.who.int/en/news-room/fact-sheets/detail/measles}.
\bibitem{ref76}Centers for Disease Control and Prevention. About chickenpox. Updated 2024. \url{https://www.cdc.gov/chickenpox/about/index.html}.
\bibitem{ref77}Centers for Disease Control and Prevention. Clinical signs and symptoms of smallpox. Updated 2024. \url{https://www.cdc.gov/smallpox/hcp/clinical-signs/index.html}.
\bibitem{ref78}World Health Organization. About hepatitis. \url{https://www.who.int/health-topics/hepatitis/about-hepatitis}.
\bibitem{ref79}U.S. Food and Drug Administration. Kohl, kajal, al-kahal, surma, tiro, tozali, or kwalli: by any name, beware of lead poisoning. \url{https://www.fda.gov/cosmetics/cosmetic-products/kohl-kajal-al-kahal-surma-tiro-tozali-or-kwalli-any-name-beware-lead-poisoning}.
\bibitem{ref80}National Center for Complementary and Integrative Health. Ayurvedic medicine: in depth. \url{https://www.nccih.nih.gov/health/ayurvedic-medicine-in-depth}.
\bibitem{ref81}National Center for Complementary and Integrative Health. Diabetes and dietary supplements: what you need to know. \url{https://www.nccih.nih.gov/health/diabetes-and-dietary-supplements-what-you-need-to-know}.
\bibitem{ref82}American Academy of Dermatology. Vitiligo: FAQs. \url{https://www.aad.org/public/diseases/a-z/vitiligo-overview}.

\bibitem{ref83}Federal Trade Commission. FTC charges marketers of Kinoki foot pads with deceptive advertising; seeks funds for consumer redress. January 28, 2009. \url{https://www.ftc.gov/news-events/news/press-releases/2009/01/ftc-charges-marketers-kinoki-foot-pads-deceptive-advertising-seeks-funds-consumer-redress}.
\bibitem{ref84}Federal Trade Commission. At FTC's request, judge imposes ban on marketers of detox foot pads. November 4, 2010. \url{https://www.ftc.gov/news-events/news/press-releases/2010/11/ftcs-request-judge-imposes-ban-marketers-detox-foot-pads}.
\bibitem{ref85}Center for Science in the Public Interest. Airborne false advertising class settlement. 2008.
\bibitem{ref86}American Dental Association. Mouthwash. MouthHealthy consumer resource.
\bibitem{ref87}Federal Trade Commission. Kellogg settles FTC charges that advertising for Frosted Mini-Wheats was false. April 20, 2009.
\bibitem{ref88}Center for Science in the Public Interest. Kellogg drops immunity claims for Cocoa Krispies. 2009.
\bibitem{ref89}Federal Trade Commission v. POM Wonderful, LLC, 777 F.3d 478 (D.C. Cir. 2015).
\bibitem{ref90}European Food Safety Authority. Scientific opinions on health claims related to probiotic dairy strains. EFSA Journal, 2010.
\bibitem{ref91}In re The Dannon Company, Inc., Activia and DanActive Litigation, class settlement. 2010.
\bibitem{ref92}Federal Trade Commission v. QRay.com, Inc., civil contempt judgment and affirmance. 2006--2007.
\bibitem{ref93}Power Balance Technologies. Statement acknowledging absence of credible scientific support for hologram claims; consumer refund program. 2010--2011.
\bibitem{ref94}Federal Trade Commission. Skechers to pay 40 million dollars for deceptive advertising of toning shoes. September 2012.
\bibitem{ref95}Federal Trade Commission. Reebok to pay 25 million dollars in customer refunds over easytone toning shoe claims. September 2011.
\bibitem{ref96}U.S. Food and Drug Administration. Public health advisory: Zicam Cold Remedy nasal gel, Zicam Cold Remedy Swabs, and loss of sense of smell. June 16, 2009. \url{http://www.fda.gov/Drugs/DrugSafety/PublicHealthAdvisories/ucm166059.htm}.
\bibitem{ref97}U.S. Food and Drug Administration. Warning letters to distributors of essential oils regarding unlawful health claims. September 2014.
\bibitem{ref98}Australian Government National Health and Medical Research Council. Statement on homeopathy. 2015.
\bibitem{ref99}Boiron USA. Class action settlement regarding Oscillococcinum claims. 2012.
\bibitem{ref100}Stanford Research into the Impact of Tobacco Advertising. ``More doctors smoke Camels than any other cigarette'' campaign materials, 1946--1952. \url{http://tobacco.stanford.edu}.
\bibitem{ref101}National Cancer Institute. Risks Associated with Smoking Cigarettes with Low Machine-Measured Yields of Tar and Nicotine. Smoking and Tobacco Control Monograph 13. NIH Publication No. 02-5074; 2001.
\bibitem{ref102}Thun MJ, Burns DM. Health impact of ``reduced yield'' cigarettes: a critical assessment of the epidemiological evidence. Tobacco Control. 2001;10(Suppl 1):i4--i11.
\bibitem{ref103}United States v. R.J. Reynolds Tobacco Co., stipulated order regarding ``natural'' and ``additive-free'' claims. December 2017.
\bibitem{ref104}El Ghissassi F, Baan R, Straif K, et al. A review of human carcinogens--part D: radiation. The Lancet Oncology. 2009;10(8):751--752. DOI: \url{https://doi.org/10.1016/S1470-2045(09)70213-X}.
\bibitem{ref105}Federal Trade Commission. Indoor Tanning Association settles FTC charges regarding health claims. March 2010.
\bibitem{ref106}Center for Science in the Public Interest. Coca-Cola fortified beverage labeling litigation and settlement. 2016.
\bibitem{ref107}Kearns CE, Schmidt LA, Glantz SA. Sugar industry and coronary heart disease research: a historical analysis of internal industry documents. JAMA Internal Medicine. 2016;176(11):1680--1685. DOI: \url{https://doi.org/10.1001/jamainternmed.2016.5394}.
\bibitem{ref108}Tye L. The Father of Spin: Edward L. Bernays and the Birth of Public Relations. Crown Publishers; 1998.
\bibitem{ref109}Heyman MB, Abrams SA; Section on Gastroenterology, Hepatology, and Nutrition; Committee on Nutrition. Fruit juice in infants, children, and adolescents: current recommendations. Pediatrics. 2017;139(6):e20170967. DOI: \url{https://doi.org/10.1542/peds.2017-0967}.
\bibitem{ref110}U.S. Food and Drug Administration. Final rule: safety and effectiveness of consumer antiseptic rubs; topically applied consumer antiseptic wash drug products, certain active ingredients. September 2016.
\bibitem{ref111}Fenton TR, Huang T. Systematic review of the association between dietary acid load, alkaline water and cancer. BMJ Open. 2016;6(6):e010438. DOI: \url{https://doi.org/10.1136/bmjopen-2015-010438}.
\bibitem{ref112}N\"urnberger F, M\"uller G. So-called cellulite: an invented disease. Journal of Dermatologic Surgery and Oncology. 1978;4(3):221--229.
\bibitem{ref113}American Academy of Dermatology. Wrinkle creams: your guide to younger looking skin. Consumer resource.

\bibitem{ref114}European Academies Science Advisory Council. Statement on homeopathic products and practices. November 2017.
\bibitem{ref115}Saper RB, Kales SN, Paquin J, et al. Heavy metal content of ayurvedic herbal medicine products. JAMA. 2004;292(23):2868--2873. DOI: \url{https://doi.org/10.1001/jama.292.23.2868}.
\bibitem{ref116}Saper RB, Phillips RS, Sehgal A, et al. Lead, mercury, and arsenic in US- and Indian-manufactured ayurvedic medicines sold via the internet. JAMA. 2008;300(8):915--923. DOI: \url{https://doi.org/10.1001/jama.300.8.915}.
\bibitem{ref117}Ang-Lee MK, Moss J, Yuan CS. Herbal medicines and perioperative care. JAMA. 2001;286(2):208--216. DOI: \url{https://doi.org/10.1001/jama.286.2.208}.
\bibitem{ref118}Ross A, Thomas S. The health benefits of yoga and exercise: a review of comparison studies. Journal of Alternative and Complementary Medicine. 2010;16(1):3--12.
\bibitem{ref119}Nortier JL, Martinez M-CM, Schmeiser HH, et al. Urothelial carcinoma associated with the use of a Chinese herb (Aristolochia fangchi). New England Journal of Medicine. 2000;342(23):1686--1692. DOI: \url{https://doi.org/10.1056/NEJM200006083422301}.
\bibitem{ref120}International Agency for Research on Cancer. Aristolochic acids. IARC Monographs on the Evaluation of Carcinogenic Risks to Humans, Volume 82. 2002.
\bibitem{ref121}U.S. Food and Drug Administration. Tainted products marketed as dietary supplements: warning letters and import alerts regarding undeclared drugs in weight-loss products.
\bibitem{ref122}Still J. Use of animal products in traditional Chinese medicine: environmental impact and health hazards. Complementary Therapies in Medicine. 2003;11(2):118--122.
\bibitem{ref123}Cheong YC, Hung Yu Ng E, Ledger WL. Acupuncture and assisted conception. Cochrane Database of Systematic Reviews. 2013, Issue 10, Art. No. CD006920.
\bibitem{ref124}Vickers AJ, Cronin AM, Maschino AC, et al. Acupuncture for chronic pain: individual patient data meta-analysis. Archives of Internal Medicine. 2012;172(19):1444--1453. DOI: \url{https://doi.org/10.1001/archinternmed.2012.3654}.
\bibitem{ref125}Rosa L, Rosa E, Sarner L, Barrett S. A close look at therapeutic touch. JAMA. 1998;279(13):1005--1010. DOI: \url{https://doi.org/10.1001/jama.279.13.1005}.
\bibitem{ref126}National Center for Complementary and Integrative Health. Chiropractic: in depth. Consumer resource.
\bibitem{ref127}Ernst E. Is reflexology an effective intervention? A review of randomised controlled trials. Medical Journal of Australia. 2009;191(5):263--266.
\bibitem{ref128}Moertel CG, Fleming TR, Rubin J, et al. A clinical trial of amygdalin (laetrile) in the treatment of human cancer. New England Journal of Medicine. 1982;306(4):201--206. DOI: \url{https://doi.org/10.1056/NEJM198201283060401}.
\bibitem{ref129}American Cancer Society. Gerson therapy. Consumer resource on unproven cancer treatments.
\bibitem{ref130}U.S. Food and Drug Administration. Warning letters regarding black salve and escharotic products marketed as cancer cures. 2014--2017.
\bibitem{ref131}U.S. Food and Drug Administration. Safety warning regarding chelation products marketed for unapproved uses, including autism spectrum disorder. October 2010.
\bibitem{ref132}National Center for Complementary and Integrative Health. Colloidal silver. Consumer resource.
\bibitem{ref133}Seely DR, Quigley SM, Langman AW. Ear candles: no evidence for their efficacy. The Laryngoscope. 1996;106(10):1226--1229.
\bibitem{ref134}Britton WB, Lindahl JR, Cooper DJ, Canby NK, Palitsky R. Defining and measuring meditation-related adverse effects in mindfulness-based programs. Clinical Psychological Science. 2021;9(6):1185--1204.
\bibitem{ref135}Sherman KJ, Cherkin DC, Wellman RD, et al. A randomized trial comparing yoga, stretching, and a self-care book for chronic low back pain. Archives of Internal Medicine. 2011;171(22):2019--2026. DOI: \url{https://doi.org/10.1001/archinternmed.2011.524}.
\bibitem{ref136}Goyal M, Singh S, Sibinga EMS, et al. Meditation programs for psychological stress and well-being: a systematic review and meta-analysis. JAMA Internal Medicine. 2014;174(3):357--368. DOI: \url{https://doi.org/10.1001/jamainternmed.2013.13018}.
\bibitem{ref137}Cramer H, Ward L, Saper R, Fishbein D, Dobos G, Lauche R. The safety of yoga: a systematic review and meta-analysis of randomized controlled trials. American Journal of Epidemiology. 2015;182(4):281--293.

\bibitem{ref138}Marcone MF. Characterization of the edible bird's nest the ``caviar of the east''. Food Research International. 2005;38(10):1125--1134.
\bibitem{ref139}U.S. Food and Drug Administration. Advice about eating fish: mercury levels in commercial fish and shellfish.
\bibitem{ref140}Feng Y, Siu K, Wang N, et al. Bear bile: dilemma of traditional medicinal use and animal protection. Journal of Ethnobiology and Ethnomedicine. 2009;5:2.
\bibitem{ref141}Shergis JL, Zhang AL, Zhou W, Xue CC. Panax ginseng in randomised controlled trials: a systematic review. Phytotherapy Research. 2013;27(7):949--965.
\bibitem{ref142}Memorial Sloan Kettering Cancer Center. Cordyceps. About Herbs monograph.
\bibitem{ref143}Kim JY, Gum SN, Paik JK, et al. Effects of nattokinase on blood pressure: a randomized, controlled trial. Hypertension Research. 2008;31(8):1583--1588.
\bibitem{ref144}Willcox BJ, Willcox DC, Suzuki M. The Okinawa Program. Clarkson Potter; 2001.
\bibitem{ref145}Hao Q, Dong BR, Wu T. Probiotics for preventing acute respiratory tract infections. Cochrane Database of Systematic Reviews. 2022, Issue 8, Art. No. CD006958.
\bibitem{ref146}Korea Consumer Protection Board. Consumer safety alert on electric fans and air conditioners. 2006.
\bibitem{ref147}Prozialeck WC, Jivan JK, Andurkar SV. Pharmacology of kratom: an emerging botanical agent with stimulant, analgesic and opioid-like properties. Journal of the American Osteopathic Association. 2012;112(12):792--799.
\bibitem{ref148}International Agency for Research on Cancer. Betel-quid and areca-nut chewing. IARC Monographs on the Evaluation of Carcinogenic Risks to Humans, Volume 85. 2004.
\bibitem{ref149}World Health Organization. WHO recommendations on maternal and newborn care for a positive postnatal experience. 2022.
\bibitem{ref150}Talbott SM, Talbott JA, George A, Pugh M. Effect of Tongkat Ali on stress hormones and psychological mood state in moderately stressed subjects. Journal of the International Society of Sports Nutrition. 2013;10(1):28.

\bibitem{ref151}Godlee F, Smith J, Marcovitch H. Wakefield's article linking MMR vaccine and autism was fraudulent. BMJ. 2011;342:c7452.
\bibitem{ref152}Van Noorden R. Science publishing: the trouble with retractions. Nature. 2011;478(7367):26--28.
\bibitem{ref153}Jarvinen TLN, Michaelsson K, Jokihaara J, et al. Overdiagnosis of bone fragility in osteoporosis. BMJ. 2015;350:h2088.
\bibitem{ref154}Moynihan R, Henry D. Too much medicine? The business of health---and its threat to sustainable healthcare. BMJ. 2002;324(7342):886--890.
\bibitem{ref155}Schwartz LM, Woloshin S. Low ``T'' as in ``template'': how to sell disease. JAMA Internal Medicine. 2015;175(1):107--108.
\bibitem{ref156}Cadegiani FA, Kater CE. Adrenal fatigue does not exist: a systematic review. BMC Endocrine Disorders. 2016;16:48.
\bibitem{ref157}American Thyroid Association. Public statement on Wilson's temperature syndrome.
\bibitem{ref158}Mayo Clinic. Candida cleanse diet: what does it treat?
\bibitem{ref159}Camilleri M. Leaky gut: mechanisms, measurement and clinical implications in humans. Gut. 2019;68(8):1516--1526.
\bibitem{ref160}American Association of Endodontists. Root canal safety: myths and facts.
\bibitem{ref161}U.S. Food and Drug Administration. Dental amalgam fillings: consumer information.
\bibitem{ref162}Van Zee A. The promotion and marketing of OxyContin: commercial triumph, public health tragedy. American Journal of Public Health. 2009;99(2):221--227.
\bibitem{ref163}Weissman DE, Haddox JD. Opioid pseudoaddiction: an iatrogenic syndrome. Pain. 1989;36(3):363--366.
\bibitem{ref164}Porter J, Jick H. Addiction rare in patients treated with narcotics. New England Journal of Medicine. 1980;302(2):123.
\bibitem{ref165}Infectious Diseases Society of America, American Academy of Neurology, and American College of Rheumatology. Clinical practice guidelines for the clinical assessment, treatment, and prevention of Lyme disease. Clinical Infectious Diseases. 2021;72(1):e1--e48.
\bibitem{ref166}U.S. Food and Drug Administration. Important patient and consumer information about regenerative medicine therapies.
\bibitem{ref167}Davis TN, O'Reilly M, Kang S, et al. Chelation treatment for autism spectrum disorders: a systematic review. Research in Autism Spectrum Disorders. 2013;7(1):90--98.
\bibitem{ref168}U.S. Food and Drug Administration. Why you should not use ivermectin to treat or prevent COVID-19.
\bibitem{ref169}Brenner DJ, Elliston CD. Estimated radiation risks potentially associated with full-body CT screening. Radiology. 2004;232(3):735--738.
\bibitem{ref170}Corona LE, Swenson CW, Sheetz KH, et al. Use of other treatments before hysterectomy for benign conditions in a hospital system. American Journal of Obstetrics and Gynecology. 2015;212(3):304.e1--304.e9.
\bibitem{ref171}World Health Organization. WHO statement on caesarean section rates. 2015.
\bibitem{ref172}Weinstein JN, Lurie JD, Olson PR, Bronner KK, Fisher ES. United States' trends and regional variations in lumbar fusion surgery: decompression versus fusion. Spine. 2006;31(23):2707--2714.
\end{thebibliography}
\endgroup
\end{document}

